\documentclass{article}

% if you need to pass options to natbib, use, e.g.:
%     \PassOptionsToPackage{numbers, compress}{natbib}
% before loading neurips_2026

% The authors should use one of these tracks.
% Before accepting by the NeurIPS conference, select one of the options below.
% 0. "default" for submission
\usepackage{neurips_2026}

\input{math_commands.tex}
% Revision markup: everything added/changed relative to the submitted version is
% shown in red via \rev{...} (inline) or {\color{red}...} (blocks). All other
% text is unchanged from the original submission.
\newcommand{\rev}[1]{{\color{red}#1}}
% Cost-aware wall-clock revision (Sept 2026): new material is wrapped in \edit{...}
\newcommand{\edit}[1]{{\color{red}#1}}
\newcommand{\newfig}[2][\linewidth]{{\setlength{\fboxrule}{1.2pt}\setlength{\fboxsep}{2pt}\fcolorbox{red}{white}{\includegraphics[width=#1]{#2}}}}
\newcommand{\catimepoisson}{\begin{tabular}{c}(results pending)\end{tabular}}
\newcommand{\caspeedpoisson}{\begin{tabular}{c}(results pending)\end{tabular}}
\newcommand{\causagepoisson}{\begin{tabular}{c}(results pending)\end{tabular}}
\newcommand{\cacostspoisson}{\begin{tabular}{c}(results pending)\end{tabular}}
\newcommand{\caaucpoisson}{\begin{tabular}{c}(results pending)\end{tabular}}
\newcommand{\caseedspoisson}{\begin{tabular}{c}(results pending)\end{tabular}}
\newcommand{\catimeconv}{\begin{tabular}{c}(results pending)\end{tabular}}
\newcommand{\caspeedconv}{\begin{tabular}{c}(results pending)\end{tabular}}
\newcommand{\causageconv}{\begin{tabular}{c}(results pending)\end{tabular}}
\newcommand{\cacostsconv}{\begin{tabular}{c}(results pending)\end{tabular}}
\newcommand{\caaucconv}{\begin{tabular}{c}(results pending)\end{tabular}}
\newcommand{\caseedsconv}{\begin{tabular}{c}(results pending)\end{tabular}}
\newcommand{\catimeaniso}{\begin{tabular}{c}(results pending)\end{tabular}}
\newcommand{\caspeedaniso}{\begin{tabular}{c}(results pending)\end{tabular}}
\newcommand{\causageaniso}{\begin{tabular}{c}(results pending)\end{tabular}}
\newcommand{\cacostsaniso}{\begin{tabular}{c}(results pending)\end{tabular}}
\newcommand{\caaucaniso}{\begin{tabular}{c}(results pending)\end{tabular}}
\newcommand{\caseedsaniso}{\begin{tabular}{c}(results pending)\end{tabular}}
\newcommand{\camainall}{
\begin{tabular}{llccccccccc}
\toprule
& & \multicolumn{3}{c}{$128\times128$} & \multicolumn{3}{c}{$256\times256$} & \multicolumn{3}{c}{$512\times512$} \\ \cmidrule(lr){3-5}\cmidrule(lr){6-8}\cmidrule(lr){9-11}
Equation & Pairing & HINTS-25 & best schedule & router (ours) & HINTS-25 & best schedule & router (ours) & HINTS-25 & best schedule & router (ours) \\ \midrule
 & Multigrid alone (no corrector) &  &  & 989\,$\mu$s & -- & -- & -- & -- & -- & -- \\
 & PCG (MG) (no corrector) &  &  & 1.94\,ms & -- & -- & -- & -- & -- & -- \\
 & Multigrid alone (no corrector) &  &  & 1.01\,ms & -- & -- & -- & -- & -- & -- \\
 & BiCGSTAB (MG) (no corrector) &  &  & 1.39\,ms & -- & -- & -- & -- & -- & -- \\
\bottomrule
\end{tabular}}
\newcommand{\caensnest}{
\begin{tabular}{lllccccccccccc}
\toprule
& & & \multicolumn{3}{c}{$\varepsilon = h^2$} & \multicolumn{3}{c}{$\varepsilon = 10^{-6}$} & \multicolumn{3}{c}{$\varepsilon = 10^{-8}$} & \multicolumn{2}{c}{$\varepsilon = 10^{-8}$: vs.\ best single solver} \\
\cmidrule(lr){4-6}\cmidrule(lr){7-9}\cmidrule(lr){10-12}\cmidrule(lr){13-14}
Equation & $N$ & $\mathcal{W}$ & router & router (WU) & oracle (WU) & router & router (WU) & oracle (WU) & router & router (WU) & oracle (WU) & router & oracle \\ \midrule
Poisson & $128^2$ & $\{Jacobi\}$ & 837\,$\mu$s & 624\,$\mu$s & \textit{614\,$\mu$s} & 967\,$\mu$s & 701\,$\mu$s & \textit{691\,$\mu$s} & 25.6\,ms$^{\dagger 13}$ & 15.6\,ms$^{\dagger 13}$ & \textit{15.4\,ms} & -- & -- \\
 &  & $\{Jacobi (0.67)\}$ & 904\,$\mu$s & 625\,$\mu$s & \textit{615\,$\mu$s} & 1.14\,ms & 717\,$\mu$s & \textit{738\,$\mu$s} & 2.93\,ms & 1.84\,ms & \textit{1.82\,ms} & (best single) & -- \\
 &  & $\{Jacobi, Jacobi (0.67)\}$ & 838\,$\mu$s & 626\,$\mu$s & \textit{614\,$\mu$s} & 1.00\,ms & 734\,$\mu$s & \textit{738\,$\mu$s} & 2.73\,ms & 1.84\,ms & \textit{1.82\,ms} & 1.00$\times$ (slower, $10^{-9}$) & 1.00$\times$ (0.001) \\
\bottomrule
\end{tabular}}
\newcommand{\caNestNum}{1}
\newcommand{\caNestWins}{0}
\newcommand{\caNestLosses}{0}
\newcommand{\caNestSpMin}{1.00$\times$}
\newcommand{\caNestSpMax}{1.00$\times$}
\newcommand{\caNestMinRatio}{1.00$\times$}
\newcommand{\caNestMaxRatio}{1.00$\times$}
\newcommand{\caNestMinRatioH}{1.00$\times$}
\newcommand{\caNestMaxRatioH}{1.00$\times$}
\newcommand{\caensscreen}{\begin{tabular}{c}(results pending)\end{tabular}}
\newcommand{\caens}{\begin{tabular}{c}(results pending)\end{tabular}}
\newcommand{\caensusage}{\begin{tabular}{c}(results pending)\end{tabular}}
\newcommand{\castatsens}{\begin{tabular}{c}(results pending)\end{tabular}}
\newcommand{\caensbig}{\begin{tabular}{c}(results pending)\end{tabular}}
\newcommand{\caoverheads}{\begin{tabular}{c}(results pending)\end{tabular}}
\newcommand{\caamort}{\begin{tabular}{c}(results pending)\end{tabular}}
\newcommand{\caassump}{\begin{tabular}{c}(results pending)\end{tabular}}
\newcommand{\caassumpB}{\begin{tabular}{c}(results pending)\end{tabular}}
\newcommand{\catheorem}{\begin{tabular}{c}(results pending)\end{tabular}}
\newcommand{\caThmMuMax}{--}
\newcommand{\caThmHolds}{--}
\newcommand{\caheldout}{
\begin{tabular}{lllcccc}
\toprule
& & & \multicolumn{2}{c}{development (seed 72, numpy kernels)} & \multicolumn{2}{c}{confirmatory (seed 73, compiled kernels)} \\ \cmidrule(lr){4-5}\cmidrule(lr){6-7}
Equation & $N$ & Pairing / ensemble & vs.\ HINTS-25 or best single & vs.\ best schedule or oracle & vs.\ HINTS-25 or best single & vs.\ best schedule or oracle \\ \midrule
Poisson & $128^2$ & $\{Jacobi, Jacobi (0.67)\}$ ($10^{-8}$, work units) & 1.00$\times$ (0.903) & 0.99$\times$ (slower, $<10^{-10}$) & 1.00$\times$ (slower, $10^{-9}$) & 0.99$\times$ (slower, $10^{-10}$) \\
\bottomrule
\end{tabular}}
\newcommand{\caT}{300}
\newcommand{\caN}{128}
\newcommand{\caLstmMs}{--}
\newcommand{\caLstmOverJacobi}{--}
\newcommand{\caVsMgMin}{--}
\newcommand{\caVsMgMax}{--}
\newcommand{\caVsKrylovMin}{--}
\newcommand{\caVsKrylovMax}{--}
\newcommand{\caVsMgEnsMin}{--}
\newcommand{\caVsMgEnsMax}{--}
\newcommand{\caRhoAMax}{--}
\newcommand{\caRhoSpecMax}{--}
\newcommand{\caRhoTwoGsMax}{--}
\newcommand{\caBandMax}{--}
\newcommand{\caAlphaHatMax}{--}
\newcommand{\caAlphaHatMed}{--}
\newcommand{\caAlphaHatMaxH}{--}
\newcommand{\caAlphaBoundMin}{--}
\newcommand{\caAlphaBoundMax}{--}
\newcommand{\caEnsVsPairMin}{--}
\newcommand{\caEnsVsPairMax}{--}
\newcommand{\caEnsVsSolverMin}{--}
\newcommand{\caEnsVsSolverMax}{--}
\newcommand{\caNumEns}{--}
\newcommand{\caSpSolverMin}{--}
\newcommand{\caSpSolverMax}{--}
\newcommand{\caSpHintsMin}{--}
\newcommand{\caSpHintsMax}{--}
\newcommand{\caSpBestMin}{--}
\newcommand{\caSpBestMax}{--}
\newcommand{\caSpSolverDeepMin}{--}
\newcommand{\caSpSolverDeepMax}{--}
\newcommand{\caSpHintsDeepMin}{--}
\newcommand{\caSpHintsDeepMax}{--}
\newcommand{\caSpBestDeepMin}{--}
\newcommand{\caSpBestDeepMax}{--}
\newcommand{\caSpOracleRatioMin}{--}
\newcommand{\caSpOracleRatioMax}{--}
\newcommand{\caNumCells}{--}
\newcommand{\caCellsRouterBeatsBest}{--}
\newcommand{\caCellsRouterBeatsHints}{--}
\newcommand{\caSpSolverBMin}{--}
\newcommand{\caSpSolverBMax}{--}
\newcommand{\caSpHintsBMin}{--}
\newcommand{\caSpHintsBMax}{--}
\newcommand{\caSpBestBMin}{--}
\newcommand{\caSpBestBMax}{--}
\newcommand{\caSpSolverDeepBMin}{--}
\newcommand{\caSpSolverDeepBMax}{--}
\newcommand{\caSpHintsDeepBMin}{--}
\newcommand{\caSpHintsDeepBMax}{--}
\newcommand{\caSpBestDeepBMin}{--}
\newcommand{\caSpBestDeepBMax}{--}
\newcommand{\caSpOracleRatioBMin}{--}
\newcommand{\caSpOracleRatioBMax}{--}
\newcommand{\caNumCellsB}{--}
\newcommand{\caCellsRouterBeatsBestB}{--}
\newcommand{\caCellsRouterBeatsHintsB}{--}
\newcommand{\caSpSolverCMin}{--}
\newcommand{\caSpSolverCMax}{--}
\newcommand{\caSpHintsCMin}{--}
\newcommand{\caSpHintsCMax}{--}
\newcommand{\caSpBestCMin}{--}
\newcommand{\caSpBestCMax}{--}
\newcommand{\caSpSolverDeepCMin}{--}
\newcommand{\caSpSolverDeepCMax}{--}
\newcommand{\caSpHintsDeepCMin}{--}
\newcommand{\caSpHintsDeepCMax}{--}
\newcommand{\caSpBestDeepCMin}{--}
\newcommand{\caSpBestDeepCMax}{--}
\newcommand{\caSpOracleRatioCMin}{--}
\newcommand{\caSpOracleRatioCMax}{--}
\newcommand{\caNumCellsC}{--}
\newcommand{\caCellsRouterBeatsBestC}{--}
\newcommand{\caCellsRouterBeatsHintsC}{--}

% the "default" option is equal to the "main" option, which is used for the Main Track with double-blind reviewing.
% 1. "main" option is used for the Main Track
%  \usepackage[main]{neurips_2026}
% 2. "position" option is used for the Position Paper Track
%  \usepackage[position]{neurips_2026}
% 3. "eandd" option is used for the Evaluations & Datasets Track
 % \usepackage[eandd]{neurips_2026}
 % if you need to opt-in for a single-blind submission in the E&D track:
 %\usepackage[eandd, nonanonymous]{neurips_2026}
% 4. "creativeai" option is used for the Creative AI Track
%  \usepackage[creativeai]{neurips_2026}
% 5. "sglblindworkshop" option is used for the Workshop with single-blind reviewing
 % \usepackage[sglblindworkshop]{neurips_2026}
% 6. "dblblindworkshop" option is used for the Workshop with double-blind reviewing
%  \usepackage[dblblindworkshop]{neurips_2026}

% After being accepted, the authors should add "final" behind the track to compile a camera-ready version.
% 1. Main Track
 % \usepackage[main, final]{neurips_2026}
% 2. Position Paper Track
%  \usepackage[position, final]{neurips_2026}
% 3. Evaluations & Datasets Track
 % \usepackage[eandd, final]{neurips_2026}
% 4. Creative AI Track
%  \usepackage[creativeai, final]{neurips_2026}
% 5. Workshop with single-blind reviewing
%  \usepackage[sglblindworkshop, final]{neurips_2026}
% 6. Workshop with double-blind reviewing
%  \usepackage[dblblindworkshop, final]{neurips_2026}
% Note. For the workshop paper template, both \title{} and \workshoptitle{} are required, with the former indicating the paper title shown in the title and the latter indicating the workshop title displayed in the footnote.
% For workshops (5., 6.), the authors should add the name of the workshop, "\workshoptitle" command is used to set the workshop title.
% \workshoptitle{WORKSHOP TITLE}

% "preprint" option is used for arXiv or other preprint submissions
 % \usepackage[preprint]{neurips_2026}

% to avoid loading the natbib package, add option nonatbib:
%    \usepackage[nonatbib]{neurips_2026}

\usepackage[utf8]{inputenc} % allow utf-8 input
\usepackage[T1]{fontenc}    % use 8-bit T1 fonts
\usepackage{hyperref}       % hyperlinks
\usepackage{url}            % simple URL typesetting
\usepackage{booktabs}       % professional-quality tables
\usepackage{amsfonts}       % blackboard math symbols
\usepackage{nicefrac}       % compact symbols for 1/2, etc.
\usepackage{microtype}      % microtypography
\usepackage{xcolor}         % colors



\usepackage{algorithm}
\usepackage{algpseudocode}

\newcommand{\ambuj}[1]{\textcolor{red}{AT: #1}}

%%%%%%%%%%%%%%%SAHANA ADDED THIS %%%%%%%%%%%
\usepackage{amsthm}
\usepackage{multirow}
\usepackage{amsmath}
\usepackage{cleveref}
\usepackage{graphicx}
\newtheorem{theorem}{Theorem}[section]
\newtheorem{lemma}[theorem]{Lemma}
\newtheorem{claim}[theorem]{Claim}
\newtheorem{proposition}[theorem]{Proposition}
\newtheorem{corollary}[theorem]{Corollary}
\newtheorem{fact}[theorem]{Fact}
\newtheorem{example}[theorem]{Example}
\newtheorem{notation}[theorem]{Notation}
\newtheorem{observation}[theorem]{Observation}
\newtheorem{conjecture}[theorem]{Conjecture}

\newtheorem{innercustomgeneric}{\customgenericname}
\providecommand{\customgenericname}{}
\newcommand{\newcustomtheorem}[2]{%
  \newenvironment{#1}[1]
  {%
   \renewcommand\customgenericname{#2}%
   \renewcommand\theinnercustomgeneric{##1}%
   \innercustomgeneric
  }
  {\endinnercustomgeneric}
}

\newcustomtheorem{customtheorem}{Theorem}
\newcustomtheorem{customlemma}{Lemma}
\newcustomtheorem{customclaim}{Claim}
\newcustomtheorem{customproposition}{Proposition}
\newcustomtheorem{customcorollary}{Corollary}
\newcustomtheorem{customfact}{Fact}
\newcustomtheorem{customexample}{Example}
\newcustomtheorem{customnotation}{Notation}
\newcustomtheorem{customobservation}{Observation}
\newcustomtheorem{customconjecture}{Conjecture}
%%%%%%%%%%%%%%%%%%%%%%%%%%%%%%%%%%%%%%%%



% Note. For the workshop paper template, both \title{} and \workshoptitle{} are required, with the former indicating the paper title shown in the title and the latter indicating the workshop title displayed in the footnote. 
\title{A Greedy PDE Router for Blending Neural Operators and Classical Methods}


% The \author macro works with any number of authors. There are two commands
% used to separate the names and addresses of multiple authors: \And and \AND.
%
% Using \And between authors leaves it to LaTeX to determine where to break the
% lines. Using \AND forces a line break at that point. So, if LaTeX puts 3 of 4
% authors names on the first line, and the last on the second line, try using
% \AND instead of \And before the third author name.


\author{Sahana Rayan \\
Department of Statistics\\
University of Michigan\\
Ann Arbor, MI 48104, USA \\
\texttt{srayan@umich.edu} \\
\And
Yash Patel \\
Department of Statistics\\
University of Michigan\\
Ann Arbor, MI 48104, USA \\
\texttt{yppatel@umich.edu} \\
\And
Ambuj Tewari \\
Department of Statistics\\
University of Michigan\\
Ann Arbor, MI 48104, USA \\
\texttt{tewaria@umich.edu} \\
}
% \

% \author{%
%   David S.~Hippocampus\thanks{Use footnote for providing further information
%     about author (webpage, alternative address)---\emph{not} for acknowledging
%     funding agencies.} \\
%   Department of Computer Science\\
%   Cranberry-Lemon University\\
%   Pittsburgh, PA 15213 \\
%   \texttt{hippo@cs.cranberry-lemon.edu} \\
%   % examples of more authors
%   % \And
%   % Coauthor \\
%   % Affiliation \\
%   % Address \\
%   % \texttt{email} \\
%   % \AND
%   % Coauthor \\
%   % Affiliation \\
%   % Address \\
%   % \texttt{email} \\
%   % \And
%   % Coauthor \\
%   % Affiliation \\
%   % Address \\
%   % \texttt{email} \\
%   % \And
%   % Coauthor \\
%   % Affiliation \\
%   % Address \\
%   % \texttt{email} \\
% }


\begin{document}


\maketitle

\edit{\noindent\textbf{Revision note.} All text, tables and figures added in this revision are typeset in red (new figures carry a red frame); everything in black is unchanged from the submitted version. The iteration-based experiments of the submitted version on $31^2$ and $63^2$ grids with the LSTM router, and their appendix analyses, have been removed: they are superseded by the wall-clock study of \Cref{sec:experiments,sec:wallclock}, which covers the same solvers and equations (plus anisotropic diffusion) on $128^2$--$512^2$ grids and also reports the iteration-based metrics. Only the Dirichlet experiment of the submitted version is kept (\Cref{sec:moreexperimentalresults}).}


\begin{abstract}
% When solving PDEs, classical numerical solvers are often computationally expensive, while machine learning methods can suffer from spectral bias, failing to capture high-frequency components. Designing an optimal hybrid iterative solver--where, at each iteration, a solver is selected from an ensemble of solvers to leverage their complementary strengths--poses a challenging combinatorial problem. While the greedy selection strategy is desirable for its constant-factor approximation guarantee to the optimal solution, it requires knowledge of the true error at each step, which is generally unavailable in practice. We address this by proposing an approximate greedy router that efficiently mimics a greedy approach to solver selection. Empirical results on the Poisson and convection-diffusion equations demonstrate that our method outperforms single-solver baselines and existing hybrid solver approaches, such as HINTS, achieving faster and more stable convergence. % \ambuj{clarify that approx guarantee is under assumptions. if you can mention the main assumptions in the abstract even better. we don't want to oversell} 
When solving PDEs, classical numerical solvers are often computationally expensive, while machine learning methods can suffer from spectral bias, failing to capture high-frequency components. Designing an optimal hybrid iterative solver--where, at each iteration, a solver is selected from an ensemble of solvers to leverage their complementary strengths--poses a challenging combinatorial problem. While greedy selection is desirable for its constant-factor approximation guarantee to the optimal solution under Lipschitz \edit{and shared-eigenbasis} assumptions, it requires knowledge of the true error at each step, which is unavailable in practice. We address this by proposing an approximate greedy router that efficiently mimics a greedy approach to solver selection. Empirical results on the Poisson, convection–diffusion and anisotropic diffusion equations show that our method reduces the final error and the area-under-the-curve (AUC) of the error trajectory relative to single-solver baselines and existing hybrid approaches such as HINTS in nearly every setting, reaching comparable error levels in substantially fewer iterations with more stable error decay. \edit{With a cost-aware instantiation of the greedy rule, a lightweight router, and a band-limited scale-equivariant corrector, we also measure end-to-end wall-clock time on $128^2$ to $512^2$ grids with compiled classical kernels, four equations (Poisson, convection--diffusion, anisotropic and variable-coefficient diffusion) and six classical solvers including geometric multigrid, against HINTS, the best fixed schedule chosen per setting on the test set, multigrid, preconditioned Krylov methods and direct solves. At truncation-level accuracy the learned router matches, without any search, the best static schedule found by a test-set search (one corrector call, then smoothing) and is \caSpHintsMin--\caSpHintsMax\ faster than HINTS with its published period; below truncation level it tracks the cost-aware oracle and beats HINTS and the one-shot schedule by large factors, while a compiled geometric multigrid remains the strongest classical method there; routing over ensembles of classical solvers is strictly better than every single solver where the error passes through regimes that favour different solvers.}
% Empirical results on the Poisson and convection-diffusion equations demonstrate that our method outperforms single-solver baselines and existing hybrid solver approaches, such as HINTS, achieving faster convergence, often reaching comparable error levels in substantially fewer iterations while exhibiting more stable error decay. 
% \ambuj{make this a bit more concrete. faster to what extent? more stable to what degree?} 
\end{abstract}

\section{Introduction}


Natural phenomena and engineered systems are often governed by ordinary and partial differential equations (PDEs). Solving these equations enables a variety of tasks - predicting the evolution of the system through simulation (e.g., forecasting weather using the Navier-Stokes equations) \citep{kalnay2003atmospheric,bauer2015quiet,staniforth1991semi}, addressing control problems (e.g., optimizing heat shield design for spacecrafts using heat transfer equations) \citep{anderson1989hypersonic,troltzsch2010optimal}, and tackling inverse problems based on external measurements (e.g., reconstructing brain activity from EEG data using electrophysiological models) \citep{baillet2001electromagnetic,grech2008review}. 

Traditionally, PDEs are solved using finite difference methods \citep{smith1985numerical,leveque2007finite}, which discretize the spatial and/or temporal domain and approximate the solution by solving a system of equations at the discrete grid points. Similarly, finite element methods \citep{hughes2003finite,bathe2006finite,brenner2008mathematical} construct a system of equations based on local basis functions, and then rely on numerical linear algebra techniques--such as the Jacobi and Gauss-Seidel method \citep{saad2003iterative}--to compute the solution. These iterative methods, however, \edit{must be re-run from scratch for every new} initial condition, boundary condition, or forcing function, as even minor changes to any of these parameters require re-solving the entire system of equations from scratch. 

These challenges have motivated the development of neural operators \citep{kovachki2023neural}--a class of machine learning models that aim to learn the solution operator directly, enabling fast inference across a range of parameters. Neural operators lift the classical linear layer in neural networks to an operator--typically a kernel integral operator--between function spaces. Despite strong approximation guarantees \citep{chen1995universal}, neural operators are prone to spectral bias \citep{liu2021multiscale,liu2024mitigating,you2024mscalefno,khodakarami2025mitigating,xu2025understanding}, similar to traditional neural networks \citep{rahaman2019spectral,xu2019frequency,luo2019theory}. As a result, they favor learning low-frequency components of the solution function and often struggle to capture high-frequency features that iterative solvers excel at. 

Recognizing this complementarity, \citet{zhang2024blending} proposed HINTS, a hybrid solver that 
% leverages the advantages of classical iterative methods and neural operators by
interleaves a neural operator pass every $\tau^{\text{th}}$ iteration of a classical solver, with $\tau$ fixed a priori.
% applying a classical update at each step $t$ and every $\tau^{\text{th}}$ step (with $\tau$ fixed a priori), it replaces that update with a neural correction. 
While HINTS reports significant empirical gains in convergence and accuracy over the standalone neural operator and Jacobi solver, this fixed schedule can be detrimental: a poorly timed neural correction may increase the error and undo recent progress. 
% \ambuj{clarify ``periodically incorporating predictions" and ``fixed cadence" so that it is easy for the reader to see what the main limitation of HINTS is}

Although state-dependent solver schedules are desirable, this naturally suggest a reinforcement learning formulation for learning routing policies. However, in our setting, the structure of iterative PDE solvers allows for a simpler approach: each solver induces a known update, and its effect can be evaluated through error-based signals.
% Although state-dependent solver schedules are desirable \ambuj{the mention of state will bring RL to the mind of many reviewers. we should mention RL ourselves (we do that in discussion but that is too late) and say that we prefer to exploit special problem structure here to keep our policy simple}, identifying the optimal sequence of solvers that minimizes final error is combinatorial since the search space grows exponentially with the number of iterations. 
Leveraging this, we adopt a greedy rule that selects, at each iteration, the solver with the largest immediate error reduction.
% that most reduces the prediction error at the current step. 
Such a rule can be near-optimal when the final error satisfies (weak) supermodularity--so that applying a beneficial solver earlier is at least as valuable as applying it later. To support this approach, we make the following contributions:
\begin{itemize}
    \item We present a general hybrid PDE solver in which a routing rule chooses, at each iteration, from a set of classical solvers and neural operators. With oracle access to the true error at every step, we show that a greedy routing rule
    achieves a constant-factor approximation to the optimal strategy for linear PDEs, provided the updates are error-reducing (Lipschitz with constant $<1$) and zero-preserving--conditions under which weak supermodularity holds\edit{---and share an orthogonal eigenbasis, which gives the postfix monotonicity the guarantee also needs}.
    \item  Given that the true error is not observed at test time, a convex surrogate loss is introduced, which, when minimized, enables the learned model to imitate the greedy router without access to true error information. This alignment is guaranteed through Bayes consistency.
    \item We empirically demonstrate that our approximate greedy solver achieves faster convergence than HINTS and individual solvers on the Poisson, convection-diffusion, anisotropic and variable-coefficient diffusion equations\edit{, in iterations and in end-to-end wall-clock time on $128^2$ to $512^2$ grids with compiled classical kernels; below truncation-level accuracy it is also faster than every fixed schedule, than geometric multigrid alone and than multigrid-preconditioned Krylov methods, and it extends to ensembles of classical solvers}.
\end{itemize}

\section{Background} \label{sec:background}



Consider the following linear PDE:
\begin{equation} \label{eq:mainpde}
    \begin{split}
    \mathcal{L}_{x}^a\left(u\right) = f ,x \in D; \quad
    \mathcal{B}_{x}\left(u\right) = b ,x \in \partial D
    \end{split}
\end{equation}
where $\mathcal{L}_x^a$ is a differential operator with respect to the spatial variable $x \in D$ parameterized by the coefficient function $a$,  $f$ is a forcing function,  $u$ is the PDE solution, $\mathcal{B}_{x}$ is the boundary operator, and $b$ is the boundary term. Such PDEs can be solved numerically using various discretization strategies, such as finite differences \citep{smith1985numerical,leveque2007finite}, finite element, and spectral methods \citep{boyd2001chebyshev,canuto2006spectral}. Here, we focus on finite differences, which replace derivatives with difference quotients (e.g., $\partial_x u(x) \approx (u(x + h) - u(x))/h$) on a uniform grid.

Formally, let $h$ be the grid size and $G_h(D)$ be a uniform grid with spacing $h$ over $D$. For a function $g: D \to \mathbb{R}$, we denote its restriction to $G_h(D)$ by $g_h \in \mathbb{R}^N$,
% , i.e. vector of $g$ evaluated at the points of $G_h(D)$.
where $N = |G_h(D)|$. If $\mathcal{L}_h^a \in \mathbb{R}^{N \times N}$ is the discretized operator, the PDE reduces to a linear system of equations
\begin{equation} \label{eq:mainpdediscrete}
    \mathcal{L}_h^au_h = f_h,
\end{equation}
with the boundary conditions incorporated in $\mathcal{L}_h^a$. Unlike much of the neural operator literature, which retains function space formulations for discretization invariance, we follow the conventions of classical numerical analysis and represent functions and operators as vectors and matrices. Discretization invariance is not essential here, since we focus on fixed-grid problems.
% and solver collaborations. 

\subsection{Iterative PDE Solvers} \label{sec:finitedifference}
Direct methods like Gaussian elimination and Thomas Algorithm can be computationally expensive in high-dimensional domains when solving systems like \Cref{eq:mainpdediscrete}. In contrast, iterative methods like Jacobi and Gauss-Seidel offer computational speedups by iteratively updating the solution, gradually converging to the true solution for the PDE. The general iterative update is
\begin{equation}\label{eq:iterativeupdate}
    u^{(t + 1)} = u^{(t)} + C\left(f_h - \mathcal{L}_h^au^{(t)}\right),
\end{equation}
where $u^{(t)}$ is the $t^{\text{th}}$ iterate of the solution and $C$ is a preconditioning matrix. Jacobi uses $C = D^{-1}$, where $D$ is the diagonal of $\mathcal{L}_h^a$, and Gauss-Seidel uses $C = (D + L)^{-1}$, where $L$ denotes the strict lower triangular part. 
However, these methods tend to damp low frequency parts of the error slowly. 

% Multigrid solvers \citep{briggs2000multigrid,trottenberg2001multigrid,hackbusch2013multi} 
% % are an advanced class of iterative methods that 
% tackle this problem by alternating smoothing
% % with specialized preconditioning matrices. Specifically, multigrid methods 
% % interleave solving problems 
% on coarse and fine grid resolutions to dampen more uniformly across frequencies. We defer the reader to \Cref{sec:multigrid} for more background on Multigrid solvers.

% and the particular form of the preconditioning matrix.

% \sahana{Might move the next two paragraphs to the appendix}



\subsection{Neural Operators} \label{sec:neuraloperators}

Suppose we observe samples $\{(a_i, f_i, u_i)\}_{i = 1}^N$ such that $(a_i, f_i) \sim \mathcal{P}$ are i.i.d. samples, and $u_i = \mathcal{G}^{*}(a_i, f_i)$ be generated by a deterministic solution operator $\mathcal{G}^{*}$. 
A neural operator (NO) $\mathcal{G}_{\theta}$ seeks to approximate $\mathcal{G}^*$ by minimizing the expected squared $L^2(D)$ error:
\begin{equation*} 
    \mathcal{R}_{\text{NO}}(\mathcal{G}_{\theta})  = \mathbb{E}_{(a, f) \sim  \mathcal{P}}\left[ \int_{D} \left\|\mathcal{G}^{*}(a, f)(x) - \mathcal{G}_{\theta}(a, f)(x)\right\|_2^2dx\right]
\end{equation*}
Neural operators use discretization-invariant layers. Prominent instantiations include Fourier Neural Operators (FNO), which use spectral transforms \citep{li2020fourier}, Graph Neural Operators, which perform graph-based aggregation over sampled points \citep{li2020multipole}, and DeepONets, which employ a trunk–branch decomposition to map input functions to output functions \citep{lu2021learning}. 


\subsection{Greedy Optimization} \label{sec:bggreedy}

Consider maximizing $g(S)$ over $S\subseteq \Omega$ with $|S| \leq T$ where $g: 2^{\Omega} \to \mathbb{R}$ is a set function defined on subsets of a set $\Omega$. An exhaustive search is infeasible, so an efficient alternative is the greedy algorithm. It builds the solution subset iteratively by adding $\omega^*$ that yields the largest marginal gain, i.e. $\omega^* =  \argmax_{\omega \in \Omega \setminus S} g(S \cup \omega)$. When $g$ is non-negative, monotone (adding elements never decreases the value), and submodular (diminishing returns), the greedy algorithm achieves a $(1 - e^{-1})$ approximation to the optimal solution \citep{nemhauser1978analysis}. Formally, a function $g$ is submodular if $g(A \ \cup \{\omega\}) - g(A) \geq g(B \ \cup \{\omega\}) - g(B)$ for all $A \subseteq B \subseteq \Omega$ and $\omega \in \Omega \setminus B$; intuitively, delaying an addition cannot increase its benefit. For set function minimization problems, supermodularity, where the inequality flips, plays an analogous role, and the greedy rule achieves constant-factor approximation guarantees \citep{liberty2017greedy}.

The suboptimality of the greedy algorithm has been extensively studied for both set-function minimization \citep{bounia2023approximating}  and maximization \citep{das2018approximate,feige2011maximizing,bian2017guarantees,harshaw2019submodular} when standard assumptions--non-negativity, monotonicity, and sub/supermodularity--are weakened. In contrast, results on greedy sequence maximization \citep{streeter2008online,alaei2021maximizing,zhang2015string,bernardini2020through,van2024performance,tschiatschek2017selecting}--where the ordering of the elements affects the function--have led to sequential analogues of submodularity and monotonicity. We leverage these tools to analyze the suboptimality of the greedy solution to minimizing the final error.

% \ambuj{main weakness of this sec is that it reads mostly as background. where is the problem setting promised in the title? problem is described actually in section 3. simplest fix might be to remove problem setting from the title of this section}

\section{General framework for Hybrid Solvers} \label{sec:hybridframework}

To solve \Cref{eq:mainpdediscrete}, consider the following hybrid iterative update:
\begin{equation}
     u^{\left(t + 1\right)}_h = u^{\left(t\right)}_h + C_{S_t}\left(f_h - \mathcal{L}_h^au^{\left(t \right)}_h\right)
\end{equation}
where $\mathcal{C} =\{C_{j}\}_{j = 1}^K$ is a set of preconditioning functions and $S_t \in [K] = \{1, \dots, K\}$ indexes the function chosen at step $t$. Here, we use ``preconditioning function'' broadly to include classical linear updates ($C_j(x) = C_jx$ where $C_j$ is a matrix), and learned models, such as neural operators. 
% Thus, each $C \in \mathcal{C}$ corresponds to a particular solver or update mechanism.
This update generalizes the classical iterative update in \Cref{eq:iterativeupdate}, by allowing $C_j$ to be non-linear and enabling the preconditioning function to be adaptively chosen at every step. $\mathcal{C}$ can also accommodate parameterized solver families (e.g., different Jacobi relaxation weights), enabling adaptive selection of solver parameters.  As the number of solvers $K$ grows, the chance of selecting a more effective update increases. Furthermore, HINTS \citep{zhang2024blending} is a special case of this hybrid solver with $K = 2$, where  $C_1$ is a neural operator and $C_2$ is a classical preconditioner. Its routing rule is given by \edit{$S_t = \1_{(t + 1) \bmod \tau > 0} + 1$}, which selects $C_1$ every $\tau$ steps (\edit{first at iteration $\tau$, after $\tau - 1$ classical steps, as in Algorithm S3 of \citet{zhang2024blending}}) and $C_2$ otherwise.

Let $e^{(t)}_h = u_h - u^{(t)}_h$ denote the error at step $t$. Then, 
% the update rule can be written in terms of the error from the previous step:
\begin{align*}
    e^{(t + 1)}_h = u_h -  u^{(t)}_h -C_{S_t}\left(\mathcal{L}_h^au_h - \mathcal{L}_h^au^{\left(t\right)}_h\right) = \left(I- C_{S_t} \circ \mathcal{L}_h^a\right)(e^{(t)}_h)
\end{align*}
where $I$ is the identity map and $I - C_{j} \circ\mathcal{L}_h^a$ is the error propagation function for solver $j$. Here, ``$\circ$'' denotes function composition (matrix multiplication for linear updates).
% composition notation for generality: when $C_{j}$ is a matrix, this coincides with matrix multiplication; for nonlinear $C_j$, it simply means function application, i.e., $(C_j \circ \mathcal{L}_h^a)(x) = C_j\left(\mathcal{L}_h^ax\right)$.
The objective is to select a sequence $S = (S_1, \dots, S_T)$ that minimizes the error norm after $T$ steps or $\|e^{(T)}_h\|_2^2$:
\begin{equation} \label{eq:errorobjective}
    \min_{S_t \in [K], |S| \leq T} h(S) := \left\|\left(I - C_{S_{|S|}} \circ \mathcal{L}_h^a\right) \circ \dots \circ \left(I - C_{S_1} \circ \mathcal{L}_h^a\right)  \left(e^{(0)}_h \right)\right\|^2_2
\end{equation}
with compositions applied from right to left, so that the $S_1$ update acts on the initial error.
% Here, the product is ordered from right to left (i.e. from $t = 1$ to $t = |S|$), reflecting the sequential application of the preconditioners. 

% If running each solver $C_{j}$ incurs a time cost $c_j$, we can formulate a budget constrained optimization problem
% \begin{align*}
%     \min_{S \in [1, K]^*} & h(S) \\
%     \text{subject to} \quad &\sum_{t = 1}^T c_{i_t} \leq B
% \end{align*}

% where $B$ is the total computational budget.

\section{Greedy Algorithm} \label{sec:greedy}

\Cref{eq:errorobjective} defines a combinatorial optimization problem with a search space exponential in $T$, making exact optimization intractable for large $T$ or $K$. As a starting point, we consider an ``omniscient'' greedy algorithm that assumes access to the true initial error $e^{(0)}_h$. This is unrealistic since knowledge of $e^{(0)}_h$ could directly recover the solution via $u_h = u^{(0)}_h + e^{(0)}_h$, but it provides a useful benchmark.
In \Cref{sec:approxgreedy}, we relax this assumption using a practical learning strategy that is Bayes consistent with this omniscient rule, thereby recovering the guarantees shown below.

As discussed in \Cref{sec:bggreedy}, when supermodularity and monotonicity hold, the greedy rule--such as the one described in \Cref{alg:greedy}--enjoys constant-factor approximation guarantees. However, classical results focus on \textit{set} functions, with only recent extensions made to sequences. Building on this line, we introduce a sequence-based notion of weak supermodularity.

\begin{algorithm} 
\caption{Greedy Algorithm for a Hybrid PDE solver}\label{alg:greedy}
\begin{algorithmic}
\Require $\{C_{j}\}_{j = 1}^{K}, T, \mathcal{L}_h^a, e^{(0)}_h$
\State $S^{0} \gets \emptyset$
\For{$t < T$}
    \State $S^{t + 1} \gets S^{t} \ \oplus \argmin_{j \in [K]}\|( I - C_{j} \circ \mathcal{L}_{h}^a)(e^{(t)}_h)\|^2_2$
    \State $e^{(t + 1)}_h \gets (I - C_{S_{t+ 1}} \circ \mathcal{L}_{h}^a) (e^{(t)}_h)$
\EndFor
\State \Return $S^{T}$
\end{algorithmic}
\end{algorithm}

Let $\Omega^*$ denote the space of sequences with elements in $\Omega$. For $S, S' \in \Omega^*$, we denote their concatenation as $S \oplus S'$. $S$ is a prefix of $S'$ or $S \preceq S'$ if $S' = S \ \oplus L$ for some $L \in \Omega^*$.
% If $S' = S \ \oplus L$ for some $L \in \Omega^*$, then $S$ is a prefix of $S'$ (denoted as $S \preceq S'$).
A function $g: \Omega^* \to \mathbb{R}$ is prefix monotonically non-increasing if $g(S \ \oplus S') \leq g(S)$ for all $S, S' \in \Omega^*$, and postfix monotonically non-increasing if $g(S' \ \oplus S) \leq g(S)$ for all $S, S' \in \Omega^*$. A prefix non-increasing function $g$ is sequence supermodular if, for all $S', S \in \Omega^*: S \preceq S'$, it holds that
\begin{equation}\label{eqn:supermod}
    g(S) - g(S \ \oplus \omega ) \geq g(S') - g(S' \ \oplus \omega ), \quad \forall \omega \in \Omega
\end{equation}

However, $h(S)$ (defined in \Cref{eq:errorobjective}) may not satisfy this property in general. Therefore, we introduce weak sequence supermodularity. A prefix non-increasing function $g$ is weakly supermodular with respect to $S' \in \Omega^*$ if, for any \edit{$S \in \Omega^{*}$}, there exists $\alpha(S') \geq 1$ such that
\begin{equation}\label{eqn:alpha_supermod}
     g(S) - g(S\ \oplus S') \leq \alpha(S') \sum_{i \in [|S'|]} g(S) - g(S \ \oplus S'_i)
\end{equation}
The parameter $\alpha(S')$ or the supermodularity ratio quantifies deviation from exact sequence supermodularity.
Expanding the $g(S) - g(S\ \oplus S')$ as a telescoping sum 
$\sum_{i = 1}^{|S'|} g(S \ \oplus \left(S'_1, \dots, S'_{i - 1}\right)) - g(S \ \oplus \left(S'_1, \dots, S'_{i}\right))$
% $$g(S) - g(S\ \oplus S')=\sum_{i = 1}^{|S'|} g(S \ \oplus \left(S'_1, \dots, S'_{i - 1}\right)) - g(S \ \oplus \left(S'_1, \dots, S'_{i}\right))$$ 
% Each term represents 
shows that the marginal decrease from appending $S_i'$ after its predecessors
% . Weak supermodularity ensures that each of these marginal terms 
is controlled by the effect of appending $S_i'$ directly to $S$. Thus, postponing the inclusion of $S_i'$ cannot yield a significantly larger benefit than adding it earlier. The supermodularity ratio $\alpha(S)$ quantifies the extent to which future gains from delays may exceed immediate gains. 
% \sahana{I may not need this. will discuss later}
% $g$ is considered weakly-$\mu$-postfix monotonic if for all $S, S' \in \Omega^*$, 
% \begin{equation*}
%     g(S' \ \oplus S) \leq \mu g(S)
% \end{equation*}
% Intuitively, $\mu$ is a ``slack'' for the factor by which $g$ violates the standard postfix monotonicity; hence, $\mu\ge 1$, with $\mu=1$ recovering standard postfix monotonicity. 
% for any $\omega \in \Omega$. However, the objective function $h(S)$ described in \Cref{eq:errorobjective} may not, in general, satisfy this property. Therefore, we introduce weakly-$\alpha$-supermodularity for sequences. A prefix non-increasing function $g$ is weakly-$\alpha$-supermodular for sequence length $T$ if, for any $S, S' \in \Omega^T$ such that $|S'| = T$, 
% \begin{equation}\label{eqn:alpha_supermod}
%      g(S) - g(S\ \oplus S') \leq \alpha T \max_{i \in [|S'|]} g(S) - g(S \ \oplus S'_i)
% \end{equation}

% See \Cref{sec:proofalpha=1} for a formal proof of this equivalence between \Cref{eqn:alpha_supermod} and \Cref{eqn:supermod} for the $\alpha=1$ case
% If we ignore the summation, by satisifying weak supermodularity, we can expect $g(S \ \oplus \left(S'_1, \dots, S'_{i - 1}\right)) - g(S \ \oplus \left(S'_1, \dots, S'_{i}\right)) \leq \alpha(S')\left(g(S) - g(S \ \oplus S'_i)\right)$ which basically says that there is no benefit of appending $S_i'$ in the future as opposed to a earlier timestep

% To see this, notice that for $\alpha=1$, \Cref{eqn:alpha_supermod} captures the notion that the total effect of concatenating a $T$-length sequence $S'$ cannot improve upon $T$ repeats of the single best step. In other words, there is no way for consecutive steps to ``complement'' one another to accrue a better improvement than the two steps individually combined. See \Cref{sec:proofalpha=1} for a formal proof of this equivalence between \Cref{eqn:alpha_supermod} and \Cref{eqn:supermod} for the $\alpha=1$ case. Thus, for $\alpha > 1$, such complementarity of steps is permitted but only within a bounded factor. Intuitively, therefore, the higher $\alpha$, the more it is possible there will be some long-range complementarity of the selection of actions that render greedy selection suboptimal.

% where $\alpha \geq 1$. When $\alpha = 1$, the functions is sequence supermodular (See \Cref{sec:proofalpha=1} for proof)

\edit{Postfix monotonicity enters the analysis only through the greedy prefixes, so we also allow a slack: $g$ is $\mu$-postfix monotone with respect to $O$ along a sequence $S^T$ if $g(S^t \oplus O) \le \mu\, g(O)$ for every prefix $S^t$, $t < T$; a postfix non-increasing function satisfies this with $\mu = 1$, and $\mu$ can be computed exactly for short horizons (\Cref{sec:wallclock_assumptions}).} Having introduced these notions, we now characterize the suboptimality of greedy solutions of weakly supermodular and 
postfix monotonic sequence functions in \Cref{th:suboptgreedy}.
% and subsequently discuss the interpretation of such results.
% \Cref{th:suboptgreedy} characterizes the suboptimality of greedy solutions to weakly-$\alpha$-supermodular, prefix monotonic, and weakly-$\mu$-postfix monotonic sequence functions. 
\begin{theorem} \label{th:suboptgreedy}
    Let $g: \Omega^* \to \mathbb{R}$ be a non-negative, prefix non-increasing function that is weakly supermodular with respect to the optimal solution $O = \argmin_{S \in \Omega^T} g(S)$ with a supermodularity ratio of $\alpha(O)$, and let the greedy solution of length $T$ be $S^T$. \edit{Suppose $g$ is $\mu$-postfix monotone with respect to $O$ along $S^T$ ($\mu = 1$ if $g$ is postfix non-increasing).} If $\phi_T(\alpha) =  \left(1 - \frac{1}{\alpha T}\right)^T$, then
    \begin{equation*}
        g(S^T) \leq \edit{\mu}\left(1 - \phi_T(\alpha(O))\right) g(O) + \phi_T(\alpha(O))\, g(\emptyset),
        \qquad\text{i.e.,}\qquad
        g(\emptyset) - g(S^T)\geq \left(1 - \phi_T(\alpha(O))\right)\left(g(\emptyset) - \edit{\mu}\, g(O)\right).
    \end{equation*}
\end{theorem}

The proof of \Cref{th:suboptgreedy} appears in \Cref{sec:proofsuboptgreedy}. As $T \to \infty$, the factor $1 - \phi_T(\alpha(O))$ decreases to $1 - e^{-1/\alpha(O)}$, so the worst-case performance of the greedy rule saturates rather than degrades with horizon. \edit{The bound concerns the \emph{reduction} of the error over $T$ steps relative to the best sequence of $T$ steps: with $\alpha = \mu = 1$ it guarantees a $(1 - e^{-1})$ fraction of the optimal reduction, uniformly over ensembles, but it is not a convergence rate---a single convergent solver applied $T$ times already reduces the error by $\rho^{2T}$, which for large $T$ is far below the $e^{-1} g(\emptyset)$ term. In the settings of \Cref{sec:wallclock} the greedy sequence is in fact optimal or within a few percent of the exact optimum over short horizons on most instances (\Cref{sec:wallclock_assumptions}), much better than the worst case.} Additionally, larger $\alpha(O)$, which indicates higher reward for delayed inclusions, loosens the suboptimality bound.  
% $\alpha(O) \to\infty$, the bound loosens, reflecting stronger long-range dependencies where delaying beneficial inclusions is rewarded more.
\Cref{th:suboptgreedy} requires that the sequence objective $h$ be weakly supermodular with respect to the optimal solution and postfix monotone--properties established in Proposition \ref{th:weaklyalphasupermodular}. We defer the proof to \Cref{sec:proofweaklyalphasupermodular}.

\begin{proposition} \label{th:weaklyalphasupermodular}
    Suppose that for all $j \in [K]$, the error propagation function $I - C_j \circ \mathcal{L}_h^a$ is $\rho_j$-Lipschitz continuous with $\rho_j < 1$, and that $(I - C_j \circ \mathcal{L}_h^a) (0_N) = 0_N$. Then, the function $h$ is weakly supermodular with respect to the optimal solution $O$, with 
    \begin{equation*}
        \alpha(O) = \max\left\{\frac{\edit{1}}{T - \sum_{i = 1}^{T} \rho_{O_i}^2} , 1\right\}
    \end{equation*}
    \edit{Furthermore, if the error propagation maps are linear and simultaneously diagonalizable by an orthogonal matrix with eigenvalues of modulus at most one (the setting of \Cref{th:simultaneoussupermodular}), $h$ is also postfix non-increasing; invertibility alone does not suffice (\Cref{sec:proofweaklyalphasupermodular}).}
\end{proposition}
The conditions of Proposition \ref{th:weaklyalphasupermodular} are natural. For classical solvers, the Lipschitz constant is $\|I_N - C_j \mathcal{L}^a_h\|$, which is typically less than $1$
% whenever the update dampens errors--a property that 
for well-posed linear elliptic PDEs (i.e., the update is damping errors). 
For neural networks, Lipschitz continuity can be enforced via weight regularization \citep{gouk2021regularisation} and a sufficiently trained model should approximate $(\mathcal{L}^a_h)^{-1}$ well enough to make the Lipschitz constant small.
% Similarly, a neural network can be trained to be Lipschitz continuous in its inputs by applying appropriate regularization techniques to its weights \citep{gouk2021regularisation}. A small Lipschitz constant implies that the preconditioning function closely approximates the inverse of $\mathcal{L}^a_h$, so a sufficiently well-trained model should satisfy the Lipschitz constant inequality. 
The requirement $(I - C_j \circ \mathcal{L}_h^a) (0_N) = 0_N$
% The condition that each preconditioning function maps $0_N$ to $0_N$ 
is both natural and desirable: it precludes spurious updates when the residual $f_h - \mathcal{L}^a_h u_h^{(t)}$ is $0$. This holds by design for classical schemes, and it can be enforced for any learned model by excluding bias terms.

The form of $\alpha(O)$ in Proposition \ref{th:weaklyalphasupermodular} highlights that the suboptimality factor in \Cref{th:suboptgreedy} is governed by the collective contraction factors of the solvers chosen by the optimal solution: as $\sum_{i = 1}^{T} \rho_{O_i}^2 \to T$, $\alpha(O)$ grows, resulting in a weaker bound. 
% This reflects the fact that when all solvers in the optimal sequence are close to the boundary of stability ($\rho_{O_i} \approx 1$), weak supermodularity degrades and greedy performance suffers. 
\edit{Postfix monotonicity ($\mu = 1$) holds when the error propagation maps share an orthogonal eigenbasis---constant-coefficient periodic operators with Jacobi-type smoothers and band-limited spectral correctors---but it can fail for non-commuting contractions even when all of them are invertible, and a Gauss--Seidel sweep or a learned corrector need not be invertible at all. For the non-commuting solvers of our experiments no a priori bound on $\mu$ is available; \Cref{sec:wallclock_assumptions} computes $\mu$, $\alpha(O)$ and the bound of \Cref{th:suboptgreedy} exactly for short horizons by exhaustive search over sequences, and finds $\mu \le 1.01$ and $\alpha(O) = 1$ on every ensemble of the study, including those with Gauss--Seidel, SOR and multigrid.} 


% in classical solvers, when the residual is $0$, the additive term in \Cref{eq:iterativeupdate} vanishes, ensuring no further correction when convergence to the true solution is achieved. For neural networks, this property holds as long as there are no bias terms included.
% Proposition \ref{th:weaklyalphasupermodular} thus ensures that all the conditions required by \Cref{th:suboptgreedy} are satisfied as long as no solver in $\mathcal{C}$ amplifies errors.
% Notably, any well-trained neural operator is expected to reduce the error norm, and traditional solvers such as Jacobi and Gauss-Seidel are known to exhibit this error-reducing property for linear elliptic PDEs under standard conditions. 
% \subsection{Suboptimality of Greedy Sequence Minimization}

% \sahana{do we need this stuff below?}
\Cref{th:suboptgreedy} thus indicates that the approximation guarantee is strongest when the sequence is nearly supermodular ($\alpha(O) \approx 1$). Generally, if all error propagation maps share an eigenbasis,
% the original problem reduces to an allocation problem--
supermodularity is established.
\edit{This occurs, for example, for linear, constant-coefficient PDEs with periodic boundary conditions when the solver ensemble consists of (damped) Jacobi iterations and band-limited spectral correctors such as a single-layer linear Fourier neural operator, all of which are diagonal in the Fourier basis; lexicographic Gauss--Seidel and SOR are not, because their sweep order breaks translation invariance.} 
% his occurs when the error propagation matrices share an eigenbasis. This happens with linear, constant-coefficient PDEs with periodic boundary conditions. and the collection of solvers are Jacobi, Gauss-Seidel, a single-layer linear Fourier Neural Operator
% This occurs, for example, for linear, constant-coefficient PDEs with periodic boundary conditions, where the Jacobi and Gauss-Seidel error propagators are circulant and hence diagonalizable by the discrete Fourier transform. Likewise, a single-layer linear Fourier Neural Operator is diagonal in the Fourier basis.
% Generally, if all error propagation maps $\left(I - C_j \mathcal{L}_{h}^a\right)$ share an eigenbasis,
% % the original problem reduces to an allocation problem--
% supermodularity is established. 
The proof of \Cref{th:simultaneoussupermodular} is deferred to \Cref{sec:proofsimultaneoussupermodular}.
\begin{proposition} \label{th:simultaneoussupermodular}
    Let $\|I_N - C_j\mathcal{L}_h^a\| \leq 1$ for all $j \in [K]$ and $\left(I - C_j \mathcal{L}_{h}^a\right) = P\Lambda_jP^{-1}$ \edit{with a common orthogonal matrix $P$}. Then, $h$ is supermodular. 
\end{proposition}

\edit{\paragraph{Cost-aware greedy routing.} \Cref{alg:greedy} treats every update as equally expensive, whereas in practice a neural-operator evaluation and a classical sweep have very different costs, and the quantity of interest is the error reached per unit of wall-clock time. We make the greedy rule cost-aware without changing the analysis: let $c_j$ denote the (measured) cost of one application of $C_j$, take the cost $u$ of one neural-operator call as the unit, and replace each preconditioner by the \emph{macro-action} $C_j^{(m_j)}$ that applies $C_j$ for $m_j = \max\{1, \mathrm{round}(u/c_j)\}$ consecutive iterations, so that all actions cheaper than the unit cost approximately $u$; an action dearer than the unit (a multigrid cycle, say) is applied once and its error reduction is compared per unit of cost, $(\|e_j\|/\|e^{(t)}\|)^{u/(m_j c_j)}$, which reduces to the plain comparison when the costs match. For a linear update, $m$ consecutive sweeps are again a stationary iteration, $(I - C_j \circ \mathcal{L}_h^a)^{m} = I - C_j^{(m)}\mathcal{L}_h^a$ with $C_j^{(m)} = \sum_{i=0}^{m-1}(I - C_j\mathcal{L}_h^a)^{i}C_j$, and for a nonlinear $C_j$ the macro-action is the $m$-fold composition of its error propagation map. Hence $\{C^{(m_j)}_j\}_{j=1}^K$ is again a set of preconditioning functions in the sense of \Cref{sec:greedy}: if $I - C_j\circ\mathcal{L}_h^a$ is $\rho_j$-Lipschitz and zero-preserving, then so is its $m_j$-fold composition (with constant $\rho_j^{m_j}$), a shared eigenbasis is inherited by powers, and \Cref{th:suboptgreedy}, Propositions \ref{th:weaklyalphasupermodular} and \ref{th:simultaneoussupermodular}, and the consistency result of \Cref{sec:approxgreedy} apply verbatim to an ensemble of macro-actions that are equalised to the unit cost. Running \Cref{alg:greedy} on such macro-actions selects, at every decision epoch, the action with the largest error reduction per unit cost, and a horizon of $T$ macro-actions corresponds to a wall-clock budget of $\approx T u$. Two caveats: the equalisation is up to rounding (the measured $m_j c_j / u$ lies in $[0.77, 1.18]$ for Jacobi, damped Jacobi, GS and SOR across all grids of \Cref{sec:wallclock}; symmetric GS is dearer than the unit on the larger grids and is compared per unit of cost like multigrid), and for an action dearer than the unit the per-unit-cost comparison is a geometric-rate heuristic rather than an equal-cost step, so the multigrid pairings of \Cref{sec:wallclock} are covered empirically, not by the equal-cost analysis. \Cref{sec:wallclock} uses this cost-aware instantiation throughout.}

\section{Approximate Greedy Router} \label{sec:approxgreedy}

The results in \Cref{sec:greedy} indicate that the error reduction from the greedy solution closely matches that of the optimal sequence, but this is predicated on having access to the initial error, $e^{(0)}_h$. Inaccurate error estimates can result in poor solver decisions, causing errors to amplify. To remedy this, we learn a router $r$ that selects solvers myopically, as in \Cref{alg:greedy}, without access to true errors. 

We adopt the following learning setup. Let $\mathcal{A}, \mathcal{F}, \mathcal{U} \subseteq \mathbb{R}^N$ denote the spaces of coefficient, forcing, and solution functions on the grid $G_h(D)$. We assume an application-specific data distribution $\mathcal{P}_{\mathcal{A} \times \mathcal{F}}$ over $\mathcal{A} \times \mathcal{F}$ that reflects test time conditions.
% the coefficients and forcings the router expects to encounter at deployment.
During training, $(a_h, f_h)$ is drawn from $\mathcal{P}_{\mathcal{A} \times \mathcal{F}}$, and a high-accuracy reference solution $u_h$ is computed, providing the true per-step error. The router $r$ is learned offline on this data; at test time, it operates without access to true errors.

% Let $\mathcal{P}_{\mathcal{A} \times \mathcal{F}}$ denote the data-generating distribution 
% \ambuj{this seems a bit abrupt -- what is *the* data generating distribution? where does it come from? ppl are used to think of solving PDEs as smth that works for every input. why is it okay to assume a distribution here? you need to properly set things up so that it is clear to the reader that there will be training phase when the router will be learned (and at the training stage we have access to the true error) and then it will be deployed. otherwise, this section might be very confusing}. 

At iteration $t$, router $r$ selects a solver using the coefficient $a_h \in \mathcal{A}$, forcing $f_h \in \mathcal{F}$, and the current iterate $u^{(t)}_h \in \mathcal{U}$. If $r(a_h, f_h, u^{(t)}_h) = j$, the next iterate is computed with solver $j$, resulting in an error of $\|( I - C_{j} \circ  \mathcal{L}_{h}^a)(e^{(t)}_h)\|^2_2$. Learning such a router requires minimizing the following loss:
\begin{equation} \label{eq:routerloss}
    l_{\text{route}}\left(r, a_h, f_h, u^{(t)}_h, u_h\right) = \sum_{j = 1}^K \left\|\left( I - C_{j} \circ \mathcal{L}_{h}^a\right)\left(u_h - u^{(t)}_h\right)\right\|^2_2\1_{r(a_h, f_h, u^{(t)}_h) = j}
\end{equation}
with risk $\mathcal{R}_{\text{route}}\left(r\right) = \mathbb{E}_{a_h, f_h \sim \mathcal{P}_{\mathcal{A} \times \mathcal{F}}}[l_{\text{route}}(r, a_h, f_h, u^{(t)}_h, u_h)]$.  
For analysis, we fix the iterate generation via teacher-forcing \citep{williams1989learning,lamb2016professor}: during training, 
% Teacher forcing is a training strategy for sequence models in which, at each step, the model is fed the ground-truth previous output rather than its own previous prediction. 
the iterates fed to the router are produced by an oracle greedy rollout, not by the router's own past choices. Formally, with $u_h^{(0)} = 0_N$, for $t \geq 1$, 
\begin{equation*}
    j^*_t \in \text{argmin}_{j \in [K]} \left\|\left( I - C_{j} \circ \mathcal{L}_{h}^a\right)\left(e^{(t - 1)}_h\right)\right\|^2_2, \quad u^{(t)}_h = u^{(t - 1)}_h + C_{j^{*}_t}\left(f_h - \mathcal{L}_h^a u_h^{(t - 1)}\right)
\end{equation*}
Thus, $\{u^{(t)}_h\}$ is a deterministic function of $(a_h, f_h)$. At each $t$, $l_{\text{route}}$ is evaluated on $r(a_h, f_h, u^{(t)}_h)$ which takes in the teacher-forced iterate. Finally, the greedy choice $j^*_{t + 1}$ is used to advance $u^{(t + 1)}_h$. Thus the input distribution seen by $r$ depends on $(a_h, f_h)$, not on the router's predictions. However, full teacher forcing induces a distributional mismatch at test time (exposure bias). In experiments, we mitigate this with scheduled sampling \citep{bengio2015scheduled}; see \Cref{sec:training details} for details.

% The router then makes a decision using the greedy decision from the previous step and $a_h, f_h$, the loss is evaluated and the next time step is again teacher forcer
% We emphasize that a distribution over all solution functions is not needed here, since each solution is fully determined by $(a_h, f_h)$. By construction, the entire sequence $\{u^{(t)}_h\}$ is fully determined by $a_h, f_h$, and the initial guess $u^{(0)}_h$. In this work, we assume that $u^{(0)}_h$ is generated by a neural operator, making it a deterministic function of $a_h$ and $f_h$. 

Minimizing $\mathcal{R}_{\text{route}}\left(r\right)$ yields the following Bayes-Optimal Router:
\begin{equation}\label{eq:bayesoptimalrouter}
    r^*(a_h, f_h, u^{(t)}_h) \in \text{argmin}_{j \in [K]} \left\|\left( I - C_{j} \circ \mathcal{L}_{h}^a\right)\left(e^{(t)}_h\right)\right\|^2_2
\end{equation} 
which aligns with the update in \Cref{alg:greedy}. Hence, $l_{\text{route}}$ is consistent with learning the greedy rule.

\subsection{Surrogate Loss}\label{section:surrogate}

Since \Cref{eq:routerloss} is discontinuous and non-convex, direct minimization is intractable in practice. We instead introduce a surrogate loss that (1) is convex, enabling efficient optimization; (2) upper bounds the original loss; and (3) is Bayes consistent, ensuring the Bayes optimal decision of the original loss is preserved upon minimization. Formally, a surrogate $\phi$ is considered to be Bayes consistent with respect to the loss $l$ if 
\begin{equation*}
    \lim_{n \to \infty} \mathcal{R}_{\phi}(f_n) - \mathcal{R}_{\phi}^* \implies \lim_{n \to \infty} \mathcal{R}_{l}(f_n) - \mathcal{R}_{l}^* 
\end{equation*}
where $\mathcal{R}_{l}(f) = \mathbb{E}\left[l(f(X), Y)\right]$ and $\mathcal{R}_{l}^* = \inf_f\mathcal{R}_{l}(f)$. In other words, in the limit of infinite data, if the risk of a sequence of learned hypotheses $\{f_n\}$ converges to the optimal risk under $\phi$, it also converges to the optimal risk with respect to the original loss $l$.

To define a surrogate loss for the routing problem, consider a set of scoring functions $\mathbf{g} =\{g_j\}_{j = 1}^K$ with $g_j: \mathcal{A} \times \mathcal{F} \times \mathcal{U} \to \mathbb{R}$ and we define the router as $r(a, f, u^{(t)}) = \text{argmax}_{j \in [K]} g_j(a, f, u^{(t)})$. For example, $\mathbf{g}$ can be a neural network with $K$ outputs. Then, minimizing the following surrogate loss yields the same decision as minimizing \Cref{eq:routerloss}:
\begin{equation} \label{eq:routersurrogate}
    \Psi\left(\mathbf{g}, a_h, f_h, u^{(t)}_h, u_h\right) = -\sum_{j = 1}^K \sum_{k = 1}^K\Tilde{c}_k(a_h, u^{(t)}_h, u_h)\1_{k\neq j}\log\left(\frac{\exp\left(g_j(a_h, f_h, u^{(t)}_h)\right)}{\sum_{m = 1}^K\exp\left(g_m(a_h, f_h, u^{(t)}_h)\right)}\right)
\end{equation}
where $\Tilde{c}_j(a_h, u^{(t)}_h, u_h) = \|( I - C_{j} \circ \mathcal{L}_{h}^a)(u_h - u^{(t)}_h)\|^2_2$.
% where 
% $\psi\left(\mathbf{g}(a, f, u^{(t)}),j\right)$ is a consistent surrogate for the multiclass $0/1$ loss and 
% $\Tilde{e} = \max_{j \in [1, K]} \left\|\left( I_N - C_{j} \mathcal{L}_{h}\right)e^{(t)}\right\|^2_2$. 
The $\Psi$-risk is denoted by $\mathcal{R}_{\Psi}(\mathbf{g}) = \mathbb{E}_{a_h, f_h \sim \mathcal{P}_{\mathcal{A} \times \mathcal{F}}}[\Psi(\mathbf{g}, a_h, f_h, u^{(t)}_h, u_h)]$. The convexity of $\Psi$ with respect to $\mathbf{g}$ follows from the convexity of log-softmax function in its inputs. Moreover, $\Psi$ upper bounds $l_{\text{route}}$ up to a constant factor, and we refer the reader to \Cref{sec:proofupperbound} for the proof. Finally, \Cref{th:consistency} shows that $\Psi$ achieves Bayes consistency with respect to $l_{\text{route}}$; the proof can be found in \Cref{sec:proofconsistency}.

\begin{theorem} \label{th:consistency}
    Let $\Tilde{c}_j(a_h, u^{(t)}_h, u_h) < \Bar{E} < \infty$  for all $j \in [K]$. If there exists $j\in [K]$ such that $\Tilde{c}_j(a_h, u^{(t)}_h, u_h) > E_{min} > 0$, then, for any collection of solvers $\{C_j\}_{j= 1}^K$ and linear discrete operator $\mathcal{L}_h^a$, $\Psi$ is Bayes consistent surrogate for $l_{\text{route}}$.
\end{theorem}
\Cref{th:consistency} is a cost-sensitive analogue of the classical Bayes-consistency of multiclass cross-entropy for 0-1 loss: in the infinite-sample limit, minimizing cross-entropy recovers the true conditional class probabilities, so the induced decision is Bayes optimal. Our result extends this to cross-entropy with instance-dependent weights $\sum_{k \neq j}^K\Tilde{c}_k(a_h, u^{(t)}_h, u_h)$. The uniform upper bound holds when all preconditioning functions are error-damping. It is also reasonable to assume at least one solver cannot annihilate the error in one step, yielding the lower bound. Under these conditions, minimizing $\Psi$ recovers the Bayes-optimal router of \Cref{eq:bayesoptimalrouter}, i.e., the greedy solution of \Cref{eq:errorobjective}. Following the work of \citet{mao2024regression}, the proof uses standard conditional-risk calibration: we relate the excess risks of $l_{\text{route}}$ and $\Psi$ and take the infinite-sample limit. \edit{This is a per-decision statement on the state distribution the router is trained on: it says that minimising $\Psi$ recovers the greedy decision at each visited state, not that the induced sequence, its time to tolerance, or its distance to the optimal sequence are controlled; the latter would require a sequential argument about the router's own state distribution, which we do not provide and mitigate empirically with DAgger (\Cref{sec:wallclock}).}
% \ambuj{ we should demystify the consistency result a bit -- if we had no weights and our non-convex loss was just 0-1 loss then this result simply follows from the fact that cross-entropy loss is proper: in the limit of infinite data, you'll learn the true probabilities of all labels and therefore argmaxing over them will yield the Bayes classifier. The result above is more general because we have costs weightings $\tilde{c}_j$ but the spirit is the . also, how much new work do we have to do beyond Mao-Mohri-Zhong? prolly a sentence on what it took to prove this will also help}


% \sahana{maybe a generalization bound here}

\section{Related Works}

\paragraph{Hybrid PDE Solvers:}
Early data-driven solvers sought convergence guarantees by predicting parameters--e.g., preconditioning matrices, multi-grid smoothers or restriction matrices--within iterative schemes \citep{taghibakhshi2021optimization,caldana2024deep,kopanivcakova2025deeponet,huang2022learning,katrutsa2020black}. These works, however, do not leverage the neural surrogates' ability to generalize across varying coefficients and forcings highlighted in \Cref{sec:neuraloperators} and thus, offer only modest speedups over classical solvers. HINTS \citep{zhang2024blending} introduced hybrid solvers that interleave a classical method with a pre-trained DeepONet on a fixed schedule (e.g., 24 Jacobi steps, then one DeepONet correction), achieving faster convergence than the standalone numerical solver. 
Subsequent works extend this idea to new geometries \citep{kahana2023geometry} and analyze error mode damping, replacing DeepONet with alternatives such as MIONet \citep{jin2022mionet} \citep{hu2025hybrid} or FNO \citep{li2020fourier,cui2022fourier}. 
% \citet{kahana2023geometry} adapted HINTS to geometries distinct from, but related to, those used to train the neural operator. Both \citet{hu2025hybrid} and \citet{cui2022fourier} characterized the dampening of error modes over iterates, with the former replacing the DeepONet of HINTS with an MIONet \citep{jin2022mionet} and the latter an FNO \citep{li2020fourier}.
However, these hybrid solvers are limited in two ways: they, firstly, only combine the trained surrogate with a \textit{single} numerical solver, and secondly rely on a fixed, heuristic schedule. Our proposed method addresses these shortcomings, resulting in significant empirical improvements (\Cref{sec:experiments}).


\paragraph{Model Routing:} Routing \citep{shnitzer2023large,hu2024routerbench,ding2024hybrid,huang2025routereval} selects, for each input, a model from a fixed set to optimize a task metric under cost or latency constraints. Simple heuristics--e.g., thresholding a cheap model's uncertainty estimates \citep{chuang2024learning,chuang2025confident}--often poorly balance the cost-accuracy tradeoff, motivating many systems to instead learn a router that maps inputs to the best model \citep{hari2023tryage,mohammadshahi2024routoo,vsakota2024fly}. We similarly learn a router, but over numerical and neural solvers for PDEs at the iteration level. Our work is the first to apply routing to the problem of learning hybrid PDE solvers, in contrast to prior approaches that use fixed schedules. Additionally, by exploiting the algebraic structure of PDE solvers, we derive theoretical guarantees for our routing strategy (\Cref{sec:greedy}), unlike typical model-routing settings.

\paragraph{Mixture of Experts:} Classical mixture-of-experts (MoE) \citep{jacobs1991adaptive,jordan1994hierarchical} uses a gating network to combine the outputs of multiple experts, trained jointly with the gate. In modern LLMs, MoE instead performs sparse, token-level routing to a subset of in-layer experts, enabling large capacity without proportional compute, typically with load-balancing mechanisms \citep{shazeer2017outrageously,lepikhin2020gshard,fedus2022switch,jiang2024mixtral}. 
Our method can be regarded as gating of a different kind: ``experts'' (solvers) are not jointly trained with the router, and a single router is reused across iterations, unlike MoE which commonly employs layer-specific gates.

\section{Experiments}\label{sec:experiments}

\edit{We evaluate the approximate greedy router on four PDEs posed on the unit domain $D = [0, 1]^2$ with periodic boundary conditions:}
\begin{equation*}
    \text{\textbf{Poisson:} } - \Delta u = f,
    \quad
    \text{\textbf{ConvDiff:} } -\Delta u + 20\,\partial_{x_1}u + 20\,\partial_{x_2}u = f,
    \quad
    \edit{\text{\textbf{AnisoDiff:} } -\epsilon\,\partial_{x_1}^2 u - \partial_{x_2}^2 u = f,\ \ \epsilon = 10^{-2},
    \quad
    \text{\textbf{VarCoeff:} } -\nabla\cdot(a(\mathbf{x})\nabla u) = f,}
\end{equation*}
\edit{where $f$ is centred to satisfy the compatibility condition $\int_D f(\mathbf{x})\,d\mathbf{x} = 0$ and is drawn from a zero-mean hierarchical Gaussian random field with covariance operator $\alpha(-\Delta + \beta I)^{-\gamma}$ and sample-dependent $(\alpha, \beta, \gamma)$ (\Cref{sec:training details}). Anisotropic diffusion is the standard problem on which point relaxation fails to smooth: an error mode oscillatory in $x_1$ and smooth in $x_2$ is damped by only $(1-\epsilon)/(1+\epsilon) \approx 0.98$ per Jacobi sweep. Variable-coefficient diffusion uses a fixed smooth coefficient field $a(\mathbf{x}) = \exp(\kappa z(\mathbf{x}))$ with $\max a/\min a = 10$: its discrete operator is not diagonalised by the FFT, so no direct spectral solve exists, the corrector's band-limited inverse has no closed form and must be learned from data, and the corrector is only approximately accurate (\Cref{sec:wc_varcoeff}). The domain is discretised on $128^2$, $256^2$ and $512^2$ grids; results with Dirichlet boundary conditions, obtained with an earlier version of the pipeline, are deferred to \Cref{sec:dirichlet}. Because iteration counts favour methods whose iterations are expensive, our primary metric is the end-to-end \emph{wall-clock time to a target relative error} $\varepsilon = \|e^{(t)}_h\|/\|u_h\|$ on $64$, $32$ and $16$ test instances per grid, with $\varepsilon = h^2$---a grid-dependent tolerance that lies below the measured discretisation error of every test instance, so that further algebraic refinement does not improve the computed solution (\Cref{sec:wallclock_protocol})---highlighted; every comparison is paired over instances and tested with a one-sided Wilcoxon signed-rank test. We also report the iteration-based metrics of the original submission, the final relative error and the area under the error curve, $\mathrm{AUC} = \sum_{t=1}^{T}\|e^{(t)}_h\|/\|u_h\|$ over $T = \caT$ iterations. \Cref{sec:wallclock} gives the full protocol and collects every result per equation.}

\edit{\paragraph{Setup.} The classical solvers are Jacobi, damped Jacobi ($\omega = 0.67$), Gauss--Seidel (GS), symmetric GS (SymGS), SOR ($\omega = 1.5$) and geometric multigrid (V(2,2)-cycles with GS smoothing), all implemented as compiled single-threaded C stencil kernels (a GS sweep costs $41\,\mu$s and a V-cycle $0.34$\,ms at $128^2$), so the classical baselines are not slowed by the implementation. The neural operator is a DeepONet corrector that acts on the residual through a $64\times 64$ sensor grid with a Fourier trunk basis and is applied scale-equivariantly: its deployed form is a linear, Fourier-structured map (a least-squares fit of the band-limited inverse of the discrete operator, which it reproduces to $\approx 10^{-6}$), so one call solves the smooth band of the spectrum at the cost of about $12$ Jacobi or $3$ GS sweeps at $128^2$ (and only $2.5$ Jacobi sweeps at $512^2$, since its cost is nearly grid-independent) and never injects error outside its band. The router is made \emph{cost-aware} through the cost-equalised macro-actions of \Cref{sec:greedy} and selects from $6{+}K$ scalar features (log residual, its recent decay rate, the previous action and call counts) with a two-layer MLP whose decision costs $\approx 5\,\mu$s independently of $N$; it is trained with the surrogate of \Cref{section:surrogate} to imitate the cost-aware greedy oracle on oracle rollouts followed by DAgger rounds. We compare against the classical solver alone, HINTS with the published period $\tau{=}25$, the best fixed schedule chosen \emph{a posteriori} on the test set for each cell from HINTS with $\tau\in\{2,5,10,25,50\}$, phase-shifted HINTS (first call at iteration $0$) and the static one-shot schedule (corrector once at iteration $0$, then the solver), the cost-aware greedy oracle (\Cref{alg:greedy} with the true error; not deployable), and the classical methods a practitioner would use on these problems without a corrector, multigrid alone and multigrid-preconditioned CG or BiCGSTAB. All configurations were fixed on a development test set and every number reported is a single confirmatory evaluation on a fresh test set with three timed replays (\Cref{sec:wallclock_heldout}).}

\edit{\paragraph{Paired solvers.} \Cref{tab:ca_main} reports, for every pairing $\{\mathrm{NO}, \text{solver}\}$, equation and grid, the median wall-clock time to $\varepsilon = h^2$ of HINTS, of the best fixed schedule, and of the learned router, together with the classical methods without corrector. Two regimes appear. \emph{At truncation level} the whole solve is one corrector call followed by a few sweeps: the router, the cost-aware oracle and the best static schedule (the corrector at iteration $0$, then smoothing) take the same decisions, so the router is within a few percent of the best fixed schedule found by a test-set search---paired ratios \caSpBestMin--\caSpBestMax\ at $128^2$, \caSpBestBMin--\caSpBestBMax\ at $256^2$ and \caSpBestCMin--\caSpBestCMax\ at $512^2$, the difference being its own decision overhead---without any search, and it is \caSpHintsMin--\caSpHintsMax\ faster than HINTS with its published period $\tau{=}25$ (\caSpHintsBMin--\caSpHintsBMax\ at $256^2$, \caSpHintsCMin--\caSpHintsCMax\ at $512^2$), whose first call comes $\tau - 1$ sweeps too late; the classical solver alone is \caSpSolverMin--\caSpSolverMax\ slower on the isotropic equations and never reaches $h^2$ on the anisotropic one. \emph{Below truncation level} the schedule matters and the picture is mixed: down to $\varepsilon = 10^{-8}$ the router is \caSpHintsDeepMin--\caSpHintsDeepMax\ faster than HINTS ($\tau{=}25$) and one to two orders of magnitude faster than the one-shot schedule, because it calls the corrector again when the smooth error has re-emerged (\Cref{fig:ca_usage_main}), but it beats the best fixed schedule in only \caCellsRouterBeatsBestDeep\ of the \caNumCells\ pairings at $128^2$ (paired ratios \caSpBestDeepMin--\caSpBestDeepMax): the best static schedule there is a phase-shifted HINTS with a short period, which calls the corrector every $5$--$10$ sweeps, more often than a cost-equalised macro-action of $12$ Gauss--Seidel or $39$ Jacobi sweeps allows the router to (\Cref{sec:wallclock_granularity} examines this granularity limit). The router tracks the cost-aware oracle within \caSpOracleRatioMax\ in median time and its decisions are identical across five independent training seeds. Against the classical methods without corrector, the compiled multigrid is a strong competitor: at $128^2$ the stationary hybrids match it at $h^2$ (router \caVsMgAllMin--\caVsMgAllMax\ of a stand-alone V-cycle solve) and are slower than it below $h^2$ (\caVsMgDeepMin--\caVsMgDeepMax), whereas routing over $\{\mathrm{NO}, \text{multigrid}\}$ is \caVsMgEnsAllMin--\caVsMgEnsAllMax\ and \caVsMgEnsDeepMin--\caVsMgEnsDeepMax\ of multigrid alone at $h^2$ and $10^{-8}$; the router is \caVsKrylovAllMin--\caVsKrylovAllMax\ faster than multigrid-preconditioned CG or BiCGSTAB, and on the constant-coefficient equations the FFT direct solve remains \caFftRatioMin--\caFftRatioMax\ faster than any iterative method including ours, whereas on variable-coefficient diffusion, which has no such solve, the router is compared with the sparse direct solve (\Cref{sec:wc_varcoeff}). The iteration-based metrics tell the same story: at $128^2$ the router's AUC over $T = \caT$ iterations is orders of magnitude below HINTS' and the solver-only baseline's for every pairing and equation (paired $t$-tests, $p < 10^{-3}$; \Cref{tab:ca_auc_poisson,tab:ca_auc_conv,tab:ca_auc_aniso,tab:ca_auc_var}).}

\edit{\paragraph{Solver ensembles.} The same construction routes over ensembles $\mathrm{NO}\cup\mathcal{W}$ of several classical solvers. Since the greedy rule is a per-decision comparison, a larger ensemble can only help where the error passes through regimes that favour \emph{different} members; an exhaustive oracle-level screening of every subset of up to three members (\Cref{sec:wallclock_ens}) shows that this never happens at $\varepsilon = h^2$---one corrector call and one macro-action of the cheapest smoother already reach truncation level---but does at tolerances far below $h^2$, where the smoother must damp the whole band outside the corrector's reach: undamped Jacobi is cheapest on the modes near the band edge but cannot damp the near-Nyquist modes, which damped Jacobi removes. We therefore evaluate the learned router over the nested sets $\{\text{Jacobi}\}$, $\{\text{Jacobi (0.67)}\}$, $\{\text{Jacobi}, \text{Jacobi (0.67)}\}$ and $\{\text{Jacobi}, \text{Jacobi (0.67)}, \text{GS}\}$ on every equation at $128^2$ and $256^2$, with the single-member routers evaluated in the same session on the same instances. Two properties emerge. At $\varepsilon = 10^{-8}$ the ensemble router is strictly faster than every single-member router in \caNestWins\ of the \caNestNum\ ensemble cells (six settings, two nested sets) (\caNestSpMin--\caNestSpMax\ over the best single solver, $p<0.01$, largest on the $256^2$ grids), reproducing its oracle: it learns a \emph{nonstationary} damping schedule---undamped Jacobi while the band-edge modes dominate, damped Jacobi once the near-Nyquist modes do---that no single member can reproduce. And enlarging the set is nearly free: adding GS changes nothing, because the oracle ignores a member that never wins a cost-aware comparison and the router imitates it, and at $h^2$ every ensemble router is within a few percent of the best single-member router; the one exception (Poisson at $128^2$, $10^{-8}$, where the ensemble router is $8\%$ slower) is inherited from the greedy oracle itself, which is not monotone in its action set (\Cref{sec:wallclock_ens}). Finally, we verify numerically that the assumptions of \Cref{sec:greedy,sec:approxgreedy} hold in all these settings, in the energy norm for the non-normal sweeps (\Cref{sec:wallclock_assumptions}).}

\begin{table}[t]
\color{red}
\centering
\caption{\edit{Median wall-clock time to truncation-level relative error $\varepsilon = h^2$ for every solver pairing $\{\mathrm{NO}, \text{solver}\}$, equation and grid ($64$, $32$ and $16$ test instances at $128^2$, $256^2$ and $512^2$): HINTS with the published period $\tau{=}25$, HINTS with the best period chosen per cell on the test set itself (subscript), and the learned router (paired median speedup over HINTS-25\,/\,over the best period in parentheses; bold: faster than both). The classical solver alone is $10^2$--$10^3\times$ slower than the router on the isotropic equations and never reaches $h^2$ on the anisotropic one (\Cref{sec:wallclock}). The last rows of each block give geometric multigrid alone and the best multigrid-preconditioned Krylov method, without corrector, on the same instances. Anisotropic diffusion is not run at $512^2$: its corrector keeps full resolution along $x_1$, and the resulting $65{,}536\times 65{,}536$ branch matrix does not fit in memory (\Cref{sec:wallclock}); multigrid is not run on it either (\Cref{sec:wallclock_protocol}). $\dagger n$: $n$ instances did not reach $h^2$ within the iteration cap.}}
\resizebox{\textwidth}{!}{\camainall}
\label{tab:ca_main}
\end{table}

\begin{figure}[t]
    \color{red}
    \centering
    \newfig[0.98\linewidth]{neurips_images/ca_router_usage.png}
    \caption{\edit{Corrector usage on the $\caN\times\caN$ grid: fraction of test runs that execute a corrector call at each iteration, for the learned router, the cost-aware oracle and HINTS ($\tau{=}25$), for every solver pairing (columns) and all three equations (rows); the same figure for the $256^2$ and $512^2$ grids is in \Cref{sec:wallclock}.}}
    \label{fig:ca_usage_main}
\end{figure}

\edit{Every result per equation and grid, the ensemble study, the accounting of overheads and training cost, and the verification of the theoretical assumptions are collected in \Cref{sec:wallclock}; \Cref{sec:moreexperimentalresults} reports the Dirichlet experiment of the original submission.}

\section{Discussion} \label{sec:discussion}

We have introduced an adaptive method for selecting solvers that efficiently minimizes the final error of a PDE in iterative methods, opening many directions for future work. Our current DeepONet is trained without accounting for its downstream role in correction term prediction; jointly training the ML solver and router could yield larger gains. Another avenue is to use reinforcement learning to learn a cost-aware routing strategy that optimizes error under compute budgets\edit{; \Cref{sec:greedy,sec:wallclock} show that a myopic cost-aware variant of the greedy rule---\Cref{alg:greedy} on cost-equalised macro-actions---already delivers uniform wall-clock gains, and RL could extend it to non-myopic budgeted routing}. RL would additionally enable extending the ensemble routing to \textit{continuous} parameterization of solvers. Finally, framing routing as an online learning problem could enable adaptation to unknown deployment conditions. 
% where offline training may degrade and fixed schedules like HINTS can perform better. 

\bibliography{mybib}
\bibliographystyle{unsrtnat}




%%%%%%%%%%%%%%%%%%%%%%%%%%%%%%%%%%%%%%%%%%%%%%%%%%%%%%%%%%%%

\appendix
% \section{Background} \label{sec:appbackground}

% \subsection{Multigrid} \label{sec:multigrid}

% Let $A_h, A_{2h}$ represent the coefficient matrix on a fine grid with discretization parameter $h$ and $2h$, $u_h$ and $f_h$ represent the PDE solution and constant vector on a grid discretized by $h$, $R^{2h}_h$ denote the restriction matrix which transfers vectors from a fine grid to a coarse one, and $I^{h}_{2h}$ is the interpolation matrix transfers vectors from a coarse grid to a fine one. In a 2-grid method, a few iterations of the smoother (e.x. Jacobi or Gauss-Seidel) are first applied on the fine grid to approximate the solution of $A_hu_h = f_h$. The residual is then computed as $r_h = f_h - A_hu_h$ and restricted to the coarse grid via $r_{2h} = R^{2h}_hr_h$. The error equation, $A_{2h}e_{2h} = r_{2h}$, is solved on the coarse grid. The resulting estimate of the error is interpolated in the fine grid by $e_h = I^{h}_{2h}e_{2h}$, and the fine grid solution is updated by adding this correction, $u_h = u_h + e_h$. Finally, additional smoothing steps are performed on the fine grid to further reduce any high frequency errors. The preconditioning matrix for two-grid solver is $C_\mathrm{2G} = I^{h}_{2h}A_{2h}^{-1}R^{2h}_h$. More complex strategies for multigrid like V-cycle and W-cycle compute error corrections recursively across multiple grids of varying coarseness %(See \Cref{sec:multigrid} for pseudocode of these methods).

\section{Proofs for \Cref{sec:greedy}} \label{sec:proofgreedy}

\subsection{Proof of Proposition \ref{th:alpha=1}} \label{sec:proofalpha=1}

\begin{proposition} \label{th:alpha=1}
    Any prefix monotonically non-increasing sequence supermodular function $g$ is weakly supermodular with respect to all sequences $S \in \Omega^*$ with $\alpha(S) = 1$
\end{proposition}

\begin{proof}

This proof is adapted from \citet{liberty2017greedy}.

If $g$ is sequence supermodular then,
\begin{align*}
    &g(S) - g(S \ \oplus S' ) \\
    &= \sum_{i = 1}^{|S'|} g(S \ \oplus \left(S'_1, \dots, S'_{i - 1}\right)) - g(S \ \oplus \left(S'_1, \dots, S'_{i}\right)) \\
    &\overset{(a)}{\leq} \sum_{i = 1}^{|S'|} g(S) - g(S \ \oplus  S'_{i}) \\
    &\leq |S'|\max_{i \in [|S'|]} g(S) - g(S \ \oplus  S'_{i})
\end{align*}

(a) by supermodularity


\end{proof}

\subsection{Proof of \Cref{th:suboptgreedy}} \label{sec:proofsuboptgreedy}



\begin{customtheorem}{\ref{th:suboptgreedy}}
    Let $g: \Omega^* \to \mathbb{R}$ be a weakly supermodular function with respect to the optimal solution $O = \argmin_{S \in \Omega^T} g(S)$ with a supermodularity ratio of $\alpha(O)$ and 
    % weak-$\mu$-
    $\mu$-postfix monotonicity with respect to $O$ along the greedy solution $S^T$ of length $T$. If $\phi_T(\alpha) =  \left(1 - \frac{1}{\alpha T}\right)^T$, then
    \begin{equation*}
        g(S^T) \leq  \mu\left(1 - \phi_T(\alpha(O))\right)
        g(O) +  \phi_T(\alpha(O))g(\emptyset)
    \end{equation*}
    % Equivalently, 
    % \begin{equation*}
    %     \frac{g(\emptyset) - g(S^T)}{g(\emptyset) - g(O)} \geq \mu\left(1 - \left(1 - \frac{1}{\alpha T}\right)^T\right) \geq \mu\left(1 - \frac{1}{e^{\alpha}}\right)
    % \end{equation*}
\end{customtheorem}

\begin{proof}

This proof strategy is inspired by \citet{streeter2008online} and \citet{liberty2017greedy}

\begin{align*}
    g(S^t) - \mu g(O) &\overset{(a)}{\leq} g(S^t) - g(S^t\ \oplus O) \\
    &=\sum_{i = 1}^{|O|} g(S^t \oplus \left(o_1, \dots o_{i - 1}\right)) -g(S^t\ \oplus \left(o_1, \dots, o_{i}\right))\\
    &\overset{(b)}{\leq}  \alpha(O)\sum_{i = 1}^{|O|} g(S^t) -g(S^t\ \oplus  o_{i}) \\
    &\leq \alpha(O)|O|\max_{i \in [|O|]} g(S^t) -g(S^t\ \oplus  o_{i}) \\
    &\leq \alpha(O) T g(S^t) -  \alpha(O) T\min_{\omega \in \Omega} g(S^t\ \oplus  \omega) \\
    &= \alpha(O) Tg(S^t) -  \alpha(O) Tg(S^{t + 1})
    % &= Tg(S^t) -  Tg(S^{t + 1}) \\
    % &= T \left(g(S^t) - g(O)\right) - T\left(g(S^{t + 1}) -g(O)\right) \\
    % g(S^{t + 1}) -g(O) &\leq \left(1 - \frac{1}{T}\right)\left(g(S^t) - g(O)\right)
\end{align*}
(a) by $\mu$-postfix monotonicity along $S^T$, i.e., $g(S^t \oplus O) \le \mu g(O)$; (b) by weak supermodularity with respect to $O$ (applied to the prefix $S^t$). 

After rearranging the inequality, we get:

\begin{align*}
    g(S^{t + 1}) &\leq \frac{1}{\alpha(O) T}\left(\mu g(O) +  \left(\alpha(O) T - 1\right)g(S^t)\right) \\
    &= \frac{\mu}{\alpha(O) T}g(O) + \left(1 - \frac{1}{\alpha(O) T}\right)g(S^t)
\end{align*}

When recursively applying this inequality, we get:

\begin{align*}
     g(S^{T}) &\leq \frac{\mu}{\alpha(O) T}g(O)\sum_{i = 0}^{T - 1}\left(1 - \frac{1}{\alpha(O) T}\right)^i + \left(1 - \frac{1}{\alpha(O) T}\right)^Tg(\emptyset) \\
     &= \mu\left(1 - \left(1 - \frac{1}{\alpha(O) T}\right)^T\right)g(O) + \left(1 - \frac{1}{\alpha(O) T}\right)^Tg(\emptyset)
\end{align*}

% \begin{align*}
%     g(S^t) -g(O) &\overset{(a)}{\leq} g(S^t) - \frac{1}{\mu}g(S^t\ \oplus O) \\
%     &= \left(1 - \frac{1}{\mu}\right)g(S^t) + \frac{1}{\mu}\left(g(S^t) - g(S^t\ \oplus O)\right) \\
%     &=\left(1 - \frac{1}{\mu}\right)g(S^t) + \frac{1}{\mu}\sum_{i = 1}^{|O|} g(S^t \oplus \left(o_1, \dots o_{i - 1}\right)) -g(S^t\ \oplus \left(o_1, \dots, o_{i}\right))\\
%     &\overset{(b)}{\leq} \left(1 - \frac{1}{\mu}\right)g(S^t) + \frac{\alpha(O)}{\mu}\sum_{i = 1}^{|O|} g(S^t) -g(S^t\ \oplus  o_{i}) \\
%     &\leq \left(1 - \frac{1}{\mu}\right)g(S^t) + \frac{\alpha(O)|O|}{\mu}\max_{i \in [|O|]} g(S^t) -g(S^t\ \oplus  o_{i}) \\
%     &\leq  \left(1 - \frac{1 - \alpha(O) T}{\mu}\right)g(S^t) -  \frac{\alpha(O) T}{\mu}\min_{\omega \in \Omega} g(S^t\ \oplus  \omega) \\
%     &=  \left(1 - \frac{1 - \alpha(O) T}{\mu}\right)g(S^t) -  \frac{\alpha(O) T}{\mu}g(S^{t + 1})
%     % &= Tg(S^t) -  Tg(S^{t + 1}) \\
%     % &= T \left(g(S^t) - g(O)\right) - T\left(g(S^{t + 1}) -g(O)\right) \\
%     % g(S^{t + 1}) -g(O) &\leq \left(1 - \frac{1}{T}\right)\left(g(S^t) - g(O)\right)
% \end{align*}
% (a) by $\mu$- postfix monotonicity, (b) by supermodularity. 

% After rearranging the inequality, we get:

% \begin{align*}
%     g(S^{t + 1}) &\leq \frac{\mu}{\alpha(O) T}\left(g(O) -  \frac{1 - \alpha(O) T}{\mu}g(S^t)\right) \\
%     &= \frac{\mu}{\alpha(O) T}g(O) + \left(1 - \frac{1}{\alpha(O) T}\right)g(S^t)
% \end{align*}

% When recursively applying this inequality, we get:

% \begin{align*}
%      g(S^{T}) &\leq \frac{\mu}{\alpha(O) T}g(O)\sum_{i = 0}^{T - 1}\left(1 - \frac{1}{\alpha(O) T}\right)^i + \left(1 - \frac{1}{\alpha(O) T}\right)^Tg(\emptyset) \\
%      &= \mu \left(1 - \left(1 - \frac{1}{\alpha(O) T}\right)^T\right)g(O) + \left(1 - \frac{1}{\alpha(O) T}\right)^Tg(\emptyset)
% \end{align*}
\end{proof}

\subsection{Proof of Proposition \ref{th:weaklyalphasupermodular}} \label{sec:proofweaklyalphasupermodular}


% \begin{customproposition}{\ref{th:weaklyalphasupermodular}}
%     If $\|I_N - C_j\mathcal{L}_h^a\| \leq 1$ for all $j \in [K]$, $h$ is (i) prefix monotonically non-increasing and (ii) f is weakly-$\alpha$-supermodular with $$\alpha = \max_{S \in [K]^*: |S| = T}\frac{\left\|I_N - \prod_{t = T}^{1}\left(I_N - C_{S_t}\mathcal{L}_h^a\right)\right\|}{T\left(1 - \max_{t \in [T]}\|I_N - C_{S_t}\mathcal{L}_h^a\|\right) }$$Furthermore, if $I_N - C_j\mathcal{L}_h^a$ is invertible for all $j \in [K]$, $h$ is also postfix monotonically non-increasing.
% \end{customproposition}

\begin{customproposition}{\ref{th:weaklyalphasupermodular}}
     Suppose that for all $j \in [K]$, the error propagation function $I - C_j \circ \mathcal{L}_h^a$ is $\rho_j$-Lipschitz continuous with $\rho_j < 1$, and that $(I - C_j \circ \mathcal{L}_h^a) (0_N) = 0_N$. Then, the function $h$ is weakly supermodular with respect to the optimal solution $O$, with 
     \begin{equation*}
         \alpha(O) = \max\left\{\frac{1}{T - \sum_{i = 1}^{T} \rho_{O_i}^2} , 1\right\}  
     \end{equation*}
     \edit{Furthermore, if the error propagation maps are linear and simultaneously diagonalizable by an orthogonal matrix with eigenvalues of modulus at most one (the setting of \Cref{th:simultaneoussupermodular}), $h$ is also postfix non-increasing; invertibility alone does not suffice (\Cref{sec:proofweaklyalphasupermodular}).}
\end{customproposition}

\begin{proof}
For brevity, we use the notation $\left(g_1 \circ \dots \circ g_T\right)(x) = \circ_{t = 1}^{T} g_t(x)$, where composition is applied from right to left so that $g_T$ acts first. In this proof, we use a few properties of Lipschitz continuous functions:
\begin{itemize}
    \item \textbf{Property 1:} If $g$ is $\rho$-Lipschitz continuous and $g(0) = 0$, then $\|g(x)\|_2 = \|g(x) - 0\|_2 = \|g(x) - g(0)\|_2 \leq \rho \|x - 0\|_2 = \rho\|x\|_2$
    \item \textbf{Property 2:} If $g_1$ and $g_2$ are Lipschitz continuous functions with Lipschitz constants of $\rho_1$ and $\rho_2$ respectively, the Lipschitz constant of $g_1 + g_2$ and $g_1 - g_2$ is $\rho_1 + \rho_2$.
    \item \textbf{Property 3:} If $g_1$ and $g_2$ are Lipschitz continuous functions with Lipschitz constants of $\rho_1$ and $\rho_2$ respectively, the Lipschitz constant of $g_1 \circ g_2$ is $\rho_1\rho_2$.
\end{itemize}

In order to prove weakly-$\alpha$-supermodularity, we must first prove prefix monotonicity.

\textbf{Prefix monotonicity: } Let $S \preceq S'$ where $S' = S \ \oplus N$.
\begin{align*}
    h(S') &= \left \|\circ_{t = |S'|}^1\left(I - C_{S'_t} \circ \mathcal{L}_h^a\right) \left(e^{(0)}_h\right)\right\|^2_2 \\
    &= \left \|\circ_{t = |S'|}^{|S| + 1}\left(I - C_{S'_t}  \circ \mathcal{L}_h^a\right)\circ \circ_{t = |S|}^1\left(I - C_{S'_t}  \circ \mathcal{L}_h^a\right) \left(e^{(0)}_h\right)\right\|^2_2 \\
    &\overset{(a)}{\leq}\left( \prod_{t = |S| + 1 }^{|S'|} \rho_{S'_t}^2 \right)\left\|\circ_{t = |S|}^1\left(I - C_{S_t} \circ \mathcal{L}_h^a\right) \left(e^{(0)}_h\right)\right\|^2_2 \\
    &\leq h(S)
\end{align*}
(a) by Property 1


\textbf{Weak supermodularity: } 
% \paragraph{ } 
We will upper bound $\alpha(O)$ by providing an upper bound for $ h(S) - h(S\ \oplus O) $ and a lower bound for $\sum_{i = 1}^T h(S) - h(S\ \oplus O_i) $. 
\edit{Since $h \ge 0$, $h(S) - h(S\ \oplus O) \le h(S) = \left\| \circ_{t = |S|}^1\left(I - C_{S_t} \circ \mathcal{L}_h^a\right) \left(e^{(0)}_h\right)\right\|^2_2$. (An earlier version of this step used $\|x\|^2 - \|Gx\|^2 \le \|x - Gx\|^2$, which is false, e.g. for $G = \tfrac12 I$; the direct bound is sharper by a factor $4$.)}


To lower bound $\sum_{i = 1}^{|O|} h(S) - h(S\ \oplus O_i) $
\begin{align*}
    \sum_{i = 1}^{|O|} h(S) - h(S\ \oplus O_i) &= \sum_{i = 1}^{|O|} \left \|\circ_{t = |S|}^1\left(I - C_{S_t} \circ \mathcal{L}_h^a\right) \left(e^{(0)}_h\right)\right\|^2_2 - \left \|\left(I - C_{O_i} \circ \mathcal{L}_h^a\right) \circ \circ_{t = |S|}^1\left(I - C_{S_t} \circ\mathcal{L}_h^a\right) \left(e^{(0)}_h\right)\right\|^2_2  \\
    &\geq \sum_{i = 1}^{|O|} \left \|\circ_{t = |S|}^1\left(I - C_{S_t} \circ\mathcal{L}_h^a\right) \left(e^{(0)}_h\right)\right\|^2_2 - \rho_{O_i}^2\left\|\circ_{t = |S|}^1\left(I - C_{S_t}\circ\mathcal{L}_h^a\right) \left(e^{(0)}_h\right)\right\|^2_2 \\
    &= \left \|\circ_{t = |S|}^1\left(I_N - C_{S_t}\circ\mathcal{L}_h^a\right) \left(e^{(0)}_h\right)\right\|^2_2 \left(T - \sum_{i = 1}^{T} \rho_{O_i}^2\right) 
\end{align*}
% \begin{align*}
%     \max_i h(S) - h(S\ \oplus S'_i) &= \left \|\circ_{t = |S|}^1\left(I - C_{S_t} \circ \mathcal{L}_h^a\right) \left(e^{(0)}_h\right)\right\|^2_2 - \min_i \left \|\left(I - C_{S_i'} \circ \mathcal{L}_h^a\right) \circ \circ_{t = |S|}^1\left(I - C_{S_t} \circ\mathcal{L}_h^a\right) \left(e^{(0)}_h\right)\right\|^2_2  \\
%     &\geq \left \|\circ_{t = |S|}^1\left(I - C_{S_t} \circ\mathcal{L}_h^a\right) \left(e^{(0)}_h\right)\right\|^2_2 - \min_i \rho_{S_i'}^2\left\|\circ_{t = |S|}^1\left(I - C_{S_t}\circ\mathcal{L}_h^a\right) \left(e^{(0)}_h\right)\right\|^2_2 \\
%     &= \left(1 - \min_i \rho_{S_i'}^2\right) \left \|\circ_{t = |S|}^1\left(I_N - C_{S_t}\circ\mathcal{L}_h^a\right) \left(e^{(0)}_h\right)\right\|^2_2 
% \end{align*}

Finally,

\begin{align*}
    \frac{h(S) - h(S\ \oplus O)}{\sum_{i = 1}^{T} \left(h(S) - h(S\ \oplus O_i)\right)} &\leq \max\left\{ \frac{\left\| \circ_{t = |S|}^1\left(I - C_{S_t} \circ \mathcal{L}_h^a\right) \left(e^{(0)}_h\right)\right\|^2_2}{\left \|\circ_{t = |S|}^1\left(I_N - C_{S_t}\circ\mathcal{L}_h^a\right) \left(e^{(0)}_h\right)\right\|^2_2 \left(T - \sum_{i = 1}^{T} \rho_{O_i}^2\right) }, 1\right\}\\
    &= \max\left\{\frac{1}{T - \sum_{i = 1}^{T} \rho_{O_i}^2} , 1\right\} 
\end{align*}

\edit{\textbf{Postfix monotonicity.} Let $S' = S \oplus N$ and suppose that $I - C_j\mathcal{L}_h^a = P\Lambda_jP^{\top}$ for all $j$ with a common orthogonal $P$ and $|\lambda_{ji}| \le 1$. By \Cref{th:simultaneousallocation},
\begin{align*}
    h(S') = \sum_{i = 1}^N z_i^2 \prod_{j = 1}^K \lambda_{ji}^{2 m_j(N)} \prod_{j = 1}^K \lambda_{ji}^{2 m_j(S)} \le \sum_{i = 1}^N z_i^2 \prod_{j = 1}^K \lambda_{ji}^{2 m_j(N)} = h(N),
\end{align*}
since every factor $\lambda_{ji}^{2 m_j(S)}$ lies in $[0, 1]$. The shared eigenbasis cannot be replaced by invertibility of the individual maps: for
\begin{align*}
    A = \begin{bmatrix} 0.0674240 & -0.1864331 \\ -0.1473150 & -0.8707252 \end{bmatrix}, \quad
    B = \begin{bmatrix} 0.7575738 & 0.4816284 \\ -0.1369844 & 0.3257283 \end{bmatrix}, \quad
    x = \begin{bmatrix} 0.2812107 \\ -0.5538228 \end{bmatrix},
\end{align*}
both matrices are invertible with spectral norm $0.9$, yet $\|BAx\|_2 = 0.330 > 0.225 = \|Bx\|_2$, i.e., $h((A) \oplus (B)) > h((B))$. (An earlier version of this proof bounded $\|G^{-1}\|$ by the reciprocal of the Lipschitz constant of $G$, which is not valid.)}

\end{proof}

\subsection{Proof of \Cref{th:simultaneousallocation}} \label{sec:proofsimultaneousallocation}

\begin{lemma} \label{th:simultaneousallocation}
     Let $\left(I - C_j \mathcal{L}_{h}^a\right) = P\Lambda_jP^{-1}$ where $P$ is an orthogonal matrix and $\Lambda_j = \mathrm{diag}(\lambda_{j1}, \dots, \lambda_{jN})$. If $P^{-1}e^{(0)}_h = z$, then the following equality holds:
     \begin{equation} \label{eq:simultaneous}
          h(S) = \sum_{i = 1}^N z_i^2 \prod_{j = 1}^K \lambda_{ji}^{2m_j(S)}
          % \qquad\mathrm{s.t}\qquad \sum_{k=1}^{K} m_k = T 
     \end{equation}
     where $m_j(S) = \sum_{t= 1}^{|S|}\1_{S_t = j}$
\end{lemma}

\begin{proof}

\begin{align*}
    h(S) &= \left\|\prod_{t = |S|}^1\left(I_N - C_{S_t} \mathcal{L}_{h}^a\right)  e^{(0)}_h \right\|^2 \\
    &=  \left\|\prod_{t = |S|}^1\left(P\Lambda_{S_t}P^{-1}\right)  e^{(0)}_h \right\|^2 \\
    &=  \left\|P\prod_{t = |S|}^1\Lambda_{S_t}P^{-1}  e^{(0)}_h \right\|^2 \\
    &= \left\|\prod_{t = |S|}^1\Lambda_{S_t}P^{-1}  e^{(0)}_h \right\|^2\\
    &=  \sum_{i = 1}^N z_i^2 \prod_{t = T}^1 \lambda_{S_ti}^2 \\
    &=  \sum_{i = 1}^N z_i^2 \prod_{j = 1}^K \lambda_{ji}^{2m_j(S)} 
    % &= \min_{\mathbf{m} \in\mathbb{N}^{K}}\quad \sum_{j = 1}^N z_i^2 \exp \left\{\sum_{k = 1}^K m_k \log \lambda_{kj}^{2}\right\}
    % \qquad\mathrm{s.t}\qquad \sum_{k=1}^{K} m_k = T
\end{align*}

\end{proof}

\subsection{Proof of \Cref{th:simultaneoussupermodular}} \label{sec:proofsimultaneoussupermodular}
\begin{customproposition}{\ref{th:simultaneoussupermodular}}
    Let $\|I_N - C_j\mathcal{L}_h^a\| \leq 1$ for all $j \in [K]$ and $\left(I - C_j \mathcal{L}_{h}^a\right) = P\Lambda_jP^{-1}$ \edit{with a common orthogonal matrix $P$}. Then, $h$ is supermodular. 
\end{customproposition}

\begin{proof}

Let $S \preceq S'$ where $S' = S \ \oplus B$. By \Cref{th:simultaneousallocation},
\begin{align*}
    h(S) &= \sum_{i = 1}^N z_i^2 \prod_{j = 1}^K \lambda_{ji}^{2m_j(S)} 
\end{align*}
where $m_j(S) = \sum_{t = 1}^{|S|} \1_{S_t = j}$ be the number of times a sequence $S$ calls the solver $j$. Recall that $h$ is considered sequence supermodular if $\forall S', S \in \Omega^*$ such that $S \preceq S'$, it holds that
\begin{align*}
    h(S) - h(S \ \oplus \omega ) \geq h(S') - h(S' \ \oplus \omega )
\end{align*}
% Without loss of generality, let's assume that $\omega$ that will be appended to $S, S'$ is $1$.
\begin{align*}
    h(S) - h(S \ \oplus \omega) &= \sum_{i = 1}^N z_i^2 \prod_{j = 1}^K \lambda_{ji}^{2m_j(S)} - \sum_{i = 1}^N z_i^2 \lambda_{\omega i}^{2}\prod_{k = 1}^K \lambda_{ji}^{2m_j(S)} \\
    &=  \sum_{i = 1}^N \left(1 - \lambda_{\omega i}^2\right) z_i^2 \prod_{j = 1}^K \lambda_{ji}^{2m_j(S)}
\end{align*}
Similarly, 
\begin{align*}
    h(S') - h(S' \ \oplus \omega) &= \sum_{i = 1}^N \left(1 - \lambda_{\omega i}^2\right) z_i^2 \prod_{j = 1}^K \lambda_{ji}^{2m_j(S')} \\
    &\overset{(a)}{=}  \sum_{i = 1}^N \left(1 - \lambda_{\omega i}^2\right) z_i^2 \prod_{j = 1}^K \lambda_{ji}^{2\left(m_j(S) + m_j(B)\right)} \\
    &\overset{(b)}{\leq } \sum_{i = 1}^N \left(1 - \lambda_{\omega i}^2\right) z_i^2 \prod_{j = 1}^K \lambda_{ji}^{2m_j(S)} \\
    &= h(S) - h(S \ \oplus \omega)
\end{align*}


(a) Since $m_j(S') = \sum_{t = 1}^{|S'|} \1_{S_t' = j} = \sum_{t = 1}^{|S|} \1_{S_t' = j} + \sum_{t = |S|}^{|S'|} \1_{S_t' = j} = \sum_{t = 1}^{|S|} \1_{S_t = k} + \sum_{t = 1}^{|B|} 1_{B_t = j} = m_j(S) + m_j(B)$, (b) since $|\lambda_{ji}| \leq 1$ (as $\|I_N - C_j\mathcal{L}_h^a\| \le 1$ and $P$ is orthogonal) and $m_j(B) \geq 0$ for all $j \in [K]$

\end{proof}


\section{Proofs for \Cref{sec:approxgreedy}} \label{sec:proofapproxgreedy}

\subsection{Proof of \Cref{th:routingequivalence}} \label{sec:proofroutingequivalence}

\begin{lemma} \label{th:routingequivalence}

For any set of preconditioning functions $\mathcal{C}$, any discrete operator $\mathcal{L}_h^a$, any router $r$, any $a_h, f_h, u_h^{(t)}, u_h \in \mathcal{A} \times \mathcal{F} \times \mathcal{U} \times \mathcal{U}$, the following equality holds true:
\begin{align*}
     l_{\text{route}}\left(r, a_h, f_h, u^{(t)}_h, u_h\right) = &\sum_{j = 1}^K \sum_{k = 1}^K\left\|\left( I - C_{k} \circ \mathcal{L}_{h}^a\right)\left(u_h - u^{(t)}_h\right)\right\|^2_2\1_{k\neq j}\1_{r(a_h, f_h, u^{(t)}_h) \neq j} \\
     &-\left(K - 2\right)\sum_{j = 1}^K\left\|\left( I - C_{j} \circ \mathcal{L}_{h}^a\right)\left(u_h - u^{(t)}_h\right)\right\|^2_2
\end{align*}
    
\end{lemma}

\begin{proof}
Note that $\sum_{j = 1}^K \1_{r(a_h, f_h, u^{(t)}_h) \neq j}= K - 1$
\begin{align*}
    l_{\text{route}}\left(r, a_h, f_h, u^{(t)}_h, u_h\right) &= \sum_{j = 1}^K \left\|\left( I - C_{j} \circ \mathcal{L}_{h}^a\right)\left(u_h - u^{(t)}_h\right)\right\|^2_2\1_{r(a_h, f_h, u^{(t)}_h) = j} \\
    &= \sum_{j = 1}^K \left\|\left( I - C_{j} \circ \mathcal{L}_{h}^a\right)\left(u_h - u^{(t)}_h\right)\right\|^2_2 - \sum_{j = 1}^K \left\|\left( I - C_{j} \circ \mathcal{L}_{h}^a\right)\left(u_h - u^{(t)}_h\right)\right\|^2_2\1_{r(a_h, f_h, u^{(t)}_h) \neq j} \\
    &= \sum_{j = 1}^K \left\|\left( I - C_{j} \circ \mathcal{L}_{h}^a\right)\left(u_h - u^{(t)}_h\right)\right\|^2_2 - \sum_{j = 1}^K \left\|\left( I - C_{j} \circ \mathcal{L}_{h}^a\right)\left(u_h - u^{(t)}_h\right)\right\|^2_2\1_{r(a_h, f_h, u^{(t)}_h) \neq j} \\
    &+ \left(K - 1\right)\sum_{k = 1}^K\left\|\left( I - C_{k} \circ \mathcal{L}_{h}^a\right)\left(u_h - u^{(t)}_h\right)\right\|^2_2 - \left(K - 1\right)\sum_{k = 1}^K\left\|\left( I - C_{k} \circ \mathcal{L}_{h}^a\right)\left(u_h - u^{(t)}_h\right)\right\|^2_2 \\
    &= \sum_{j = 1}^K \left\|\left( I - C_{j} \circ \mathcal{L}_{h}^a\right)\left(u_h - u^{(t)}_h\right)\right\|^2_2 - \sum_{j = 1}^K \left\|\left( I - C_{j} \circ \mathcal{L}_{h}^a\right)\left(u_h - u^{(t)}_h\right)\right\|^2_2\1_{r(a_h, f_h, u^{(t)}_h) \neq j} \\
    &+ \sum_{j = 1}^K\1_{r(a_h, f_h, u^{(t)}_h) \neq j}\sum_{k = 1}^K\left\|\left( I - C_{k} \circ \mathcal{L}_{h}^a\right)\left(u_h - u^{(t)}_h\right)\right\|^2_2 \\
    &- \left(K - 1\right)\sum_{k = 1}^K\left\|\left( I - C_{k} \circ \mathcal{L}_{h}^a\right)\left(u_h - u^{(t)}_h\right)\right\|^2_2 \\
\end{align*}
\begin{align*}
    &= \sum_{j = 1}^K\left(\sum_{k = 1}^K\left\|\left( I - C_{k} \circ \mathcal{L}_{h}^a\right)\left(u_h - u^{(t)}_h\right)\right\|^2_2  - \left\|\left( I - C_{j} \circ \mathcal{L}_{h}^a\right)\left(u_h - u^{(t)}_h\right)\right\|^2_2\right)\1_{r(a_h, f_h, u^{(t)}_h) \neq j} \\
    &- \left(K - 2\right)\sum_{k = 1}^K\left\|\left( I - C_{k} \circ \mathcal{L}_{h}^a\right)\left(u_h - u^{(t)}_h\right)\right\|^2_2 \\
    &=\sum_{j = 1}^K\sum_{k = 1}^K\left\|\left( I - C_{k} \circ \mathcal{L}_{h}^a\right)\left(u_h - u^{(t)}_h\right)\right\|^2_2\1_{k\neq j}\1_{r(a_h, f_h, u^{(t)}_h) \neq j} \\
    &- \left(K - 2\right)\sum_{k = 1}^K\left\|\left( I - C_{k} \circ \mathcal{L}_{h}^a\right)\left(u_h - u^{(t)}_h\right)\right\|^2_2
 \end{align*}
\end{proof}


\subsection{Proof of \Cref{th:upperbound}} \label{sec:proofupperbound}

\begin{proposition} \label{th:upperbound}
    For any router $r$ defined by $r(a, f, u^{(t)}) = \text{argmax}_{j \in [K]} g_j(a, f, u^{(t)})$, any $a_h \in \mathcal{A}$, $f_h \in \mathcal{F}$, and $u_h^{(t)}, u_h \in \mathcal{U}$, the routing loss $l_{\text{route}}$ satisfies:
    \begin{equation*}
        \log(2)l_{\text{route}}\left(r, a_h, f_h, u^{(t)}_h, u_h\right)\leq \Psi(\mathbf{g}, a_h, f_h, u_h^{(t)}, u_h) 
    \end{equation*}
\end{proposition}


\begin{proof}

By \Cref{th:routingequivalence}, we know that

\begin{align*}
    \log(2)l_{\text{route}}\left(r, a_h, f_h, u^{(t)}_h, u_h\right) = &\log(2)\sum_{j = 1}^K \sum_{k = 1}^K\left\|\left( I - C_{k} \circ \mathcal{L}_{h}^a\right)\left(u_h - u^{(t)}_h\right)\right\|^2_2\1_{k\neq j}\1_{r(a_h, f_h, u^{(t)}_h) \neq j} \\
     &-\log(2)\left(K - 2\right)\sum_{j = 1}^K\left\|\left( I - C_{j} \circ \mathcal{L}_{h}^a\right)\left(u_h - u^{(t)}_h\right)\right\|^2_2 \\
     &\leq \log(2)\sum_{j = 1}^K \sum_{k = 1}^K\left\|\left( I - C_{k} \circ \mathcal{L}_{h}^a\right)\left(u_h - u^{(t)}_h\right)\right\|^2_2\1_{k\neq j}\1_{r(a_h, f_h, u^{(t)}_h) \neq j} \\
     &\overset{(a)}{\leq} -\sum_{j = 1}^K \sum_{k = 1}^K\left\|\left( I - C_{k} \circ \mathcal{L}_{h}^a\right)\left(u_h - u^{(t)}_h\right)\right\|^2_2\1_{k\neq j}\log\left(\frac{\exp\left(g_j(a, f, u^{(t)})\right)}{\sum_{k = 1}^K\exp\left(g_k(a, f, u^{(t)})\right)}\right) \\
     &= \Psi(\mathbf{g}, a_h, f_h, u_h^{(t)}, u_h)
\end{align*}
(a) if $r(a_h, f_h, u^{(t)}_h) \neq j$, $\frac{\exp\left(g_j(a, f, u^{(t)})\right)}{\sum_{k = 1}^K\exp\left(g_k(a, f, u^{(t)})\right)} < 0.5$ which implies that $-\log\left(\frac{\exp\left(g_j(a, f, u^{(t)})\right)}{\sum_{k = 1}^K\exp\left(g_k(a, f, u^{(t)})\right)}\right)  \geq \log(2)\1_{r(a_h, f_h, u^{(t)}_h) \neq j}$

\end{proof}


\subsection{Proof of \Cref{th:consistency}} \label{sec:proofconsistency}

\begin{customtheorem}{\ref{th:consistency}}
    Let $\Tilde{c}_j(a_h, u^{(t)}_h, u_h) < \Bar{E} < \infty$  for all $j \in [K]$. If there exists $j\in [K]$ such that $\Tilde{c}_j(a_h, u^{(t)}_h, u_h) > E_{min} > 0$, then, for any collection of solvers $\{C_j\}_{j= 1}^K$ and linear discrete operator $\mathcal{L}_h^a$, $\Psi$ is Bayes consistent surrogate for $l_{\text{route}}$.
\end{customtheorem}

\begin{proof}

For a given $a_h, f_h$, let $u_h$ be $\mathcal{G}_h\left(a_h, f_h\right)$ where $\mathcal{G}_h$ denotes the solution operator acting on the grid $G_h$. Furthermore, let's consider routers of the form
\begin{equation*}
    r(a, f, u^{(t)}) = \text{argmax}_{j \in [K]} g_j(a, f, u^{(t)})
\end{equation*}

For a given $a_h, f_h, u_h^{(t)} \in \mathcal{A} \times \mathcal{F} \times \mathcal{U}$, let the optimal loss under $l_{\text{route}}$ be $l_{\text{route}}^*\left(a_h, f_h, u^{(t)}_h\right) = \inf_{\Tilde{r}} l_{\text{route}}\left(\tilde{r}, a_h, f_h, u^{(t)}_h, \mathcal{G}_h\left(a_h, f_h\right)\right)$. Similarly, let the optimal loss under $\Psi$ be $\Psi^*\left(a_h, f_h, u^{(t)}_h\right) = \inf_{\Tilde{\mathbf{g}}} \Psi\left(\tilde{\mathbf{g}}, a_h, f_h, u^{(t)}_h, \mathcal{G}_h\left(a_h, f_h\right)\right)$. Let $B_j(a_h, f_h, u^{(t)}_h) = \sum_{k = 1}^K\left\|\left( I - C_{k} \circ \mathcal{L}_{h}^a\right)\left(\mathcal{G}_h\left(a_h, f_h\right) - u^{(t)}_h\right)\right\|^2_2\1_{k\neq j}$.

\begin{align*}
    &l_{\text{route}}\left(r, a_h, f_h, u^{(t)}_h, \mathcal{G}_h\left(a_h, f_h\right)\right) - l_{\text{route}}^*\left(a_h, f_h, u^{(t)}_h\right) \\
    &\overset{(a)}{=} \sum_{j = 1}^K\sum_{k = 1}^K\left\|\left( I - C_{k} \circ \mathcal{L}_{h}^a\right)\left(u_h - u^{(t)}_h\right)\right\|^2_2\1_{k\neq j}\1_{r(a_h, f_h, u^{(t)}_h) \neq j}  -\left(K - 2\right)\sum_{j = 1}^K\left\|\left( I - C_{j} \circ \mathcal{L}_{h}^a\right)\left(u_h - u^{(t)}_h\right)\right\|^2_2 \\
    &\quad- \inf_{\Tilde{r}} \sum_{j = 1}^K\sum_{k = 1}^K\left\|\left( I - C_{k} \circ \mathcal{L}_{h}^a\right)\left(u_h - u^{(t)}_h\right)\right\|^2_2\1_{k\neq j}\1_{\tilde{r}(a_h, f_h, u^{(t)}_h) \neq j}  +\left(K - 2\right)\sum_{j = 1}^K\left\|\left( I - C_{j} \circ \mathcal{L}_{h}^a\right)\left(u_h - u^{(t)}_h\right)\right\|^2_2 \\
    &= \sum_{j = 1}^K B_j(a_h, f_h, u^{(t)}_h)\1_{r(a_h, f_h, u^{(t)}_h) \neq j} - \inf_{\Tilde{r}} \sum_{j = 1}^KB_j(a_h, f_h, u^{(t)}_h)\1_{\tilde{r}(a_h, f_h, u^{(t)}_h) \neq j}  \\
    &= \sum_{k = 1}^K B_k(a_h, f_h, u^{(t)}_h)\left(\sum_{j = 1}^K \frac{B_j(a_h, f_h, u^{(t)}_h)}{\sum_{k = 1}^K B_k(a_h, f_h, u^{(t)}_h)}\1_{r(a_h, f_h, u^{(t)}_h) \neq j} - \inf_{\Tilde{r}} \sum_{j = 1}^K \frac{B_j(a_h, f_h, u^{(t)}_h)}{ \sum_{k = 1}^K B_k(a_h, f_h, u^{(t)}_h)}\1_{\tilde{r}(a_h, f_h, u^{(t)}_h) \neq j}\right)
\end{align*}
(a) by \Cref{th:routingequivalence}

Let $\mathcal{X} = \mathcal{A} \times \mathcal{F} \times \mathcal{U}$ and $\mathcal{Y} = [K]$. Let $\mathcal{P}_{\mathcal{X}}$ denote the degenerate distribution supported at the point $(a_h, f_h, u_h^{(t)})$. We define the conditional distribution - $P(Y = j \mid X = (a_h, f_h, u_h^{(t)})) = \frac{B_j(a_h, f_h, u^{(t)}_h)}{ \sum_{k = 1}^K B_k(a_h, f_h, u^{(t)}_h)}$ for $j \in [K]$. The risk and optimal risk of $0-1$ loss under this distribution can be written as:
\begin{align*}
    \mathcal{R}_{0-1}(r) &= \sum_{j = 1}^K \frac{B_j(a_h, f_h, u^{(t)}_h)}{\sum_{k = 1}^K B_k(a_h, f_h, u^{(t)}_h)}\1_{r(a_h, f_h, u^{(t)}_h) \neq j} \\
    \mathcal{R}_{0-1}^* &= \inf_r \sum_{j = 1}^K \frac{B_j(a_h, f_h, u^{(t)}_h)}{\sum_{k = 1}^K B_k(a_h, f_h, u^{(t)}_h)}\1_{r(a_h, f_h, u^{(t)}_h) \neq j} 
    % \mathcal{R}_{ce}(r) &= \sum_{j = 1}^K \frac{B_j(a_h, f_h, u^{(t)}_h)}{\sum_{k = 1}^K B_k(a_h, f_h, u^{(t)}_h)}\1_{f(a_h, f_h, u^{(t)}_h) \neq j} \\
    % \mathcal{R}_{ce}^* &= \inf_f \sum_{j = 1}^K \frac{B_j(a_h, f_h, u^{(t)}_h)}{\sum_{k = 1}^K B_k(a_h, f_h, u^{(t)}_h)}\1_{f(a_h, f_h, u^{(t)}_h) \neq j}
\end{align*}

If $r(a_h, f_h, u^{(t)}_h) = \argmax_{j \in [k]} g_j(a_h, f_h, u^{(t)}_h)$ for all $x \in \mathcal{X}$, then the risk and optimal risk of cross entropy loss ($l_{ce}(\mathbf{g}, x, y)-\log\left(\frac{\exp\left(g_y(x)\right)}{\sum_{k = 1}^K\exp\left(g_k(x)\right)}\right)$) under this distribution can be written as:

\begin{align*}
    \mathcal{R}_{ce}(\mathbf{g}) &= -\sum_{j = 1}^K \frac{B_j(a_h, f_h, u^{(t)}_h)}{\sum_{k = 1}^K B_k(a_h, f_h, u^{(t)}_h)}\log\left(\frac{\exp\left(g_j(a_h, f_h, u^{(t)}_h)\right)}{\sum_{k = 1}^K\exp\left(g_k(a_h, f_h, u^{(t)}_h)\right)}\right) \\
    \mathcal{R}_{ce}^* &= \inf_{\mathbf{g}} -\sum_{j = 1}^K \frac{B_j(a_h, f_h, u^{(t)}_h)}{\sum_{k = 1}^K B_k(a_h, f_h, u^{(t)}_h)}\log\left(\frac{\exp\left(g_j(a_h, f_h, u^{(t)}_h)\right)}{\sum_{k = 1}^K\exp\left(g_k(a_h, f_h, u^{(t)}_h)\right)}\right) 
\end{align*}

From Theorem 3.1 of \citep{mao2023cross}, $\mathcal{R}_{0-1}(r) - \mathcal{R}_{0-1}^* \leq \Gamma^{-1}\left(\mathcal{R}_{ce}(\mathbf{g}) - \mathcal{R}_{ce}^*\right)$ if $r(a_h, f_h, u^{(t)}_h) = \argmax_{j \in [k]} g_j(a_h, f_h, u^{(t)}_h)$ where $\Gamma(z) = \frac{1 + z}{2}\log(1 + z) + \frac{1 - z}{2}\log(1 - z)$. Then,

\begin{align*}
    &l_{\text{route}}\left(r, a_h, f_h, u^{(t)}_h, \mathcal{G}_h\left(a_h, f_h\right)\right) - l_{\text{route}}^*\left(a_h, f_h, u^{(t)}_h\right) \\
    &= \sum_{k = 1}^K B_k(a_h, f_h, u^{(t)}_h)\left(\sum_{j = 1}^K \frac{B_j(a_h, f_h, u^{(t)}_h)}{\sum_{k = 1}^K B_k(a_h, f_h, u^{(t)}_h)}\1_{r(a_h, f_h, u^{(t)}_h) \neq j} - \inf_{\Tilde{r}} \sum_{j = 1}^K \frac{B_j(a_h, f_h, u^{(t)}_h)}{ \sum_{k = 1}^K B_k(a_h, f_h, u^{(t)}_h)}\1_{\tilde{r}(a_h, f_h, u^{(t)}_h) \neq j}\right) \\
    &\leq \sum_{k = 1}^K B_k(a_h, f_h, u^{(t)}_h)\Gamma^{-1}\left(-\sum_{j = 1}^K \frac{B_j(a_h, f_h, u^{(t)}_h)}{\sum_{k = 1}^K B_k(a_h, f_h, u^{(t)}_h)}\log\left(\frac{\exp\left(g_j(a_h, f_h, u^{(t)}_h)\right)}{\sum_{k = 1}^K\exp\left(g_k(a_h, f_h, u^{(t)}_h)\right)}\right) \right.\\&\quad\quad\quad\quad\quad\quad\quad\quad\quad\quad\quad\quad- \left.\inf_{\mathbf{g}} -\sum_{j = 1}^K \frac{B_j(a_h, f_h, u^{(t)}_h)}{\sum_{k = 1}^K B_k(a_h, f_h, u^{(t)}_h)}\log\left(\frac{\exp\left(g_j(a_h, f_h, u^{(t)}_h)\right)}{\sum_{k = 1}^K\exp\left(g_k(a_h, f_h, u^{(t)}_h)\right)}\right) \right) \\
    &\overset{(a)}{\leq} \Bar{E}K\left(K- 1\right)\Gamma^{-1}\left(-\sum_{j = 1}^K \frac{B_j(a_h, f_h, u^{(t)}_h)}{\sum_{k = 1}^K B_k(a_h, f_h, u^{(t)}_h)}\log\left(\frac{\exp\left(g_j(a_h, f_h, u^{(t)}_h)\right)}{\sum_{k = 1}^K\exp\left(g_k(a_h, f_h, u^{(t)}_h)\right)}\right) \right.\\&\quad\quad\quad\quad\quad\quad\quad\quad\quad\quad\quad\quad- \left.\inf_{\mathbf{g}} -\sum_{j = 1}^K \frac{B_j(a_h, f_h, u^{(t)}_h)}{\sum_{k = 1}^K B_k(a_h, f_h, u^{(t)}_h)}\log\left(\frac{\exp\left(g_j(a_h, f_h, u^{(t)}_h)\right)}{\sum_{k = 1}^K\exp\left(g_k(a_h, f_h, u^{(t)}_h)\right)}\right) \right) \\
    &\overset{(b)}{\leq} \Bar{E}K\left(K- 1\right)\Gamma^{-1}\left(-\sum_{j = 1}^K \frac{B_j(a_h, f_h, u^{(t)}_h)}{(K - 1)E_{min}}\log\left(\frac{\exp\left(g_j(a_h, f_h, u^{(t)}_h)\right)}{\sum_{k = 1}^K\exp\left(g_k(a_h, f_h, u^{(t)}_h)\right)}\right) \right.\\&\quad\quad\quad\quad\quad\quad\quad\quad\quad\quad\quad\quad- \left.\inf_{\mathbf{g}} -\sum_{j = 1}^K \frac{B_j(a_h, f_h, u^{(t)}_h)}{(K - 1)E_{min}}\log\left(\frac{\exp\left(g_j(a_h, f_h, u^{(t)}_h)\right)}{\sum_{k = 1}^K\exp\left(g_k(a_h, f_h, u^{(t)}_h)\right)}\right) \right) \\
    &= \Bar{E}K\left(K- 1\right)\Gamma^{-1}\left(\frac{\Psi\left(\mathbf{g}, a_h, f_h, u^{(t)}_h, \mathcal{G}_h\left(a_h, f_h\right)\right)- \Psi^*\left(a_h, f_h, u^{(t)}_h\right)}{(K - 1)E_{min}} \right)
\end{align*}
(a) since $\left\|\left( I - C_{j} \circ \mathcal{L}_{h}^a\right)\left(e^{(t)}_h\right)\right\|^2_2 < \Bar{E}$ for all $j \in [K]$, (b) since $\Gamma^{-1}$ is non-decreasing and $\exists j \in [K]$ such that $\left\|\left( I - C_{j} \circ \mathcal{L}_{h}^a\right)\left(e^{(t)}_h\right)\right\|^2_2 > E_{min}$

Finally,

\begin{align*}
    &\lim_{n \to \infty }\mathcal{R}_{\text{route}}\left(r_n\right) - \mathcal{R}_{\text{route}}^* \\
    &\overset{(a)}{=} \lim_{n \to \infty } \mathbb{E}_{a_h, f_h \sim \mathcal{P}_{\mathcal{A} \times \mathcal{F}}}\left[l_{\text{route}}\left(r_n, a_h, f_h, u^{(t)}_h, \mathcal{G}_h\left(a_h, f_h\right)\right) - l_{\text{route}}^*\left(a_h, f_h, u^{(t)}_h\right)\right] \\
    &\leq \lim_{n \to \infty } \mathbb{E}_{a_h, f_h \sim \mathcal{P}_{\mathcal{A} \times \mathcal{F}}}\left[\Bar{E}K\left(K- 1\right)\Gamma^{-1}\left(\frac{\Psi\left(\tilde{\mathbf{g}}, a_h, f_h, u^{(t)}_h, \mathcal{G}_h\left(a_h, f_h\right)\right)- \Psi^*\left(a_h, f_h, u^{(t)}_h\right)}{(K - 1)E_{min}} \right)\right] \\
    &\overset{(b)}{\leq} \lim_{n \to \infty } \Bar{E}K\left(K- 1\right)\Gamma^{-1}\left(\frac{\mathbb{E}_{a_h, f_h \sim \mathcal{P}_{\mathcal{A} \times \mathcal{F}}}\left[\Psi\left(\mathbf{g}_n, a_h, f_h, u^{(t)}_h, \mathcal{G}_h\left(a_h, f_h\right)\right)- \Psi^*\left(a_h, f_h, u^{(t)}_h\right) \right]}{(K - 1)E_{min}}\right) \\
\end{align*}
\begin{align*}
    &= \lim_{n \to \infty } \Bar{E}K\left(K- 1\right)\Gamma^{-1}\left(\frac{\mathcal{R}_{\Psi}\left(\mathbf{g}_n\right) - \mathcal{R}_{\Psi}^*}{(K - 1)E_{min}}\right) \\
    &\overset{(c)}{=} \Bar{E}K\left(K- 1\right)\Gamma^{-1}\left(\frac{\lim_{n \to \infty } \mathcal{R}_{\Psi}\left(\mathbf{g}_n\right) - \mathcal{R}_{\Psi}^*}{(K - 1)E_{min}}\right) \\
    &= \Bar{E}K\left(K- 1\right)\Gamma^{-1}\left(0\right) \\
    &\overset{(d)}{=} 0
\end{align*}

(a) $\mathcal{R}_{\text{route}}^*  = \mathbb{E}_{a_h, f_h \sim \mathcal{P}_{\mathcal{A} \times \mathcal{F}}}\left[l_{\text{route}}^*\left(a_h, f_h, u^{(t)}_h\right)\right]$ since the infimum is taken over all measurable functions, (b) by Jensen's inequality since $\Gamma^{-1}$ is concave, (c) by continuity of $\Gamma^{-1}$ at $0$, (d) $\Gamma^{-1}(0) = 0$

% \begin{align*}
%     \mathcal{R}_{\Psi}\left(\mathbf{g}_n\right) = \mathbb{E}_{a, f \sim \mathcal{P}_{\mathcal{A} \times \mathcal{F}}}\left[\Psi\left(\mathbf{g}, a, f, u^{(t)}, u\right)\right]
% \end{align*}


\end{proof}



\section{Training Details} \label{sec:training details}

\edit{The periodic experiments of \Cref{sec:experiments} use the corrector and the lightweight router described in \Cref{sec:wallclock_protocol}; the data distribution below is shared by all experiments, and the DeepONet and LSTM router described afterwards are those of the Dirichlet experiment in \Cref{sec:dirichlet}.}

Data for both DeepONet and the routers is sampled from a zero-mean Hierarchical Gaussian Random Field on the periodic domain with covariance operator $\alpha(- \Delta + \beta I)^{-\gamma}$, where $\alpha, \beta, \gamma$ are sampled as:
\begin{align*}
    \alpha &\sim \text{Log-Uniform}(0.01, 100)\\
    \beta &\sim \text{Log-Uniform}(0.1, 1000) \\
    \gamma &\sim \text{Uniform}\left(\{0.5, 1.0, 1.5, 2.0, 2.5, 3.0, 4.0\}\right)
\end{align*}

Given $(\alpha, \beta, \gamma)$, samples are generated in the Fourier space. For each non-zero frequency mode $k$, we draw an independent complex coefficient from a Gaussian distribution with mean 0 and variance $\alpha(4\pi^2 \|k\|^2_2 + \beta)^{-\gamma}$. We enforce a Hermitian symmetry to obtain a real-valued field, set the zero-frequency (DC) mode to $0$ to obtain a zero mean field, and apply the inverse Discrete Fourier Transform to obtain the field in physical space. 

For each sample, we compute reference solutions with a least squares solver and treat them as ground truth. 

This data is used to trained our DeepONet models and LSTM routers. All the models were implemented using PyTorch and all the models were trained on one Nvidia A40 GPU.

\Cref{tab:hpsettings} contains all hyperparameter details for the DeepONet. DeepONet took 30 minutes to train. We then use the model with the best validation loss. 

\begin{table}[H]
    \centering
    \caption{Hyperparameter settings for DeepONet}
    \begin{tabular}{lr}
        \hline
        Hyperparameter & Value\\
        \hline \hline
        Learning rate & 1e-3 \\
        Branch Dimension & 128 \\
        Hidden dimension for branch net & 256 \\
        No. of hidden layers in branch net & 4\\
        Hidden dimension for trunk net & 256 \\
        No. of hidden layers in trunk net & 4\\
        Gradient Clipping Norm & 1.0  \\
        Weight Decay & 0.005 \\
        Batch size & 256 \\
        Training samples & 10000 \\
        Validation samples & 2000 \\
        Epochs & 1000 \\
        % Hidden layers & 1 & 3 & 1  \\
        % Hidden units & 8 & 128 & 8   \\
        % Dropout & 0.2 & 0.4 & 0.2   \\
        % Epochs & 200 & 100 & 200  \\
        % Training Samples & 2400 & 800 & 2400  \\
        % Gradient Clipping Norm & 1.0 & 1.0 & 1.0  \\
        % Weight Decay & 0.005 & 0.001 & 0.005  \\
        \hline
    \end{tabular}
    \label{tab:hpsettings}
\end{table}


The routers are LSTM models. Along with the inputs specified in \Cref{sec:approxgreedy}, namely $a_h, f_h, u_h^{(t)}$, we also supply the routers with the iteration index $t$, current residual $f_h - \mathcal{L}_h^au_h^{(t)}$, and the routing decisions from the previous iteration. While the iteration index and residual are deterministic functions of $(a_h, f_h, u_h^{(t)})$ and align with the theory presented in \Cref{sec:approxgreedy}, incorporating previous routing decisions is a mild deviations from the idealized setting. Nevertheless, this modification is motivated by the same considerations that justify the use of scheduled sampling: under deployment, routing decisions are predicted autoregressively and influence future inputs, whereas training under teacher forcing assumes independence from past predictions. Including previous routing decisions bridges the gap induced by exposure bias by providing the model with trajectory-level information. 


All the routers are trained with scheduled sampling. We use a warm-up of $e_{w}$ epochs with teacher-forcing probability $p_{tf}(e) = ss_{\text{start}}$. After the warm-up, the $p_{tf}$ decays geometrically by a factor of $\gamma_{tf} < 1$ per epoch and is floored by $s_{\text{end}}$:
\begin{equation*}
    p_{tf}(e) = \begin{cases}
        ss_{\text{start}} & e \leq e_w\\
        \max(ss_{start}\gamma_{tf}^{e - e_{w}}, ss_{end}) & e > e_w
    \end{cases}
\end{equation*}
At each time step, with probability $p_{tf}(e)$, we feed the teacher-forced greedy iterate; otherwise, we feed the router's own predicted iterate.


Since LSTMs on long rollouts can suffer from exploding/vanishing gradients, we use truncated backpropagation through time (TBPTT) \citep{mozer2013focused,robinson1987utility,werbos1988generalization}: the forward pass unrolls the entire trajectory, but gradients are propagated only through the most recent $w_{\text{bptt}}(e)$ steps at epoch $e$. Hidden states are passed forward between segments, while earlier segments are treated as stop-gradient.

We employ a curriculum learning approach analogous to scheduled sampling. Let $T_{\max}$ be the horizon ($300$ iterations). With a warm-up of $e_w$ epochs,
\begin{equation}
    w_{\text{bptt}}(e)  = 
    \begin{cases}
        w_{\text{start}} & e \leq e_w \\
        \min\left(T_{\max}, w_{\text{start}}\gamma_{\text{bptt}}^{\left\lfloor \frac{e - e_w}{f_{\text{bptt}}}\right\rfloor} \right) & e > e_w
    \end{cases}
\end{equation}
so the window grows geometrically by a factor of $\gamma_{\text{bptt}} > 1$ every $f_{\text{bptt}}$ epochs and is capped at the full trajectory length.

\Cref{tab:routerhpsettings} contains all hyperparameter details for the LSTM routers. The routers for the Paired Solver experiment took a maximum of 4 hours and 30 minutes to train while the routers for the Solver Ensemble experiment took a maximum of 6 hours to train. We then use the model with the best validation loss for testing.  The architecture and dataset size is larger for the Solver ensemble experiments to encourage the model to learn some of the nuanced differences between the classes.

\begin{table}[H]
    \centering
    \caption{Hyperparameter settings for routers in Paired solver and Solver Ensemble Experiments}
    \begin{tabular}{lrr}
        \hline
        Hyperparameter & Paired Solver & Solver Ensemble\\
        \hline \hline
        Learning rate & 1e-3 & 1e-3\\
        Hidden dimension & 256 & 512\\
        No. of hidden layers & 4 & 3\\
        Gradient Clipping Norm & 1.0 & 1.0 \\
        Weight Decay & 0.005 & 0.005 \\
        Batch size & 32 & 32 \\
        Training samples & 256 & 1024\\
        Validation samples & 32 & 128 \\
        Epochs & 200 & 200\\
        $ss_{\text{start}}$ & 1.0 & 1.0\\
        $\gamma_{tf}$ & $0.95$ & $0.95$ \\
        $ss_{\text{end}}$ & 0.0 & 0.0 \\
        $w_{\text{start}}$ & $50$ & $50$ \\
        $\gamma_{\text{bptt}}$ & $1.25$ & $1.25$\\
        $e_w$ & 10 & 10\\
        $f_{\text{bptt}}$ & $1$ & $1$\\
        % Hidden layers & 1 & 3 & 1  \\
        % Hidden units & 8 & 128 & 8   \\
        % Dropout & 0.2 & 0.4 & 0.2   \\
        % Epochs & 200 & 100 & 200  \\
        % Training Samples & 2400 & 800 & 2400  \\
        % Gradient Clipping Norm & 1.0 & 1.0 & 1.0  \\
        % Weight Decay & 0.005 & 0.001 & 0.005  \\
        \hline
    \end{tabular}
    \label{tab:routerhpsettings}
\end{table}



\section{Additional Experimental Results} \label{sec:moreexperimentalresults}

\edit{This appendix keeps the one experiment of the original submission that the wall-clock study of \Cref{sec:wallclock} does not supersede: the paired-solver comparison under Dirichlet boundary conditions, which was run with the original pipeline (a $31\times 31$ grid, a DeepONet with an MLP branch and an LSTM router trained with scheduled sampling; \Cref{sec:training details}). The other results of the original submission---iteration-based comparisons on $31^2$ and $63^2$ periodic grids with the LSTM router, residual and Fourier-mode analyses, and ensemble decision traces---are superseded by the per-equation results, the usage analyses and the ensemble study of \Cref{sec:wallclock} and have been removed.}

\subsection{Results with Dirichlet Boundary Conditions} \label{sec:dirichlet}

We replicated the paired-solver experiment of the original submission for solutions with zero Dirichlet boundary conditions, with the original pipeline; the metrics are the final error and AUC over $300$ iterations, with paired $t$-tests against the learned router. \edit{The dense symmetric-GS reference of that pipeline mis-implemented the SSOR scaling (it added the factor $\omega/(2-\omega)$ to $D^{-1}$ instead of multiplying it), so the SymGS rows below use a convergent but non-standard symmetric sweep; the implementation has since been corrected and the wall-clock study uses the standard SSOR.}

In \Cref{tab:errorcomparison_dirichlet}, we notice that, similar to PDEs with periodic boundary conditions, the learned greedy router closely matches the performance of the true greedy oracle and outperforms its single solver and HINTS baselines with great statistical significance. Despite the learned router matching the performance of the true greedy oracle, we observe that the learned router doesn't seem to provide much improvement over single-solver baseline in the case of SOR (1.5)-related solvers. This may be a result of the deeponet correction always being weaker relative to the SOR (1.5) correction in terms of error reduction. 


\begin{table}[H]
    \centering
    \caption{Final error and AUC of squared $L^2$ error (lower is better). Values are mean ($\pm$ standard deviation) over 128 test instances with Dirichlet boundaries; both are reported in $\times 10^{-3}$. If a standard error is not shown, it is $< 10^{-3}$ in the reported units (raw $< 10^{-6}$). Statistical significance is assessed via paired $t$-tests comparing each baseline (single-solver only and HINTS) against the learned greedy router, using a one-sided alternative that the learned router achieves lower error/AUC. Reported $p$-values correspond to these tests.}
    \resizebox{\textwidth}{!}{%
    \begin{tabular}{ccccccccc}
        \hline \hline
         Equation &  \multicolumn{4}{c}{Poisson} & \multicolumn{4}{c}{ConvDiff} \\
         \hline 
         Methods & $\|e^{(T)}_h\| \times 10^{3}$ & $\|e^{(T)}_h\|$ $p$-value & AUC & AUC $p$-value & $\|e^{(T)}_h\| \times 10^{3}$ & $\|e^{(T)}_h\|$ $p$-value & AUC & AUC $p$-value\\
         \hline
         \multicolumn{9}{c}{Jacobi-related solvers}\\
         \hline
Jacobi Only & 0.704 (2.151) & 0.001 & 0.703 (1.929) & $<10^{-3}$ & 0.003 (0.006) & 0.057 & 0.169 (0.430) & $<10^{-3}$ \\
HINTS-Jacobi & 2.674 (7.315) & $<10^{-3}$ & 0.869 (2.396) & $<10^{-3}$ & 0.474 (0.001) & $<10^{-3}$ & 0.199 (0.254) & $<10^{-3}$ \\
Learned Greedy-Jacobi & 0.243 (0.851) & - & 0.236 (0.719) & - & 0.003 (0.006) & - & 0.076 (0.162) & - \\
True-Greedy-Jacobi & 0.042 (0.047) & - & 0.099 (0.252) & - & 0.003 (0.007) & - & 0.042 (0.089) & - \\
\hline
         \multicolumn{9}{c}{GS-related solvers}\\
         \hline
GS only & 0.138 (0.431) & 0.001 & 0.408 (1.139) & $<10^{-3}$ & 0.003 (0.006) & 0.028 & 0.057 (0.146) & $<10^{-3}$ \\
HINTS-GS & 2.668 (7.316) & $<10^{-3}$ & 0.841 (2.360) & $<10^{-3}$ & 0.470 (0.001) & $<10^{-3}$ & 0.115 (0.137) & $<10^{-3}$ \\
Learned Greedy-GS & 0.032 (0.158) & - & 0.115 (0.433) & - & 0.003 (0.006) & - & 0.024 (0.048) & - \\
True-Greedy-GS & 0.010 (0.015) & - & 0.063 (0.166) & - & 0.003 (0.006) & - & 0.017 (0.036) & - \\
\hline
         \multicolumn{9}{c}{SymGS-related solvers}\\
         \hline
SymGS only & 2.924 (8.865) & 0.01 & 0.926 (2.740) & 0.006 & 1.331 (3.687) & 0.004 & 0.413 (1.133) & 0.004 \\
HINTS-SymGS & 2.748 (7.375) & 0.008 & 0.856 (2.427) & 0.008 & 0.917 (1.015) & 0.006 & 0.332 (0.697) & 0.003 \\
Learned Greedy-SymGS & 1.333 (5.715) & - & 0.416 (1.722) & - & 0.628 (1.468) & - & 0.196 (0.449) & - \\
True-Greedy-SymGS & 0.342 (0.906) & - & 0.147 (0.375) & - & 0.184 (0.323) & - & 0.067 (0.118) & - \\
\hline
         \multicolumn{9}{c}{Jacobi (0.67)-related solvers}\\
         \hline
Jacobi (0.67) only & 1.265 (3.768) & $<10^{-3}$ & 0.910 (2.468) & $<10^{-3}$ & 0.008 (0.021) & 0.003 & 0.253 (0.642) & $<10^{-3}$ \\
HINTS-Jacobi (0.67) & 2.686 (7.311) & $<10^{-3}$ & 0.886 (2.412) & $<10^{-3}$ & 0.487 (0.002) & $<10^{-3}$ & 0.220 (0.279) & $<10^{-3}$ \\
Learned Greedy-Jacobi (0.67) & 0.134 (0.275) & - & 0.215 (0.660) & - & 0.004 (0.010) & - & 0.127 (0.393) & - \\
True-Greedy-Jacobi (0.67) & 0.096 (0.220) & - & 0.145 (0.408) & - & 0.004 (0.013) & - & 0.061 (0.131) & - \\
\hline
         \multicolumn{9}{c}{SOR (1.5)-related solvers}\\
         \hline
SOR (1.5) only & 0.004 (0.011) & 0.012 & 0.145 (0.412) & $<10^{-3}$ & 0.003 (0.006) & 1.0 & 0.006 (0.014) & 1.0 \\
HINTS-SOR (1.5) & 2.650 (7.321) & $<10^{-3}$ & 0.791 (2.295) & $<10^{-3}$ & 0.468 (0.001) & $<10^{-3}$ & 0.014 (0.014) & $<10^{-3}$ \\
Learned Greedy-SOR (1.5) & 0.004 (0.011) & - & 0.039 (0.091) & - & 0.003 (0.006) & - & 0.006 (0.014) & - \\
True-Greedy-SOR (1.5) & 0.004 (0.011) & - & 0.028 (0.070) & - & 0.003 (0.006) & - & 0.006 (0.014) & - \\
\hline \hline 
\end{tabular}
}
    \label{tab:errorcomparison_dirichlet}
\end{table}

{\color{red}% ---------------------------------------------------------------------------
% Cost-aware wall-clock study (Sept 2026 revision). Everything in this file is
% new material and is rendered in red through the enclosing colour group.
% Numbers are pulled from costaware_tables.tex (generated by make_tables.py).
% ---------------------------------------------------------------------------
\section{Cost-Aware Routing and Wall-Clock Evaluation} \label{sec:wallclock}

The experiments of \Cref{sec:experiments} compare methods per iteration, which favors methods whose iterations are expensive. This appendix evaluates the end-to-end \emph{wall-clock time to a target accuracy} of the learned router against HINTS, the classical solver alone and the classical methods a practitioner would reach for on these problems (geometric multigrid and preconditioned Krylov methods), for all five classical solvers and three PDEs, on $128^2$, $256^2$ and $512^2$ grids, using the cost-aware instantiation of the greedy rule introduced at the end of \Cref{sec:greedy}. Besides the Poisson and convection--diffusion equations of \Cref{sec:experiments} we add \emph{anisotropic diffusion}, $-\epsilon\,\partial_{x_1}^2 u - \partial_{x_2}^2 u = f$ with $\epsilon = 10^{-2}$ (diffusion tensor $\mathrm{diag}(\epsilon, 1)$), on the same periodic domain with the same forcing distribution. It is the standard model problem on which point-relaxation methods fail to smooth: an error component oscillatory in $x_1$ and smooth in $x_2$ is damped by only $(1-\epsilon)/(1+\epsilon) \approx 0.98$ per Jacobi sweep, and geometric multigrid with point smoothing loses its grid-independent convergence (line relaxation or semi-coarsening would be needed). \Cref{sec:wallclock_protocol} describes the protocol; \Cref{sec:wc_poisson,sec:wc_conv,sec:wc_aniso} collect all results per equation (every grid, every tolerance, every baseline, with significance tests and seed trials); \Cref{sec:wallclock_ens} studies ensembles of one, two and three classical solvers; \Cref{sec:wallclock_overheads} accounts for overheads and training cost; and \Cref{sec:wallclock_assumptions} verifies numerically that the assumptions of \Cref{th:suboptgreedy}, Propositions~\ref{th:weaklyalphasupermodular}--\ref{th:simultaneoussupermodular} and \Cref{th:consistency} hold in these settings.

\subsection{Protocol} \label{sec:wallclock_protocol}

\paragraph{Fast, matched implementations.} A wall-clock comparison is meaningful only if the classical baseline is not handicapped by implementation. All solvers are therefore re-implemented for the periodic constant-coefficient setting as $O(N^2)$ operations: Jacobi and damped Jacobi as 5-point stencil sweeps, Gauss--Seidel, SOR and symmetric GS via cached sparse triangular factorizations, geometric multigrid as V(2,2)-cycles with lexicographic Gauss--Seidel smoothing, full-weighting restriction, bilinear prolongation, rediscretised coarse operators and a direct solve on the $32\times 32$ coarsest grid (error contraction $\approx 0.03$ per cycle on the isotropic equations). Their one-step iterates match the dense implementations used in the rest of the paper to $10^{-12}$ in double precision. Reference solutions are computed by FFT diagonalization of the same discrete operator, i.e., they are exact solutions of the discrete system, so true errors are available for evaluation at no charged cost; the FFT solve is used only as this reference (it is a direct method that exists because the test problems have constant coefficients and periodic boundaries, and is not a competitor of the iterative methods compared here). All timings are single-threaded CPU measurements (Apple M4 Pro, numpy/scipy/PyTorch, one thread) on an otherwise idle machine.

\paragraph{Corrector.} The neural corrector is a DeepONet applied to the residual $r = f_h - \mathcal{L}_h^a u^{(t)}_h$ and evaluated on a $64\times 64$ sensor grid, i.e., on the coarser grid obtained by band-limited (FFT) restriction to the modes $|k_1|,|k_2| \le 31$ at every grid size (for the anisotropic equation the sensor grid is $N\times N/4$, i.e., full resolution along $x_1$, so that the band contains every mode the point smoothers cannot damp; the sensor grid is the only PDE-specific hyperparameter of the pipeline); its output is prolongated back to the fine grid by exact trigonometric interpolation, so the correction is band-limited by construction and never injects error into the high-frequency band that the classical smoother owns (\Cref{fig:ca_predictions}). The trunk net is the fixed orthonormal Fourier basis of that band ($p = 3969$ basis functions, as in POD-DeepONet \citep{lu2022comprehensive}), and the branch net is a linear layer on the $4096$ sensor values; it is fit by least squares on $32{,}000$--$64{,}000$ exact (residual, error) pairs of the discrete operator drawn from a mixture of spectral shapes (forcing-like, residual-like, white and randomly tilted spectra), and reaches a validation relative error of $\approx 2\times 10^{-6}$ on every mode of the band. The corrector is applied scale-equivariantly, $C_{\mathrm{NO}}(r) = \|Rr\|\,\mathcal{G}_\theta(Rr/\|Rr\|)$, which preserves $C_{\mathrm{NO}}(0) = 0$ (Proposition~\ref{th:weaklyalphasupermodular}). Its cost is dominated by a $4096\times 4096$ matrix--vector product and is therefore almost independent of $N$ ($0.65$\,ms at $128^2$, $3.6$\,ms at $512^2$, \Cref{tab:ca_overheads}), i.e., about $12$ Jacobi or $3$ GS sweeps at $128^2$ and about $2.5$ Jacobi sweeps at $512^2$ (\Cref{tab:ca_costs_poisson,tab:ca_costs_conv,tab:ca_costs_aniso}). The same corrector is used by HINTS and by all routing policies.

\paragraph{Cost-equalised macro-actions.} Per-iteration costs $c_j$ (residual evaluation plus the update) are measured once per pairing and grid on the benchmark machine (minimum over seven repeated blocks of the block median, cached), the unit of cost $u$ is one corrector call, and every operation is applied $m_j = \max\{1,\mathrm{round}(u/c_j)\}$ times per decision; for the five stationary solvers this equalises costs to within $8\%$. An operation dearer than the unit (a multigrid cycle) is applied once and compared per unit of cost, $\|e_j\|^{u/(m_j c_j)}$; the surrogate weights $\tilde c_j$ use the same exponent. Ensembles reuse the cached pairwise costs of their members, so every member has the same macro-action size in the ensemble as in its pairwise run. The cost-aware oracle runs \Cref{alg:greedy} on these macro-actions with the true error; the \emph{greedy oracle} rows run the cost-agnostic \Cref{alg:greedy} on single operations, as in the main text.

\paragraph{Lightweight router.} The learned router selects macro-actions from $6 + K$ scalar features, all $O(1)$ given the residual norm that any stopping test computes: the log relative residual, its per-iteration change over the last macro-action and an exponential moving average thereof, the log epoch index, a one-hot encoding of the previous macro-action, and the log counts of corrector calls so far and of iterations since the last call. It is a two-hidden-layer MLP (64 units) evaluated in numpy, costing $\approx 5\,\mu$s per decision independently of $N$. It is trained with the surrogate of \Cref{section:surrogate}, with the per-state weights $\tilde c_k$ normalised by $\sum_k \tilde c_k$ (which does not change the Bayes decision), on states visited by oracle rollouts with $\varepsilon$-exploration (probability $0.1$), followed by two to three DAgger rounds \citep{ross2011reduction} in which the current router's own rollouts are relabelled by the oracle---the batch analogue of the scheduled sampling used in \Cref{sec:training details}. Training uses $32$--$128$ instances per round (fewer on the larger grids) drawn from the same GRF distribution as the test set but with a different seed; the router of every pairing is trained in seconds to a few minutes. Ensemble routers, which must additionally learn to choose among classical members, use a larger imitation budget ($256$ rollouts at $128^2$ and $128$ at $256^2$, six DAgger rounds, hidden width $128$, $300$ epochs): with the pairwise recipe the router over $\{\text{Jacobi}, \text{Jacobi (0.67)}\}$ imitated its oracle imperfectly at tight tolerances (convection--diffusion, $128^2$, $\varepsilon = 10^{-8}$: $3.34$ against the oracle's $2.67$ work-unit ms), which the larger budget removes ($2.70$). At $512^2$ the routers are trained from $32$ oracle rollouts of at most $400$ iterations (stopping error $10^{-8}$); for the convection--diffusion Jacobi pairing, where the first such router missed $h^2$ within the cap on $4$ of $16$ instances, the reported router uses $64$ rollouts of at most $800$ iterations---the only per-cell change to the recipe.

\paragraph{Baselines and metric.} HINTS applies the corrector every $\tau$-th iteration; we report $\tau = 25$ (the schedule used in the main text and in \citet{zhang2024blending}) and the best of $\tau \in \{5, 10, 25, 50\}$ selected \emph{per cell and per tolerance} on the test instances themselves (an optimistic baseline). The classical baselines without corrector run on the same stencil with the same accounting: multigrid V(2,2) alone; conjugate gradients, CG preconditioned by symmetric GS and by a symmetric V-cycle for Poisson; BiCGSTAB, multigrid-preconditioned BiCGSTAB and restarted GMRES(20) for convection--diffusion (GMRES is accounted at the granularity of its restart cycles and was not run at $512^2$, where it was two orders of magnitude slower than the router on the smaller grids). They are not run for anisotropic diffusion, where point-smoothed multigrid loses its grid-independent contraction ($\approx 0.85$ per cycle) and would not be informative. For each test instance ($64$ at $128^2$, $32$ at $256^2$, $16$ at $512^2$) we record, after every iteration, the true relative error $\|e^{(t)}_h\|/\|u_h\|$ (untimed, benchmark-only) and the cumulative charged wall-clock time of a replay that executes only the chosen operations (residual, update, and---for the learned router---every feature and decision computation; oracle decisions are not charged). We report the median first time at which the error drops below a tolerance $\varepsilon$, paired per-instance median speedups, and one-sided Wilcoxon signed-rank tests on the paired log time ratios (alternative: the router is faster); a run that does not reach $\varepsilon$ within the iteration cap ($60{,}000$ at $128^2$, $40{,}000$ at $256^2$, $8{,}000$ at $512^2$) is marked $\dagger$, reported as a lower bound when it is the majority, counted against the router when it is the router's, and entered at the time it spent up to the cap on the baseline side of the tests (a lower bound, hence conservative). We highlight $\varepsilon = h^2$: the discretization error of the second-order scheme scales as $O(h^2)$, so driving the algebraic error far below $h^2$ does not improve the computed solution; we also report $10^{-3}$ and $10^{-8}$, and the iteration-based metrics of the main text (final relative error and AUC $=\sum_{t=1}^{T}\|e^{(t)}_h\|/\|u_h\|$ over $T = \caT$ iterations) on the $128^2$ grid. Every router was additionally retrained from scratch with five independent seeds (different training instances, exploration and initialisation; the corrector and the test set are fixed) and re-run on the test set; because those trials ran in separate sessions, whose live timings are not comparable with the main run's, they are compared in \emph{work units} (measured per-iteration costs times executed operations, plus $5\,\mu$s per decision), which are determined by the decision sequence alone.

\begin{figure}[H]
\color{red}
    \centering
    \newfig[0.93\linewidth]{neurips_images/ca_deeponet_predictions.png}
    \caption{One-shot corrector predictions on the $128\times128$ grid (two test instances per equation): forcing, exact discrete solution, the DeepONet correction applied to the initial residual ($u^{(0)}_h = 0$), and the remaining error, which is confined to the modes outside the corrector's band (the relative error of the prediction is $10^{-6}$--$10^{-5}$ on the isotropic equations and $10^{-6}$ on anisotropic diffusion, whose solution is dominated by modes that are smooth in $x_2$).}
    \label{fig:ca_predictions}
\end{figure}

\begin{figure}[H]
\color{red}
    \centering
    \newfig[0.98\linewidth]{neurips_images/ca_convergence.png}
    \caption{Relative error versus charged wall-clock time on the $128\times128$ grid for one representative test instance of every pairing (columns) and equation (rows); the dashed line is $\varepsilon = h^2$. The learned router and the cost-aware oracle are indistinguishable.}
    \label{fig:ca_convergence}
\end{figure}

\subsection{Poisson} \label{sec:wc_poisson}

\Cref{tab:ca_time_poisson} collects the median time to $10^{-3}$, $h^2$ and $10^{-8}$ for every pairing, every policy and the classical baselines on all three grids; \Cref{tab:ca_speed_poisson} the paired speedups of the learned router with their $p$-values; \Cref{fig:ca_usage_poisson,tab:ca_usage_poisson} where the iterations go; \Cref{tab:ca_costs_poisson} the measured costs and macro-action sizes; \Cref{tab:ca_auc_poisson} the iteration-based metrics; and \Cref{tab:ca_seeds_poisson} the seed trials.

\paragraph{Uniform gains and their mechanism.} On the $128^2$ grid the learned router is the fastest deployable method at $\varepsilon = h^2$ in every pairing: hundreds to thousands of times faster than the solver alone, several times faster than HINTS ($\tau{=}25$) and still faster than the best fixed period chosen per cell and per tolerance on the test set, with all paired $p < 0.01$ (\Cref{tab:ca_speed_poisson}); at $10^{-8}$ the margins over the best fixed period shrink but remain significant in every pairing, and the best period itself varies across pairings and tolerances ($\tau{=}5$ at truncation level, $\tau \in \{5, 10, 50\}$ at $10^{-8}$), so no single fixed schedule is competitive everywhere. \Cref{tab:ca_usage_poisson} explains the numbers: with a corrector that removes the smooth band to $10^{-6}$ in one call, the cost-aware oracle and the router call it first and reach $h^2$ after one or two further sweeps, i.e., in about two iterations, whereas HINTS reaches the same accuracy only at its first scheduled call (iteration $\tau$), having spent $\tau{-}1$ sweeps on an error that is almost entirely smooth. Below $h^2$ the remaining error is the high-frequency band that the corrector cannot see; the router therefore stops calling it (one to three calls in total down to $10^{-8}$; \Cref{fig:ca_usage_poisson}), whereas a fixed schedule keeps paying for a call every $\tau$ iterations. The cost-agnostic greedy oracle (\Cref{alg:greedy} on single operations) makes the same decisions as its cost-aware counterpart at $128^2$, because after the first call the corrector is never better than a single sweep on the remaining high-frequency error; the learned router reproduces the cost-aware oracle's decisions almost exactly (same median iteration and call counts, median times within a few percent, the difference being its own decision cost and timer noise), and the five retrained routers reproduce the main router's decision sequence to $h^2$ on essentially every test instance (\Cref{tab:ca_seeds_poisson}).

\paragraph{Classical baselines and larger grids.} Against the classical methods without corrector (bottom block of \Cref{tab:ca_time_poisson}), the router with its best stationary pairing reaches truncation-level accuracy several times faster than multigrid alone (three V-cycles) and than multigrid-preconditioned CG at $128^2$ and $256^2$, and an order of magnitude faster than unpreconditioned or SymGS-preconditioned CG. The reason is the corrector's accuracy on the smooth band: one call removes the modes $|k|\le 31$ to $10^{-6}$ at the price of a single $4096\times4096$ matrix--vector product, whereas a V(2,2)-cycle contracts the whole spectrum by $\approx 0.03$ at the price of four Gauss--Seidel sweeps on the fine grid plus the coarse-grid work. Routing over $\{\mathrm{NO}, \text{multigrid}\}$ is itself about twice as fast as multigrid alone on every grid: the oracle and the router call the corrector once and let the V-cycle remove the rest, which is the case in which the per-unit-cost comparison of \Cref{sec:greedy} matters (the cycle is dearer than the unit). On $256^2$ the router remains the fastest iterative method in every pairing (\caSpHintsBMin--\caSpHintsBMax\ over HINTS, \caSpBestBMin--\caSpBestBMax\ over the best period across the three equations). At $512^2$ the margins shrink but the ordering holds for the stationary pairings, and the $\{\mathrm{NO}, \text{multigrid}\}$ router is the fastest iterative method we measured at this size. With Jacobi and Gauss--Seidel the gains at $512^2$ are small because the corrector's fixed band $|k|\le 31$ leaves the smoother all modes $k \ge 32$, whose slowest member contracts by only $0.98$ per Jacobi sweep at $N=512$ (against $0.71$ at $N=128$): every policy then needs the same $\approx 80$ sweeps after a single corrector call, the wall-clock time is dominated by sweeps, and the schedule barely matters (with Jacobi the router misses $h^2$ within the $8{,}000$-iteration cap on $2$ of $16$ instances, marked in the table). This is the expected limit of a corrector whose band does not grow with the grid: the hybrid's advantage comes from removing the modes the smoother cannot damp, and a corrector with a band that scales with $N$ (at the price of a dense matrix--vector product that grows like the fourth power of the band) or a multigrid member restores the gains, as the multigrid pairing shows.

\begin{table}[H]
\color{red}
\centering
\caption{Poisson: median wall-clock time to relative error $\varepsilon \in \{10^{-3}, h^2, 10^{-8}\}$ on the $128^2$, $256^2$ and $512^2$ grids ($64$, $32$ and $16$ test instances) for every pairing $\{\mathrm{NO}, \text{solver}\}$ and policy, and for the classical baselines without corrector. ``HINTS (best $\tau$)'' is the best of $\tau\in\{5,10,25,50\}$ chosen per cell on the test set itself (the chosen $\tau$ is the subscript). The greedy oracle is the cost-agnostic \Cref{alg:greedy}, the cost-aware oracle is \Cref{alg:greedy} on cost-equalised macro-actions; both use the true error and are not deployable. Bold marks the router where it is the fastest deployable method of its pairing; $\dagger n$: $n$ instances did not reach the tolerance within the iteration cap ($>$: lower bound).}
\resizebox{\textwidth}{!}{\catimepoisson}
\label{tab:ca_time_poisson}
\end{table}

\begin{table}[H]
\color{red}
\centering
\caption{Poisson: paired median speedup of the learned router at $\varepsilon = h^2$ and $10^{-8}$ over HINTS ($\tau{=}25$), over the best fixed period for that cell, over multigrid alone and over multigrid-preconditioned CG, with the one-sided Wilcoxon $p$-value that the router is faster (or, where its median is larger, that it is slower). Bold: $\ge 1.1\times$ with $p < 0.01$.}
\resizebox{\textwidth}{!}{\caspeedpoisson}
\label{tab:ca_speed_poisson}
\end{table}

\begin{figure}[H]
\color{red}
    \centering
    \newfig[0.98\linewidth]{neurips_images/ca_usage_poisson.png}
    \caption{Poisson: fraction of test runs that execute a corrector call at each iteration, for the learned router, the cost-aware oracle and HINTS ($\tau{=}25$), for every pairing (columns) and grid (rows; all test instances of each grid).}
    \label{fig:ca_usage_poisson}
\end{figure}

\begin{table}[H]
\color{red}
\centering
\caption{Poisson: median iterations (and corrector calls) to $\varepsilon = h^2$ for HINTS ($\tau{=}25$), the cost-aware oracle and the learned router on every grid.}
\resizebox{0.9\textwidth}{!}{\causagepoisson}
\label{tab:ca_usage_poisson}
\end{table}

\begin{table}[H]
\color{red}
\centering
\caption{Poisson: measured cost of one classical iteration $c_{\text{solver}}$ and of one corrector call $c_{\mathrm{NO}}$ (residual evaluation included), and the resulting macro-action size $m_{\text{solver}}$ (sweeps per decision; the corrector is always applied once).}
\resizebox{0.9\textwidth}{!}{\cacostspoisson}
\label{tab:ca_costs_poisson}
\end{table}

\begin{table}[H]
\color{red}
\centering
\caption{Poisson, $128\times128$: iteration-based metrics of the main text---relative error after $T = \caT$ iterations and AUC over the $T$ iterations (mean and standard error over $64$ instances; $p$: one-sided paired $t$-test that the method's AUC exceeds the router's). Bold: router better than both solver-only and HINTS.}
\resizebox{0.6\textwidth}{!}{\caaucpoisson}
\label{tab:ca_auc_poisson}
\end{table}

\begin{table}[H]
\color{red}
\centering
\caption{Poisson, $128\times128$: five independent router-training seeds per pairing---mean and standard deviation over seeds of the median work-unit time to $h^2$, the share of test instances on which the seed router's decision sequence to $h^2$ coincides with the main router's, the range over seeds of the paired speedup over HINTS ($\tau{=}25$), and the number of seeds significantly faster ($p<0.01$) than both HINTS ($\tau{=}25$) and the best fixed period.}
\resizebox{0.85\textwidth}{!}{\caseedspoisson}
\label{tab:ca_seeds_poisson}
\end{table}

\begin{figure}[H]
\color{red}
    \centering
    \newfig[0.98\linewidth]{neurips_images/ca_router_decisions_poisson.png}
    \caption{Poisson, $128\times128$: routing decisions (dark: corrector call) of the learned router (top row) and the cost-aware oracle (bottom row) over the first $60$ iterations of every test instance, for every pairing.}
    \label{fig:ca_decisions_poisson}
\end{figure}

\subsection{Convection--diffusion} \label{sec:wc_conv}

The convection--diffusion results (\Cref{tab:ca_time_conv,tab:ca_speed_conv,fig:ca_usage_conv,tab:ca_usage_conv,tab:ca_costs_conv,tab:ca_auc_conv,tab:ca_seeds_conv}) mirror the Poisson ones: at $128^2$ the router is the fastest deployable method in every pairing at $h^2$ and at $10^{-8}$, reproduces the cost-aware oracle, and is $2$--$3\times$ faster than multigrid alone and than multigrid-preconditioned BiCGSTAB; on $256^2$ the same holds in every pairing; at $512^2$ the router is faster than HINTS and than the best fixed period in the Jacobi, GS and SOR pairings and the $\{\mathrm{NO}, \text{multigrid}\}$ router is again the fastest iterative method, with the same mechanism as on Poisson (the Jacobi and GS pairings are limited by the modes outside the corrector's band: the router, the oracle and HINTS all reach $h^2$ after $84$--$85$ Jacobi iterations). The nonsymmetric operator changes only the details: the best fixed period is $\tau{=}10$ rather than $5$ in several cells, the undamped-Jacobi pairing stalls at the near-Nyquist band below $10^{-6}$ for every policy, and GMRES(20) is two orders of magnitude slower than the router.

\begin{table}[H]
\color{red}
\centering
\caption{Convection--diffusion: median wall-clock time to relative error $\varepsilon \in \{10^{-3}, h^2, 10^{-8}\}$ on the $128^2$, $256^2$ and $512^2$ grids for every pairing and policy, and for the classical baselines without corrector (conventions as in \Cref{tab:ca_time_poisson}).}
\resizebox{\textwidth}{!}{\catimeconv}
\label{tab:ca_time_conv}
\end{table}

\begin{table}[H]
\color{red}
\centering
\caption{Convection--diffusion: paired median speedup of the learned router at $\varepsilon = h^2$ and $10^{-8}$ over HINTS ($\tau{=}25$), the best fixed period, multigrid alone and multigrid-preconditioned BiCGSTAB, with one-sided Wilcoxon $p$-values (conventions as in \Cref{tab:ca_speed_poisson}).}
\resizebox{\textwidth}{!}{\caspeedconv}
\label{tab:ca_speed_conv}
\end{table}

\begin{figure}[H]
\color{red}
    \centering
    \newfig[0.98\linewidth]{neurips_images/ca_usage_convdiff.png}
    \caption{Convection--diffusion: corrector usage per iteration of the learned router, the cost-aware oracle and HINTS ($\tau{=}25$) for every pairing (columns) and grid (rows).}
    \label{fig:ca_usage_conv}
\end{figure}

\begin{table}[H]
\color{red}
\centering
\caption{Convection--diffusion: median iterations (and corrector calls) to $\varepsilon = h^2$ on every grid.}
\resizebox{0.9\textwidth}{!}{\causageconv}
\label{tab:ca_usage_conv}
\end{table}

\begin{table}[H]
\color{red}
\centering
\caption{Convection--diffusion: measured per-iteration costs and macro-action sizes.}
\resizebox{0.9\textwidth}{!}{\cacostsconv}
\label{tab:ca_costs_conv}
\end{table}

\begin{table}[H]
\color{red}
\centering
\caption{Convection--diffusion, $128\times128$: iteration-based metrics (conventions as in \Cref{tab:ca_auc_poisson}).}
\resizebox{0.6\textwidth}{!}{\caaucconv}
\label{tab:ca_auc_conv}
\end{table}

\begin{table}[H]
\color{red}
\centering
\caption{Convection--diffusion, $128\times128$: five-seed router-training trials (conventions as in \Cref{tab:ca_seeds_poisson}).}
\resizebox{0.85\textwidth}{!}{\caseedsconv}
\label{tab:ca_seeds_conv}
\end{table}

\begin{figure}[H]
\color{red}
    \centering
    \newfig[0.98\linewidth]{neurips_images/ca_router_decisions_convdiff.png}
    \caption{Convection--diffusion, $128\times128$: routing decisions (dark: corrector call) of the learned router (top row) and the cost-aware oracle (bottom row) over the first $60$ iterations of every test instance.}
    \label{fig:ca_decisions_conv}
\end{figure}

\subsection{Anisotropic diffusion} \label{sec:wc_aniso}

The anisotropic equation is where the two components of the hybrid are most complementary and where the classical solvers are weakest: the point smoothers cannot damp the $x_1$-oscillatory modes at all (\Cref{tab:ca_time_aniso}: none of the five solvers alone reaches $h^2$ within the iteration cap on either grid), whereas the corrector, whose sensor grid keeps full resolution along $x_1$, removes exactly those modes together with the smooth ones in a single call and leaves the smoother only the $x_2$-oscillatory band, which it damps at the isotropic rate. Consequently the learned router reaches $h^2$ in under $1$\,ms at $128^2$ for every pairing---the same absolute time as on the isotropic equations---while HINTS ($\tau{=}25$) needs $3$--$10$\,ms and the best fixed period $1.2$--$2.4$\,ms; at $256^2$ (Jacobi, GS and SOR pairings, $16$ instances) the router needs $9$--$10$\,ms against $15$--$30$\,ms for HINTS and $10$--$13$\,ms for the best period (\Cref{tab:ca_speed_aniso}). The router again matches the cost-aware oracle, its decisions are identical across the five training seeds (\Cref{tab:ca_seeds_aniso}), and it spends the iterations after the single corrector call on the $x_2$-oscillatory band (\Cref{fig:ca_usage_aniso,tab:ca_usage_aniso}). The multigrid and Krylov baselines are not run for this equation (\Cref{sec:wallclock_protocol}).

\begin{table}[H]
\color{red}
\centering
\caption{Anisotropic diffusion: median wall-clock time to relative error $\varepsilon \in \{10^{-3}, h^2, 10^{-8}\}$ on the $128^2$ ($64$ instances, five pairings) and $256^2$ ($16$ instances; Jacobi, GS and SOR pairings) grids for every policy (conventions as in \Cref{tab:ca_time_poisson}; the solver-only runs are capped at $60{,}000$ and $40{,}000$ iterations and never reach $h^2$).}
\resizebox{0.9\textwidth}{!}{\catimeaniso}
\label{tab:ca_time_aniso}
\end{table}

\begin{table}[H]
\color{red}
\centering
\caption{Anisotropic diffusion: paired median speedup of the learned router at $\varepsilon = h^2$ and $10^{-8}$ over HINTS ($\tau{=}25$) and the best fixed period, with one-sided Wilcoxon $p$-values (no classical baselines without corrector are run for this equation).}
\resizebox{0.9\textwidth}{!}{\caspeedaniso}
\label{tab:ca_speed_aniso}
\end{table}

\begin{figure}[H]
\color{red}
    \centering
    \newfig[0.85\linewidth]{neurips_images/ca_usage_anisodiff.png}
    \caption{Anisotropic diffusion: corrector usage per iteration of the learned router, the cost-aware oracle and HINTS ($\tau{=}25$) for every pairing (columns) and grid (rows).}
    \label{fig:ca_usage_aniso}
\end{figure}

\begin{table}[H]
\color{red}
\centering
\caption{Anisotropic diffusion: median iterations (and corrector calls) to $\varepsilon = h^2$.}
\resizebox{0.75\textwidth}{!}{\causageaniso}
\label{tab:ca_usage_aniso}
\end{table}

\begin{table}[H]
\color{red}
\centering
\caption{Anisotropic diffusion: measured per-iteration costs and macro-action sizes.}
\resizebox{0.75\textwidth}{!}{\cacostsaniso}
\label{tab:ca_costs_aniso}
\end{table}

\begin{table}[H]
\color{red}
\centering
\caption{Anisotropic diffusion, $128\times128$: iteration-based metrics (conventions as in \Cref{tab:ca_auc_poisson}).}
\resizebox{0.6\textwidth}{!}{\caaucaniso}
\label{tab:ca_auc_aniso}
\end{table}

\begin{table}[H]
\color{red}
\centering
\caption{Anisotropic diffusion, $128\times128$: five-seed router-training trials (conventions as in \Cref{tab:ca_seeds_poisson}).}
\resizebox{0.85\textwidth}{!}{\caseedsaniso}
\label{tab:ca_seeds_aniso}
\end{table}

\begin{figure}[H]
\color{red}
    \centering
    \newfig[0.98\linewidth]{neurips_images/ca_router_decisions_anisodiff.png}
    \caption{Anisotropic diffusion, $128\times128$: routing decisions (dark: corrector call) of the learned router (top row) and the cost-aware oracle (bottom row) over the first $60$ iterations of every test instance.}
    \label{fig:ca_decisions_aniso}
\end{figure}

\subsection{Solver ensembles} \label{sec:wallclock_ens}

\paragraph{When can a second classical solver help?} The cost-aware oracle is a per-decision comparison of macro-actions, so a router over $\mathrm{NO}\cup\mathcal{W}$ can only beat the best pairwise router of its members if, along the way to the tolerance, the error passes through regimes in which \emph{different} members give the largest reduction per unit cost. We screened this ceiling exhaustively before training any router: for every subset $\mathcal{W}$ of up to three members of $\{$Jacobi, Jacobi (0.67), GS, SymGS, SOR (1.5), multigrid$\}$ (multigrid on the isotropic equations only), every equation and grid, we ran the cost-aware oracle over $\mathrm{NO}\cup\mathcal{W}$ on the test instances and compared its work-unit time to that of the oracle over each member alone (\Cref{tab:ca_ens_screen}). Two facts emerge. At $\varepsilon = h^2$ no ensemble beats the best single solver anywhere: one corrector call followed by one or two macro-actions of the cheapest smoother already reaches truncation-level accuracy, so there is no second regime to exploit. At $\varepsilon = 10^{-8}$, after the corrector has removed the smooth band, the smoother has to damp the whole band $|k| > 31$, whose sub-bands favour different solvers: undamped Jacobi is the cheapest per unit cost on the modes near the band edge (it contracts them by $\cos(2\pi\cdot 32/N)$-like factors that damped Jacobi cannot match), but it does not damp the near-Nyquist modes at all, which damped Jacobi and GS remove in a few sweeps. The oracle over $\{\text{Jacobi}, \text{Jacobi (0.67)}\}$---a \emph{nonstationary} Jacobi iteration whose damping parameter is chosen per macro-action---is therefore faster than every single solver, by $13$--$28\%$ on four of the six (equation, grid) settings at $128^2$ and $256^2$, and adding a third member (GS) changes nothing, i.e., the gain saturates. At $512^2$ multigrid dominates every subset (its cycle removes all sub-bands at once), so the oracle over any ensemble containing it equals the oracle over $\{\mathrm{NO}, \text{multigrid}\}$: adding members neither helps nor hurts. The ``difficulty'' that matters for ensembles is thus not the equation but the width of the spectrum left to the classical members after the corrector, which grows with the grid (fixed band) and with the tolerance.

\paragraph{Nested ensembles.} \Cref{tab:ca_ens_nest} evaluates the learned router over the nested sets $\{\text{Jacobi}\}$, $\{\text{Jacobi (0.67)}\}$, $\{\text{Jacobi}, \text{Jacobi (0.67)}\}$ and $\{\text{Jacobi}, \text{Jacobi (0.67)}, \text{GS}\}$ on every equation at $128^2$ and $256^2$; the single-member routers are the pairwise routers of \Cref{sec:wc_poisson,sec:wc_conv,sec:wc_aniso}, evaluated in the same session on the same instances, and all members keep the macro-action sizes of their pairwise runs. The table reports live and work-unit times to $h^2$, $10^{-6}$ and $10^{-8}$ for the router and its oracle, and the paired speedup at $10^{-8}$ over the best single-member router with its $p$-value (work units are the primary measure here: the nested runs share the machine with other jobs, so their live times are uniformly inflated relative to \Cref{sec:wc_poisson,sec:wc_conv,sec:wc_aniso}, while work units are determined by the decision sequence alone). Two properties are visible. \emph{Monotonicity}: the router over a larger set is never slower than the best of its members beyond timer noise, because the oracle ignores a member that never gives the largest reduction per unit cost and the router imitates it (adding GS to $\{\text{Jacobi}, \text{Jacobi (0.67)}\}$ leaves the decision sequence unchanged); the earlier ensembles of \Cref{tab:ca_ens} show the same for four- and five-member sets at $h^2$. \emph{Strict separation}: at $10^{-8}$ the two-member router is significantly faster than every single-member router in \caNestWins\ of the \caNestNum\ ensemble cells (paired Wilcoxon $p<0.01$; speedups \caNestSpMin--\caNestSpMax\ over the best single solver), reproducing its oracle, and it is never significantly slower. The learned policy is the nonstationary damping schedule described above: undamped Jacobi macro-actions while the band-edge modes dominate the residual, damped Jacobi once the near-Nyquist modes do, which the router detects from the decay rate of the residual over the last macro-action.

\begin{table}[H]
\color{red}
\centering
\caption{Oracle-level screening of ensembles: for every equation and grid, the fastest single solver and the fastest ensemble of up to three members under the cost-aware oracle (median work-unit time over $16$ test instances), at $\varepsilon = h^2$ and $10^{-8}$; the ratio is the speedup of the ensemble over the best single solver, and the last column is the ratio for $\{\text{Jacobi}, \text{Jacobi (0.67)}\}$.}
\resizebox{\textwidth}{!}{\caensscreen}
\label{tab:ca_ens_screen}
\end{table}

\begin{table}[H]
\color{red}
\centering
\caption{Nested ensembles: median time to $\varepsilon \in \{h^2, 10^{-6}, 10^{-8}\}$ of the learned router (live and work units) and of the cost-aware oracle (work units) over $\mathrm{NO}\cup\mathcal{W}$, for the single members and the nested two- and three-member sets, all evaluated in one session on the same instances ($64$ at $128^2$, $32$ at $256^2$, $16$ for anisotropic diffusion at $256^2$). The last columns give the paired speedup at $10^{-8}$ over the best single-member router (resp.\ oracle) with the one-sided Wilcoxon $p$-value; bold: $\ge 1.1\times$ with $p<0.01$.}
\resizebox{\textwidth}{!}{\caensnest}
\label{tab:ca_ens_nest}
\end{table}

\paragraph{Larger ensembles at truncation level.} \Cref{tab:ca_ens} repeats the ensembles of \Cref{sec:experiments} ($\{\text{Jacobi}, \text{GS}\}$ up to all five solvers) at $128^2$ and $\varepsilon = h^2$, where the screening predicts no gain: the ensemble router matches the best pairwise router of its members to within about $\pm 20\%$ in live time (ratio of medians \caEnsVsPairMin--\caEnsVsPairMax\ over the \caNumEns\ cells) and within $1.5\%$ in work units, without knowing in advance which member is fastest, and is \caEnsVsSolverMin--\caEnsVsSolverMax\ faster than the fastest member alone; \Cref{tab:ca_ensusage} lists how its iterations are distributed over the members. Its oracle is up to $25\%$ faster in a few cells because, at $h^2$, several cost-equalised smoothers are near-ties and the oracle's choice among them depends on the spectral content of the current error, which the router's residual-norm features reflect only indirectly (agreement with the oracle $55$--$83\%$ on its own rollouts against $88$--$99\%$ for pairwise routers). More capacity or imitation data does not change this: routers with twice the hidden width, three times the DAgger rounds and twice the oracle rollouts (\Cref{tab:ca_ens_big}) take the same number of iterations and corrector calls to $h^2$ as the default routers on \caEnsBigSameMin--\caEnsBigSameMax\ of the test instances and are significantly faster in \caEnsBigNFaster\ and significantly slower in \caEnsBigNSlower\ of the \caNumEnsBig\ cells ($p<0.05$); the residual disagreements are near-ties whose resolution does not affect the time to $h^2$. Deciding \emph{when} to call the corrector---the decision that dominates the wall-clock gains---is learned equally well in the ensembles.

\begin{table}[H]
\color{red}
\centering
\caption{Ensembles of \Cref{sec:experiments} at $128\times128$ and $\varepsilon = h^2$: median live time of the best member solver alone, of the best pairwise router among the members (paired speedup of the ensemble router over it in parentheses), and of the router and the cost-aware oracle over $\mathrm{NO}\cup\mathcal{W}$ ($64$ instances; pairwise baselines from the pairwise runs on the same instances).}
\resizebox{\textwidth}{!}{\caens}
\label{tab:ca_ens}
\end{table}

\begin{table}[H]
\color{red}
\centering
\caption{Ensembles at $128\times128$: mean (standard deviation over instances) fraction of the iterations up to the $h^2$ crossing that the ensemble router spends in each member and in the corrector.}
\resizebox{0.9\textwidth}{!}{\caensusage}
\label{tab:ca_ensusage}
\end{table}

\begin{table}[H]
\color{red}
\centering
\caption{Ensembles at $128\times128$: one-sided Wilcoxon $p$-values at $\varepsilon = h^2$ for the ensemble router against the fastest member solver alone (censored runs at their lower bound), against the best pairwise router in both directions, and against the cost-aware oracle over the same ensemble (alternative: the ensemble router is slower).}
\resizebox{0.9\textwidth}{!}{\castatsens}
\label{tab:ca_stats_ens}
\end{table}

\begin{table}[H]
\color{red}
\centering
\caption{Control on router capacity for the ensembles ($128\times128$, $64$ instances): the default ensemble router (hidden width $64$, $2$ DAgger rounds, $128$ oracle rollouts) against one with hidden width $128$, $6$ DAgger rounds and $256$ rollouts. Times to $h^2$ are median work units net of the per-decision cost; ``same decisions'' is the share of instances on which both routers take the same number of iterations and corrector calls to $h^2$; $p$ is the one-sided Wilcoxon $p$-value on the paired net times; ``agreement'' is each router's agreement with the oracle on its own rollouts during training.}
\resizebox{\textwidth}{!}{\caensbig}
\label{tab:ca_ens_big}
\end{table}

\subsection{Overheads and amortisation} \label{sec:wallclock_overheads}

\Cref{tab:ca_overheads} lists the measured cost of every operation as a function of $N$, including one decision of the LSTM router of \Cref{sec:experiments} (full-state input of size $4N^2$, hidden size 256, 4 layers): at $128^2$ one LSTM decision costs \caLstmMs\,ms, about \caLstmOverJacobi\ Jacobi sweeps, which is why the wall-clock study uses the scalar-feature router ($\approx 5\,\mu$s, independent of $N$). The learned router's decisions therefore account for well under $1\%$ of its solve time, and its corrector calls for $50$--$70\%$ of the time to $h^2$ (one or two calls). \Cref{tab:ca_amort} accounts for training: the corrector (shared with HINTS) is fit in minutes, the router in seconds to a few minutes per pairing, and the break-even number of solves after which the router's own training cost is recovered relative to HINTS ($\tau{=}25$) with the same corrector is in the thousands at $128^2$ and in the hundreds on the larger grids.

\begin{table}[H]
\color{red}
\centering
\caption{Measured single-thread cost of one operation as a function of the grid size (Poisson; median of repeated calls). Multigrid uses a direct solve below $64^2$ and is not listed there.}
\resizebox{0.9\textwidth}{!}{\caoverheads}
\label{tab:ca_overheads}
\end{table}

\begin{table}[H]
\color{red}
\centering
\caption{Training-cost amortisation for the GS pairing: corrector training (data generation and least-squares fit, CPU), router training (oracle rollouts and DAgger rounds), median wall-clock saved per solve at $\varepsilon = h^2$ relative to HINTS ($\tau{=}25$) and to the solver alone, and the number of solves after which the router's training cost is recovered relative to HINTS.}
\resizebox{\textwidth}{!}{\caamort}
\label{tab:ca_amort}
\end{table}

\subsection{Verification of the theoretical assumptions} \label{sec:wallclock_assumptions}

\Cref{th:suboptgreedy} and Propositions~\ref{th:weaklyalphasupermodular}--\ref{th:simultaneoussupermodular} assume that every error propagation map $G_j = I - C_j\circ\mathcal{L}_h^a$ is Lipschitz with constant $\rho_j < 1$ (respectively $\le 1$ with a shared eigenbasis), zero-preserving, and---for postfix monotonicity---invertible; \Cref{th:consistency} assumes that the per-step errors $\tilde c_j$ are bounded above and that at least one solver cannot annihilate the error in one step. \Cref{tab:ca_assump} verifies these conditions numerically for every operation of the study on the mean-free subspace (the constant mode is the null space of the periodic operator and is excluded from the data and from every iteration): the Euclidean Lipschitz constant $\|G_j\|_2$ (power iteration on $G_j^{\top}G_j$; exact from the Fourier symbol for Jacobi), the Lipschitz constant $\|G_j\|_A = \|A_s^{1/2} G_j A_s^{-1/2}\|_2$ in the energy norm of the symmetric part $A_s$ of the discrete operator, the same for the macro-action $G_j^{m_j}$, the spectral radius, the smallest observed contraction $\sigma_{\min} = \min\|G_j x\|/\|x\|$, $\|C_j(0)\|$, and the commutator $\|G_j G_{\mathrm{NO}} - G_{\mathrm{NO}} G_j\|$ on random vectors. \Cref{tab:ca_assump_paths} evaluates the quantities of \Cref{th:suboptgreedy} and \Cref{th:consistency} along the cost-aware oracle's own path to $\varepsilon = 10^{-8}$ (a horizon of tens of macro-actions; to $h^2$ the path has one or two steps).

\paragraph{What holds, and in which norm.} (i) Every classical solver is zero-preserving exactly and contracts in the sense that matters for convergence: the spectral radius of $G_j$ is below $1$ in every setting (at most \caRhoSpecMax), and so is the energy-norm Lipschitz constant of the solver and of its macro-action (at most \caRhoAMax). The \emph{Euclidean} Lipschitz constant, however, exceeds $1$ for Gauss--Seidel, SOR and symmetric GS (up to \caRhoTwoGsMax\ for SOR), whose iteration matrices are non-normal; the proofs of Propositions~\ref{th:weaklyalphasupermodular}--\ref{th:simultaneoussupermodular} use only the triangle inequality and the composition rule for Lipschitz constants, so they hold verbatim with $\|\cdot\|_2$ replaced by any norm, and in the energy norm the Lipschitz condition is satisfied by all classical solvers in all settings. For (damped) Jacobi both constants agree (the maps are Fourier-diagonal); the values are close to $1$ because the slowest modes (the smoothest ones, and the near-Nyquist mode for undamped Jacobi) are barely damped by a point smoother---which is exactly why the corrector is needed. (ii) The corrector's map is zero-preserving, contracts every band mode to $\caBandMax$ or better (on average to $10^{-4}$--$10^{-6}$) and leaves the modes outside the band untouched, so its Lipschitz constant is exactly $1$ (not $<1$): \Cref{th:weaklyalphasupermodular} in its strict form does not apply to the corrector, but \Cref{th:simultaneoussupermodular} does, and its second condition---a shared eigenbasis---holds exactly for the ensembles $\{\mathrm{NO}, \text{Jacobi}\}$, $\{\mathrm{NO}, \text{Jacobi (0.67)}\}$ and $\{\mathrm{NO}, \text{Jacobi}, \text{Jacobi (0.67)}\}$ (all maps are diagonal in the Fourier basis; measured commutators $\approx 10^{-6}$, at the corrector's float32 precision), so $h$ is supermodular for them and \Cref{th:suboptgreedy} holds with $\alpha = 1$; for GS, SOR, symmetric GS and multigrid the commutators are small but nonzero ($10^{-3}$--$10^{-1}$). (iii) Invertibility fails for most maps in exact arithmetic: a Gauss--Seidel, SOR or symmetric-GS sweep annihilates an error concentrated at the first grid point ($\sigma_{\min} = 0$), undamped Jacobi annihilates the mode $(N/4, N/4)$, and the corrector annihilates the band; postfix monotonicity is therefore not guaranteed, which \Cref{sec:greedy} anticipates (``our experiments show that greedy routers remain effective even when invertibility is not met''). (iv) The worst-case ratio $\alpha(O)$ of \Cref{th:weaklyalphasupermodular}, which only uses the Lipschitz constants, is loose here: along the oracle's path to $10^{-8}$ (\Cref{tab:ca_assump_paths}) $\sum_i \rho_{O_i}^2$ is close to $T$ because each $\rho$ is dominated by the slowest mode, so the bound ranges from $\caAlphaBoundMin$ (the multigrid pairing, whose cycle contracts everything) to several hundred, and is infinite on the few instances where $\sum_i \rho_{O_i}^2 \ge T$. The empirically measured supermodularity ratio of \Cref{eqn:alpha_supermod} along the same path---with $S$ the oracle's prefix and $S'$ its suffix---is at most \caAlphaHatMax\ (median \caAlphaHatMed) in every setting, and at most \caAlphaHatMaxH\ on the short paths to $h^2$: the objective behaves supermodularly on the states the greedy rule actually visits, which is the regime in which \Cref{th:suboptgreedy} gives the $(1-e^{-1})$ guarantee. (v) For \Cref{th:consistency}, the per-step errors are bounded by the current error ($\max_j \tilde c_j / \|e^{(t)}\|^2 \le 1$ at every visited state) and at least one solver leaves a positive fraction of it ($\min_t \max_j \tilde c_j / \|e^{(t)}\|^2 > 0$), so the surrogate's Bayes consistency applies.

\begin{table}[H]
\color{red}
\centering
\caption{Numerical verification of the assumptions of Propositions~\ref{th:weaklyalphasupermodular}--\ref{th:simultaneoussupermodular} for every operation of the study (mean-free subspace). $\|\cdot\|_2$: Euclidean Lipschitz constant of the error propagation map; $\|\cdot\|_A$: Lipschitz constant in the energy norm of the symmetric part of $\mathcal{L}_h$, for one iteration and for the macro-action of $m_j$ iterations; $\rho$: spectral radius; $\sigma_{\min}$: smallest observed contraction (a unit error at the first grid point for the triangular sweeps, the minimum over Fourier modes otherwise; for the corrector, the largest remaining fraction of a band mode); $\|C_j(0)\|$: zero preservation; last column: commutator with the corrector's map on random unit vectors (Proposition~\ref{th:simultaneoussupermodular}). For (damped) Jacobi the constants are exact from the Fourier symbol (the value with the Nyquist modes included; without them undamped Jacobi has $\|G\|_2 = \rho = 0.9994$ at $128^2$).}
\resizebox{\textwidth}{!}{\caassump}
\label{tab:ca_assump}
\end{table}

\begin{table}[H]
\color{red}
\centering
\caption{Quantities of \Cref{th:suboptgreedy}, Proposition~\ref{th:weaklyalphasupermodular} and \Cref{th:consistency} along the cost-aware oracle's path to $\varepsilon = 10^{-8}$ ($8$ test instances; medians unless stated): number of macro-actions $T$, $\sum_i\rho_{O_i}^2$ with the energy-norm constants of the chosen macro-actions, the resulting $\alpha(O)$ of Proposition~\ref{th:weaklyalphasupermodular} (median over the instances on which $T - \sum_i\rho_{O_i}^2 > 0$; $\infty$ otherwise), the empirical supermodularity ratio $\hat\alpha$ of \Cref{eqn:alpha_supermod} evaluated with $S$ the oracle's prefix and $S'$ its suffix (maximum and median over prefixes and instances), and the normalised bounds of \Cref{th:consistency}: the largest ratio $\max_j \tilde c_j/\|e^{(t)}\|^2$ over visited states ($\bar E$) and the smallest value of $\max_j \tilde c_j/\|e^{(t)}\|^2$ ($E_{\min}$).}
\resizebox{\textwidth}{!}{\caassumpB}
\label{tab:ca_assump_paths}
\end{table}
}  % \edit-equivalent group (a macro argument this large breaks the PDF page tree)

\section{Limitations}

A primary limitation of our approach is the higher on-time computational cost requires to learn the greedy routing strategy. In contrast, the simpler baselines such as HINTS incur no training overhead. While this cost is amortized during inference time and leads to improved solution quality, it may be prohibitive in settings with limited computational resources or when rapid deployment is required. \edit{On the deployment side, \Cref{sec:wallclock} shows that the gains survive cost accounting: with a corrector that covers the smooth half of the spectrum and a classical smoother that covers the rest, one or two corrector calls suffice and the router's advantage over fixed schedules is uniform across solvers, equations and tolerances. The corrector we use is specific to the constant-coefficient periodic setting (a fixed Fourier trunk basis on a coarse sensor grid), its cost grows with the fourth power of the sensor-grid size, and undamped Jacobi remains limited far below truncation level by the near-Nyquist modes that neither component damps. Our wall-clock study is restricted to single-thread CPU execution; with a corrector band that does not grow with the grid, the gains over fixed schedules shrink at $512^2$ for the point smoothers, whereas pairing the corrector with geometric multigrid restores them (\Cref{sec:wallclock}). Ensembles of classical solvers pay off only where the error passes through regimes that favour different members; at truncation level a single well-chosen smoother suffices. For the constant-coefficient periodic test problems the discrete operator is diagonalised by the FFT, and the resulting direct solve, reported in every table, is faster than every iterative method we measured, including ours (\caFftRatioMin--\caFftRatioMax\ against the best router); the hybrid is meant for problems on which no such direct diagonalisation exists and the corrector must be learned rather than computed, which is why \Cref{sec:wallclock} also reports a variable-coefficient diffusion problem.}

Additionally, the effectiveness of the learned router depends on the neural operator exhibiting heterogeneous performance across samples. When the neural operators produces predictions of relatively uniform quality, the benefit of adaptive routing diminishes and it is more practical to set a fixed schedule based on the quality of the neural operator. In practice, achieving such heterogeneity can be sensitive to training choices, and neural operators themselves can be finicky to train.


\section{LLM Usage}

LLMs, specifically ChatGPT and Gemini, supported the writing process in an iterative manner. We drafted paragraphs and asked the models for feedback on grammar and clarity. We then incorporated selected suggestions into the writing and repeated this process until we were satisfied with the writing. 

The code developed for the experiments was written by the authors with the help of occasional code completions. The central components (e.g., the hybrid solver implementation and the greedy-router training pipelines) were implemented exclusively by the authors.

All substantive intellectual contributions, which include ideas, theorems, and analyses, are our own. LLMs were occasionally used to verify the correctness of proofs, but all proof strategies originated from the authors and relevant literature.

% \section{New Multiple Solver Usage Results}

% \begin{table}[t]
% \centering
% \caption{
% Final error and AUC of squared $L^2$ error for ensembles of increasing size. Here $\mathcal{W}$ denotes the set of numerical solvers and
% $\mathrm{NO}$ the neural operator (DeepONet). Each $\mathrm{Router}(\mathrm{NO}\cup\mathcal{W})$
% is compared against the best pairwise routing using our method, i.e., $\mathrm{Router}(\mathrm{NO}\cup \{\text{solver}\})$ for $\text{solver} \in \mathcal{W}$.
% Values are mean ($\pm$ s.e.) over 64 test instances; Error is reported in
% $\times 10^{-3}$.
% }
% \label{tab:errorcomparison2}

% \resizebox{\textwidth}{!}{
% \begin{tabular}{llcccc}
% \hline
% \multirow{2}{*}{\textbf{$\mathcal{W}$}}
% & \multirow{2}{*}{\textbf{Method}}
% & \multicolumn{2}{c}{\textbf{Poisson}}
% & \multicolumn{2}{c}{\textbf{ConvDiff}} \\
% \cline{3-6}
% &
% & $\|e^{(T)}_h\|\times 10^3$ & AUC
% & $\|e^{(T)}_h\|\times 10^3$ & AUC \\
% \hline

% \multirow{2}{*}{$\{\text{Jacobi, GS}\}$}
% & Best $\mathrm{Router}(\mathrm{NO}\cup \{\text{solver}\})$
% & 0.003 (0.010) &  0.115 (0.285) &   $< 10^{-3}$ &  0.035 (0.095)\\
% & Router$(\mathrm{NO}\cup\mathcal{W})$
% & 0.001 (0.001) & 0.080 (0.153) & $<10^{-3}$ & 0.031 (0.076) \\
% \hline

% \multirow{2}{*}{$\{\text{Jacobi, GS, SymGS}\}$}
% & Best $\mathrm{Router}(\mathrm{NO}\cup \{\text{solver}\})$
% & $< 10^{-3}$ & 0.074 (0.194)  & $< 10^{-3}$ & 0.026 (0.050)  \\
% & Router$(\mathrm{NO}\cup\mathcal{W})$
% & $<10^{-3}$  & 0.058 (0.101) & $<10^{-3}$  & 0.023 (0.044) \\
% \hline

% \multirow{2}{*}{$\{\text{Jacobi, GS, SymGS, Jacobi (0.67)}\}$}
% & Best $\mathrm{Router}(\mathrm{NO}\cup \{\text{solver}\})$
% & $< 10^{-3}$ & 0.074 (0.194)  & $< 10^{-3}$ & 0.026 (0.050) \\
% & Router$(\mathrm{NO}\cup\mathcal{W})$
% & 0.005 (0.016) & 0.085 (0.218) & $<10^{-3}$  & 0.023 (0.046)\\
% \hline

% \multirow{2}{*}{$\{\text{Jacobi, GS, SymGS, Jacobi (0.67), SOR (1.5)}\}$}
% & Best $\mathrm{Router}(\mathrm{NO}\cup \{\text{solver}\})$
% & $<10^{-3}$ & 0.051 (0.108)
% & $<10^{-3}$ & 0.015 (0.043) \\
% & Router$(\mathrm{NO}\cup\mathcal{W})$
% & $<10^{-3}$  & 0.049 (0.102) & $<10^{-3}$  & 0.010 (0.024)\\
% \hline

% \end{tabular}
% }
% \end{table}


% \paragraph{Solver Ensembles Experiments:} We now consider routing over a collection of more than two solvers. We study the performance of our router across ensembles of increasing size, denoted by $\mathcal{W}$, consisting of a trained DeepONet and a diverse set of numerical solvers, including Jacobi, Gauss-Seidel (GS), Symmetric Gauss Seidel (SymGS), Damped Jacobi with $\omega = 0.67$ and Successive Over-relaxation (SOR) with $\omega = 1.5$ . Our goal is to evaluate whether enlarging the solver set yields gains beyond any pairwise configurations. To this end, we compare $\mathrm{Router}(\mathrm{NO}\cup\mathcal{W})$ against the best pairwise version of our method, i.e., $\mathrm{Router}(\mathrm{NO}\cup\{\text{solver}\})$ for solver $\in \mathcal{W}$. As shown in \Cref{tab:errorcomparison2}, routing over larger ensembles often improves both final error and AUC relative to the best pairwise configuration, indicating that additional solver diversity can be leveraged. As in \Cref{fig:router_strat}, the router achieves these gains by learning to route to each member of the ensemble for a non-trivial proportion of the iterates, as measured in \Cref{sec:solverdecisions}. Finally, despite the increased size of the action space, the learned router continues to track these improvements, demonstrating robustness of both the training procedure and the convex surrogate of \Cref{section:surrogate}. These results demonstrate that our method can scale to richer solver sets, enabling improved performance without sacrificing learnability.
% Additional experiments are provided in \Cref{sec:moreexperimentalresults}.

\clearpage



\input{checklist.tex}


\end{document}